\chapter{动态场景与4D}

\paragraph*{本章阅读锚点}
静态3DGS把空间中的高斯作为可微渲染基元，但默认每个高斯只对应一个不随时间变化的场景。动态场景的核心困难不是把时间简单加到输入里，而是要回答三个更具体的问题：高斯是否在所有时间持久存在，运动是由逐高斯轨迹、规范空间变形场、显式4D基元还是控制节点描述，外观和几何在遮挡、模糊、单目欠约束和长视频中如何保持一致。本章沿着这些问题展开。代表性方法主要构成四类锚点：Dynamic 3D Gaussians把高斯看成持久粒子，Deformable 3D Gaussians把动态建模转成规范空间到目标时间的变形，4D-GS把规范高斯和时空编码结合，Spacetime Gaussian Feature Splatting则把高斯本身扩展到时空基元。后续工作大多是在这些锚点上改变运动参数化、监督信号、静动态分解、流式压缩或生成补全方式。

\section{变形场与4D高斯表示}

这一节关注动态3DGS的主干表示。最直接的路线是为每个时间训练一份3DGS，但这样会让存储和优化随帧数线性增长，并且不能天然给出时间连续轨迹。因此，主流方法通常保留一套规范空间高斯，再用时间条件模块把它变换到目标帧；或者把高斯直接提升为4D时空基元，让时间中心、持续时间、轨迹和外观随时间共同优化。两条路线的区别在于：前者更像“静态高斯加变形网络”，后者更像“高斯本身拥有时空支撑域”。



% METHOD: 2310.08528
\paragraph{4D-GS（HexPlane 驱动的规范高斯变形）（2023.10）}
4D-GS~\cite{2310_08528} 仍使用规范 3D Gaussian，但以结构化时空特征替代逐高斯逐帧参数或纯坐标 MLP。对坐标 $(x,y,z,t)$，多分辨率 HexPlane 将其投影到三张空间--空间平面和三张空间--时间平面；每张平面双线性采样，成对特征相乘并跨分辨率拼接，再由轻量 MLP 解码位置、旋转和尺度偏移。规范高斯的颜色和 opacity 保持在显式基元中，变形后走标准 splatting。

\textbf{训练与渲染。} 先以静态 3DGS 建立规范点云，再联合优化高斯、平面特征和多个形变头。每个训练图像携带时间戳：查询平面、预测偏移、投影变形后的协方差，按深度 alpha 合成并以图像重建项反传；densification/pruning 作用于规范高斯而不是为每帧复制基元。六张二维平面把直接四维体的存储和查询拆开，使邻近空间和时间共享特征。

该方法与 Deformable 3DGS 都是 canonical-plus-deformation，但它把时空局部性编码在显式平面中；与原生 4D Gaussian 则不同，时间不是高斯协方差的一维，而是生成变形的查询条件。平面分解提高效率却隐含低秩/可分解假设，快速大幅运动、遮挡或超出训练时间的变化仍可能使共享特征混叠。

% METHOD: 2310.10642
\paragraph{Real-time Photorealistic Dynamic Scene Representation and Rendering with 4D Gaussian Splatting（原生时空高斯）（2023.10）}
Real-time Photorealistic Dynamic Scene Representation and Rendering with 4D Gaussian Splatting~\cite{2310_10642} 不以规范空间和外部 deformation field 表达运动，而在 $(x,y,z,t)$ 中直接放置各向异性 4D Gaussian。每个基元学习 4D 均值和经 4D 缩放、4D 旋转构成的协方差；空间--时间协方差交叉项使固定时间的条件分布给出相应时刻的三维中心和协方差，随后即可按 3DGS 投影、排序和 alpha 合成。

\textbf{时间切片与外观。} 渲染时间 $t$ 时，系统不是移动一套 canonical Gaussian，而是从每个 4D 分布计算以该时间为条件的 3D Gaussian，并把它送入光栅化器。外观采用 Spherindrical Harmonics：以视线方向的球谐基和时间的 Fourier 基组合，使颜色可同时依赖方向和时间。这样时空几何及时间变化的方向外观都归入单一原语，而非另由变形网络预测。

\textbf{端到端密度控制。} 训练以渲染图像误差端到端更新 4D 原语，并交替进行 density control。除视空间位置梯度外，论文把时间均值梯度纳入 densification 指示，并在 split 时联合采样空间和时间位置；对剧烈变化场景使用时间批采样以减轻闪烁。它可在任意训练时间与视角切片，但四维协方差、旋转和方向--时间外观增加每原语参数，长序列及复杂变化的密度和优化代价仍是直接时空表示的边界。

% METHOD: 2312.03431
\paragraph{Gaussian-Flow（双域显式属性轨迹）（2023.12）}
Gaussian-Flow~\cite{2312_03431} 直接为每个 Gaussian 的时间变化属性拟合曲线，而不查询隐式 deformation network。以参考时刻的中心、旋转与辐射为基值，位置、四元数和颜色的残差写为低阶多项式与 Fourier 级数之和；尺度与 opacity 保持不变。多项式覆盖平滑趋势，频域项补充快速变化，每一属性维度有自己的系数。

为避免短时间大运动需要不稳定的大系数，方法还为每粒子学习时间缩放与偏移，再将缩放时间送入双域曲线。当前时刻的属性由解析曲线计算后，仍按 3DGS 投影和 alpha 合成。时间扰动一致性约束残差平滑，KNN 刚性约束保持局部空间连贯；优化还交替处理邻域随 densification 变化的问题。它保留了显式粒子的快速查询，却把长序列复杂运动的容量交给每粒子的系数，因而频率/阶数和正则必须在拟合能力与振荡之间权衡。

% METHOD: 2312.16812
\paragraph{Spacetime Gaussian Feature Splatting（时间可见性、参数轨迹与特征渲染）（2023.12）}
Spacetime Gaussian Feature Splatting~\cite{2312_16812} 将动态基元写为带时间寿命的 Spacetime Gaussian。每个高斯用一维 Gaussian/RBF 时间 opacity 决定何时可见，并用多项式参数化中心与旋转随时间的变化；因此静态、高速移动和短暂出现内容可以由不同的时间中心与宽度区分，而不是共享一套始终可见的高斯。

外观不再以逐高斯 SH 直接解码：高斯携带低维 neural feature，先按 alpha 合成 splat 到图像特征，再由小网络结合视线与时间解码 RGB。密度增长同时参考训练误差图和粗深度，在难收敛区域新增高斯。时间查询时计算有效 opacity、轨迹和旋转后走同一 rasterizer。该设计比纯多项式位置模型多出显式存在区间，但多项式轨迹和时间 RBF 的阶数/尺度限制了极复杂长时动态。

% METHOD: 2404.03613
\paragraph{Per-Gaussian Embedding-Based Deformation（逐高斯身份码的粗细变形）（2024.04）}
Per-Gaussian Embedding-Based Deformation~\cite{2404_03613} 认为坐标加时间的单一 deformation field 忽略了 3DGS 是许多局部基元的混合。方法为每个 Gaussian 学习身份 embedding，并将其与时间 embedding 输入变形网络，直接预测该基元的位置、旋转、尺度、不透明度和外观的时间变化。

变形分为 coarse 与 fine 两个残差：前者承担慢速或大尺度运动，后者在其上补快速局部变化；相邻高斯 embedding 的 KNN 平滑项使动态区域不出现各自独立的噪声轨迹。查询时间先生成属性残差，再使用常规 Gaussian splatting 渲染并由图像误差反传。与坐标 MLP 相比，逐基元码增加局部自由度；代价是 embedding 存储和优化规模随高斯数增长，拓扑突变仍不由连续残差自然解决。

% METHOD: 2404.06270
\paragraph{3D Geometry-aware Deformable Gaussian Splatting（几何条件变形）（2024.04）}
3D Geometry-aware Deformable Gaussian Splatting~\cite{2404_06270} 将动态场景拆为 canonical Gaussian field 和 deformation field，但不只以坐标和时间预测位移。canonical 分支从高斯附近的三维结构提取局部 geometry feature；变形网络以该特征和时间为条件输出位置与旋转变化，使相同时间下不同局部形状可有不同、但几何连贯的运动。

训练将 canonical 高斯、几何特征与变形联合放入可微 splatting，通过图像监督更新；当前时间的变换后高斯按深度合成。显式几何条件把空间邻域信息带入 deformation，而不是仅希望 MLP 从坐标中自行发现它。其有效性仍受 canonical 几何和局部特征质量约束，遮挡、稀疏单目观测或显著拓扑改变会使几何条件本身不稳定。

% METHOD: 2405.17421
\paragraph{MoSca（4D Motion Scaffold 锚定的高斯融合）（2024.05）}
MoSca~\cite{2405_17421} 从随意单目视频重建动态场景，先以基础视觉模型的跟踪/深度先验建立稀疏 Motion Scaffold 图。图的节点存储随时间变化的 SE(3) 轨迹及控制半径，边按轨迹曲线距离连接近邻；该图以远少于场景点的变量编码低秩、平滑的运动结构。

任一空间点的稠密变形由附近节点的 dual-quaternion blending 产生，并以 ARAP 类正则维持局部合理性。各时刻观测中的高斯锚到 scaffold 后被带到共同表示中全局融合，几何和外观与变形分离优化；渲染误差经 splatting 更新二者。系统还在 bundle adjustment 中联合估计相机焦距和位姿，而不把外部 pose 工具作为前提。Motion scaffold 能把不同时间的观测合并，却依赖二维先验、图拓扑和节点数；极快变形、严重遮挡或错误相机初始化会影响全链路。

% METHOD: 2406.10324
\paragraph{L4GM（视频到逐帧 Gaussian 的前馈大重建模型）（2024.06）}
L4GM~\cite{2406_10324} 将 3D LGM 扩展为单视角视频到动画 3D 资产的前馈模型，不进行逐场景 4D 优化。模型从低帧率采样的视频帧预测每帧的 3D Gaussian 表示；在预训练 LGM 的基础网络中加入 temporal self-attention，使不同时间的高斯预测共享视频上下文，并以每个时间点的多视图渲染损失训练。

为输出更高帧率，独立插值网络接收相邻两帧 Gaussian 表示并预测中间表示，而不是只在图像空间插帧。训练数据由多视角动画对象视频构成，推理时输入单视角视频即可一次产生时间序列高斯。其速度来自前馈化和强训练先验；输出质量与对象类别、合成训练分布及逐帧表示的时间一致性密切相关，复杂场景或未见动态并没有显式几何保证。

% METHOD: 2408.13036
\paragraph{H3D-DGS（异构控制点的可观测--不可观测运动分解）（2024.08）}
H3D-DGS~\cite{2408_13036} 不让所有动态控制点的六自由度完全由渲染梯度猜测。以局部相机射线坐标分解运动后，投影到图像平面的平移分量及相应可观测旋转由 optical flow 反投影直接给出；沿视线的平移与其余不可观测旋转再由可微渲染的梯度优化补全。两类信息合成 H3D control points，并只部署于动态物体，背景保持静态。

控制点动作通过局部射线近似传播到邻域高斯，当前时刻再由变形后的高斯进行 splatting。该混合策略保留稀疏控制表示的紧凑性，并将二维中直接可见的运动从非凸优化中剥离；它仍依赖光流质量、相机标定和局部近似，深度方向及遮挡后的运动没有直接观测，需由优化和正则推断。

% METHOD: 2409.11211
\paragraph{SplatFields（2024.09）\cite{2409_11211}}
SplatFields 处理的是另一个常被动态方法忽略的问题：3DGS 的每个高斯几乎是独立参数表，位置接近的高斯可以学到完全不相关的颜色、尺度、透明度和位移。在密集多视角里，这种自由度能换来细节；在 3-view、6-view 或稀疏多相机动态捕获里，它也会把训练视角里的噪声、遮挡和偶然颜色记成几何。SplatFields 的改变不是换渲染器，而是把“逐高斯自由参数”改成“由连续神经场生成的高斯属性”：高斯仍按 3DGS 的方式投影和 $\alpha$-混合，但特征、位置修正、尺度、不透明度等由共享的 MLP 或三平面 CNN 特征场预测，附近高斯因此受到同一个连续函数约束。

具体流程可以理解为两层。第一层是基础高斯集合，它仍由稀疏点、视觉壳或已有 3DGS 初始化，负责提供可渲染的显式基元。第二层是 SplatField：网络以高斯坐标为输入，静态重建时预测该位置的 splat 特征和位移；动态重建时再接入时间或序列条件，使同一局部结构在时空上共享参数。训练时图像重建损失、mask 损失和动态场景中的退火平滑约束仍通过 splatting 反传，但梯度先更新共享神经场，再落到一批空间相邻的高斯属性上。这样做牺牲的是训练阶段的一部分速度和显存，换来稀疏视角下更强的空间自相关先验；训练完成后，神经场预测出的高斯属性可以固定下来，推理仍接近普通 3DGS 的实时渲染速度。

因此，SplatFields 在本章里是“稀疏输入正则化”的局部锚点。它与 4D-GS、Deformable 3DGS 的关系不是替代运动模型，而是补上一个参数来源约束：4D-GS 关心时间变形如何表示，SplatFields 关心高斯属性是否应该由相邻空间共享的函数生成。论文在 Blender、DTU、Owlii 和 NeRF-DS 等静态/动态稀疏设置中验证了这一点，尤其在视角越少时，相比原始 3DGS 和多种 splatting 基线越能减少过拟合。
\textbf{方法信息：} Meta 参与了这篇论文；代码：\url{https://github.com/markomih/splatfields}；GitHub 关注度较高，约 187 stars；引用很多，Semantic Scholar 约 56 次。



% METHOD: 2410.15629
\paragraph{Fully Explicit Dynamic Gaussian Splatting（2024.10）\cite{2410_15629}}
把动态表示推向完全显式。它把高斯分成静态和动态两组：静态高斯只保留简单的时间变化，动态高斯则在稀疏关键时间存储位置和旋转，查询任意时间时用三次Hermite插值生成平滑位置轨迹，用Slerp插值旋转，并用时间不透明度的高斯混合描述物体出现、消失或遮挡变化。训练不是一次覆盖完整长序列，而是先在短时间窗口内稳定关键帧参数，再逐步延长时间范围；如果某些高斯在回溯过程中累积误差过大，点回溯机制会把它们识别为错误动态基元并移除。相对于4D-GS，它取消了隐式时空编码器，把动态性改写成“关键帧状态 + 插值规则 + 可见性函数”，因此轨迹更可解释，也更适合强调高速渲染的场景。
\textbf{方法信息：} 引用很多，Semantic Scholar 约 75 次。
\textbf{设计动机与边界：} 相对原始 3DGS 以独立高斯和图像重建误差为主的设定，Fully Explicit Dynamic Gaussian Splatting 把问题落在以下可实现的改动上：方法包含四个核心设计：(1)静态/动态高斯体分离：静态高斯体(Gs)建模为位置随时间线性变化的3DGS；动态高斯体(Gd)基于关键帧插值，在稀疏时间戳采样位置和旋转然后插值；(2)时间插值：位置使用三次Hermite插值器(CHip)实现平滑连续运动，旋转使用球面线性插值(Slerp)，时间不透明度使用高斯混合模型处理物体可见性变化；(3)渐进训练方案：初始使用短时间序列训练，逐步延长时间，使模型能在稀疏点云条件下工作良好；(4)点回溯技术：量化每个高斯体随时间的累积误差，检测并移除动态场景中的错误高斯体。其设计目的不是单独增加一个模块，而是用上述表示、监督或渲染约束消除原流程在目标任务中的具体失配。提出完全显式的4DGS表示，通过关键帧插值和静/动态分离实现高效动态场景渲染，在2080Ti上达到62fps的SOTA渲染质量适用时还应保留论文给出的边界：关键帧间隔需均匀假设限制了灵活性；对极端快速非刚性变形的建模能力有限；渐进训练增加训练时间
% METHOD: 2410.17249
\paragraph{SpectroMotion（2024.10）\cite{2410_17249}}
处理动态镜面场景。普通动态3DGS即使几何和颜色对齐，也很难让镜面高光随法线和光照正确变化，因为高光不是固定贴图，而是由表面朝向、相机方向和环境光共同决定。SpectroMotion在Deformable 3DGS上加入物理渲染：先用标准3DGS稳定基础几何，再训练动态变形；变形过程中对法线做残差校正，避免镜面物体运动后法线和几何脱节；随后引入随时间变化的环境图，并用基于图像的光照模型渲染镜面反射。相对于只优化RGB重建的动态高斯，它把法线、环境光和镜面反射的物理关系纳入最终颜色生成，因此适合金属、反光塑料等动态物体。
\textbf{实现细节：} SpectroMotion（2024.10） （CVPR 2025，Fan et al.）指出现有 3DGS 动态扩展在镜面表面表示上存在根本性缺陷：动态场景方法（Deformable 3DGS、4DGS）依赖低阶球谐表示颜色，其固有的低频特性无法捕捉镜面高光等复杂视觉效果；而镜面渲染方法（如 GaussianShader、GS-IR）虽然引入了环境贴图与 BRDF 模拟镜面反射，但标准环境贴图无法建模时变光照，因此仅适用于静态场景。镜面反射对法线极度敏感：法线微小偏差 $\delta n$ 可导致反射方向 $\omega_r^t = 2(\omega_o^t \cdot \mathbf{n}^t)\mathbf{n}^t - \omega_o^t$ 的剧烈偏转，因此精确的法线估计对动态镜面场景至关重要。SpectroMotion 提出三项核心创新：残差法线校正（residual correction）用于变形过程中的精确法线计算、可变形环境贴图（deformable environment map）适应时变光照条件、三阶段粗到精训练策略显著增强场景几何和镜面颜色预测。它是首个也是目前唯一能够合成真实世界动态镜面场景的 3DGS 方法。
\textbf{实现补充：} 实现上还要核对输入、输出、参数、特征、初始化、训练损失、正则、渲染或推理接口，以及优化中的剪枝、分裂、采样、编码或解码步骤；这些抓手决定该方法到底改了 3DGS 管线的哪一环。
\textbf{方法信息：} NVIDIA 参与了这篇论文。
\textbf{设计动机与边界：} 相对原始 3DGS 以独立高斯和图像重建误差为主的设定，SpectroMotion 把问题落在以下可实现的改动上：方法结合3DGS、物理渲染(PBR)和变形场：(1)残差法线校正：在变形过程中对表面法线进行残差校正，确保动态物体法线计算的准确性，这对镜面反射渲染至关重要；(2)可变形环境图：设计随时间变化的环境贴图适应动态场景中的时变光照条件，而非使用固定的环境光照；(3)三阶段粗到细训练策略：静态阶段使用vanilla 3DGS稳定场景几何；动态阶段使用Deformable 3DGS建模动态物体，通过法线优化提升几何质量；镜面阶段采用基于图像的光照(IBL)模型进行镜面渲染。整个框架实现端到端优化。其设计目的不是单独增加一个模块，而是用上述表示、监督或渲染约束消除原流程在目标任务中的具体失配。提出首个能合成真实世界动态镜面场景的3DGS方法，通过残差法线校正、可变形环境图和三阶段粗到细训练，实现动态镜面场景的SOTA渲染适用时还应保留论文给出的边界：三阶段训练增加训练时间；对极端镜面（完美镜面）的渲染仍有挑战；依赖环境图的表达能力
% METHOD: 2410.20815
\paragraph{Grid4D（4D 分解哈希动态高斯）（2024.10）\cite{2410_20815}}
Grid4D 直接接在 4D-GS 的时空编码问题之后。4D-GS 用六个二维平面来分解 $(x,y,z,t)$，这样训练和查询都很轻，但它默认动态特征可以由若干平面低秩组合表达；当场景组件运动方向复杂、局部形变差异很大时，不同平面上的特征容易互相重叠，变形解码器会把本来属于不同运动的信号混在一起。Grid4D 的改变不是放弃“规范高斯 + 时间变形”框架，而是把时空特征底座从 2D 平面换成更高表达力的 3D 哈希分解。

具体做法是把 4D 输入拆成四个 3D 哈希编码：一个纯空间编码 $(x,y,z)$，以及三个时空编码 $(x,y,t)$、$(y,z,t)$、$(x,z,t)$。纯 4D 哈希表会有很高存储成本，2D 平面又过强依赖低秩假设；四个 3D 哈希表在二者之间折中，既保留三维局部邻域，又避免完整 4D 网格的 $O(n^4)$ 开销。每个编码采用多分辨率哈希查询得到特征后，Grid4D 再用 directional attention 聚合空间特征和时空特征：空间编码给出“这个高斯属于哪个局部结构”，时空编码给出“该结构在不同方向上的时间变化”，注意力权重决定哪些时空分支更应该影响当前高斯的变形。

训练流程仍然沿用动态 3DGS 的可微渲染闭环：从规范高斯开始，查询 Grid4D 编码和轻量解码器，预测位置、旋转或尺度等变形残差，再把变形后的高斯交给 rasterizer，与目标帧图像监督。为避免显式哈希在相邻时间产生抖动，论文加入时间平滑正则，使相邻时间步的变形保持连续。相对 4D-GS，Grid4D 的锚点意义在于说明动态高斯的瓶颈不只在“有没有时间编码”，还在“时空编码的分解假设是否足够表达复杂运动”；后续 DASH 等方法继续沿着哈希编码路线，只是进一步把静态区域排除在动态哈希之外。官方代码为 \url{https://github.com/JiaweiXu8/Grid4D}，GitHub 约 94 stars；Semantic Scholar 约 47 次引用，接近当前引用阈值，因此这里作为动态编码路线的重要补写方法。

\textbf{实现细节：} 回顾上文，DaRePlane 在 HexPlane 三对平面乘法的框架内优化了每个 2D 平面的参数化方式（用方向感知小波系数替代灰度系数），但保留了将 4D 时空编码分解为六对 2D 平面的基本结构。Grid4D（2024.10） 则提出一个更根本的问题：平面分解本身的低秩假设是否合理？问题：平面分解的特征重叠。Grid4D 指出，平面方法将 4D 坐标 $(x,y,z,t)$ 分解为六对 2D 平面（如 $(x,y)\times(z,t)$、$(y,z)\times(x,t)$ 等），暗含的假设是 4D 变形特征可以被低秩因子化。但实际场景中，空间坐标大量重叠的高斯极为常见。这意味着即使 A 和 B 在时间轴上有完全不同的变形行为，它们在这些平面上的特征表示却无法区分。DaRePlane 的方向感知小波提升了每个平面内部的方向编码质量，但无法解决平面间特征重叠这一结构性问题——因为重叠的根源在于分解维度本身，而非平面系数的表示方式。动机与方案：从 2D 平面到 3D 网格。直觉上，既然 2D 平面分解导致了过多的维度丢弃（每次只保留两个坐标），一个自然的方向是提高分解粒度——保留更多维度。
\textbf{方法信息：} 代码已公开，仓库为 jiaweixu8/grid4d (94 stars)。入选依据：企业:否；代码:是；星标:否；引用:否。


% METHOD: 2411.11941
\paragraph{TimeFormer（2024.11）\cite{2411_11941}}
是对Deformable 3DGS时间建模方式的插件式改造。原始变形场通常独立查询每个时间戳，容易在剧烈运动和反射表面上产生不连续。TimeFormer把一批随机时间戳组成时间序列，把每个高斯的位置和时间编码送入跨时间Transformer，让自注意力显式学习多个时间之间的关系。它采用双流训练：TimeFormer流和基础变形流共享权重，训练时用Transformer提供全局运动知识，推理时移除该模块以保持速度。其本质改变是把时间一致性从局部正则变成跨时间注意力建模。
相对Deformable 3DGS，它没有替换规范高斯和变形场，而是在训练阶段给变形场增加跨时间关系教师。
\textbf{实现细节：} 工程瓶颈方面，FRoGS 系统分析了粗-细双变形使渲染速度减半、稀疏初始化点云造成远距运动伪影、暗场景前景阴影高斯引入色偏三个反复出现的瓶颈，并逐一给出针对性改进——用 Hadamard 早期融合将两次前向传播压缩为一次、用中值深度反向投影补充规范场锚点、用激进不透明度衰减抑制阴影高斯。TimeFormer 则从架构层面入手，借鉴 BatchFormer 在批次维度建模样本关系的思想，在变形场之前插入沿时间维度的 Transformer 自注意力，让每个时间戳的运动模式能"看到"其它时间戳的线索，并通过双流共享权重的训练-推理分离避免推理时引入额外开销。CDGS 把视角提升到容量预算层面，将动态重建形式化为"在固定高斯数量预算下最大化渲染质量"的约束优化，提出可微人口控制器（温度退火硬 Sigmoid 门控有效不透明度）和统一重要性分数（融合几何稳定性、运动显著性、感知影响），证明与其无约束增殖再剪枝不如让优化从一开始就在预算约束下学习。
\textbf{方法信息：} Tencent 参与了这篇论文。



% METHOD: 2501.03714
\paragraph{MoDec-GS（2025.01）\cite{2501_03714}}
MoDec-GS 把 Scaffold-GS 的锚点思想引入动态视频，目标是降低 4DGS 这类方法的存储。普通 4D-GS 在规范空间直接维护大量高斯，再为它们学习时空变形；MoDec-GS 改成先维护稀疏 anchor，每个 anchor 存 local context feature、offset 和尺度等信息，并通过神经高斯解码器生成附近 3D Gaussians。这样，静态 Scaffold-GS 中“少量锚点生成多个高斯”的压缩思想被带到动态重建里。

运动建模分成 Global-to-Local Motion Decomposition。第一层是 Global Anchor Deformation：小型全局 HexPlane 根据 anchor 位置和时间查询特征，再由多头解码器输出 anchor 位置、offset、local context feature 等属性的变化，把 Global Canonical Scaffold 变成某个时间段的 Local Canonical Scaffold。它主要处理人物整体移动、相机视角下的大范围场景变化等全局运动。第二层是 Local Gaussian Deformation：Local CS 先解码出具体高斯，再用共享局部 HexPlane 对这些高斯做目标时间的细粒度变形，处理局部表情、衣物或小物体运动。Temporal Interval Adjustment 在训练中动态调整每个 Local CS 覆盖的时间范围，使运动剧烈的片段获得更短、更专门的局部规范空间，而平稳片段可共享更长区间。相对 Swift4D 的静动态二分，MoDec-GS 更强调 anchor 层级压缩和全局/局部运动分摊。
\textbf{实现细节：} 前述 Incoherent Deformation, Not Capacity 从正则化角度解决变形不一致性——用 EER 施加局部平滑先验，用 GAD（损失率感知致密化）控制高斯增殖。但它并未改变变形场本身的表示结构：一个六面体平面仍然需要同时编码全局刚体运动和局部非刚性变形，容量瓶颈依然存在。MoDec-GS（2025.01） 采取了正交路线：不做正则化，而是将变形场拆成两级，让全局运动和局部运动由不同容量的表示分别承担。核心设计：两级变形替代单一六面体平面。现实视频中的运动通常是全局运动（如物体整体位移、摄像机运动）和局部运动（如手指弯曲、衣角飘动）的复杂组合。MoDec-GS 同样以锚点表示（Scaffold-GS）作为规范空间基础，但提出全局到局部运动分解（Global-to-Local Motion Decomposition, GLMD），将变形分为两个串行阶段。第一阶段是全局锚点变形（Global Anchor Deformation, GAD）：直接变形锚点本身的位置和属性，而非变形重建后的每个 3D 高斯。锚点 $v$ 的位置 $(x_v, y_v, z_v)$ 连同规范时间 $t_c$ 输入一个小型全局六面体平面 $H_G$，查询的特征经 MLP 解码后得到锚点属性偏移量（如局部上下文特征偏移 $\Delta f_v = \lambda_f[F_G(H_G(x_v, y_v, z_v, t_c))]$），直接加到锚点属性上。每个锚点还附加两个可学习的动态参数 $d_G$（全局）和 $d_L$（局部），$d_G$ 通过可微二值掩码 $M(d_G) = (\mathbb{I}[\sigma(d_G) > \tau] - \sigma(d_G)) \cdot \mathrm{sg}(\sigma(d_G)) + \sigma(d_G)$ 控制锚点是否参与全局运动，实现静态/动态区域的自动分离，无需预计算的外部掩码。消融实验表明，加入这对可学习动态参数对存储的增加几乎可以忽略，但能带来一致的质量提升。
\textbf{方法信息：} 代码：\url{https://github.com/skwak-kaist/modec-gs}。
\textbf{设计动机与边界：} 相对原始 3DGS 以独立高斯和图像重建误差为主的设定，MoDec-GS 把问题落在以下可实现的改动上：MoDec-GS采用基于变形的实时动态场景渲染方法，但使用锚点表示（Scaffold-GS）作为规范3DGS以获得紧凑性。GLMD分两阶段：（1）Global Anchor Deformation（GAD）：用微小全局hexplane变形锚点位置和属性，将Global CS变换为Local CS，高效捕捉全局运动。（2）Local Gaussian Deformation（LGD）：从Local CS通过神经高斯推导重建3DGs，用共享局部hexplane显式变形到每个目标时间戳。TIA在训练中动态重新平衡各Local CS的时间间隔，无需预计算光流或跟踪等外部信息。嵌入可学习参数到锚点以反映运动特征。其设计目的不是单独增加一个模块，而是用上述表示、监督或渲染约束消除原流程在目标任务中的具体失配。提出Global-to-Local Motion Decomposition（GLMD）以粗到细方式捕捉复杂动态运动。利用Global Canonical Scaffolds和Local Canonical Scaffolds扩展静态Scaffold表示。引入Temporal Interval Adjustment（TIA）自动控制每个Local CS的时间覆盖。模型大小平均减少70%。适用时还应保留论文给出的边界：两阶段运动建模增加训练复杂度；对极快速运动的捕捉能力受hexplane容量限制；TIA需要预设时间分段数量。
% METHOD: 2503.12307
\paragraph{Swift4D（2025.03）\cite{2503_12307}}
Swift4D 从 4D-GS 的效率瓶颈出发：真实视频里大部分区域是静态背景或近似静态物体，给所有高斯都配时空编码会把存储和训练时间浪费在不运动的基元上。它的第一步是动态-静态分解：在 canonical 3DGS 上给每个高斯增加一个可学习动态参数 $d$，再用训练视频的时间标准差生成 2D 动态/静态像素掩码。渲染时 $d$ 也按 alpha compositing 投到图像上形成 dynamic map，训练用二分类损失让投影到动态像素的高斯获得更高 $d$，投影到静态像素的高斯获得更低 $d$。

分解完成后，Swift4D 只让动态高斯进入多分辨率 4D Hash 变形器。该变形器以高斯位置和时间为查询，输出位置、旋转、尺度等时变残差；静态高斯则直接按 3DGS 参数渲染，不保存额外动态属性。最终图像由静态集合和变形后的动态集合共同 splat 得到。为了减少 canonical/deformation 耦合带来的 floaters，方法还加入 temporal pruning，删除在时间变形中贡献异常或不稳定的基元。相对 4D-GS 的“所有高斯都查询同一个时空结构”，Swift4D 的增量是把动态建模限制在真正需要变形的子集上，因此训练更快、模型更小；代码中的 \texttt{train\_dynamic}、dynamic map 渲染和 hash 变形训练也对应这个三阶段流程。
\textbf{实现细节：} 前述 DASH 和 Grid4D 都采用了动静分离策略来缓解 4D 哈希编码的碰撞与存储问题，但它们的动静分类方式各不相同。Swift4D（2025.03） 沿着动静分离的思路进一步推进，但其核心贡献不在于分离策略本身，而在于提出了一种基于渲染的动静分类机制——将分类问题嵌入 3DGS 的 $\alpha$-混合渲染管线中，利用多视角的自然遮挡关系自动纠正误判，使得分类过程既快速（约 1 分钟）又鲁棒。渲染式动静分类。Swift4D 为每个高斯点引入一个可学习的动态参数 $d$（初始化为 0），并通过 3DGS 标准的 $\alpha$-混合将其投影到 2D： 其中 $\sigma$ 为 Sigmoid 函数，将动态值映射到 $(0,1)$。这个设计的直觉在于：如果一个高斯覆盖了动态像素，反向传播会赋予其正梯度（推动 $d$ 增大）；如果覆盖了静态像素，则赋予负梯度（推动 $d$ 减小）。由于 $\alpha$-混合本身包含深度排序和透明度累积，梯度会随着遮挡关系和距离自然衰减——远处的误判高斯不会获得强梯度。监督信号来自训练视频的 2D 动静像素掩码，通过对每个像素计算时间标准差 $S(\mathbf{x}) = \sqrt{\frac{1}{T}\sum_{t=1}^{T}(C(\mathbf{x},t) - \bar{C}(\mathbf{x}))^{}2}$ 并以阈值 $\tau = 0.02$ 二值化得到。整个优化通过二值交叉熵损失 $\mathcal{L}_d = \mathbb{E}_\mathbf{x}[-\hat{D}(\mathbf{x})\log D(\mathbf{x}) - (1-\hat{D}(\mathbf{x}))\log(1-D(\mathbf{x}))]$ 驱动，将动静分类转化为可微的二元分类问题。这种方法的一个关键优势是多视角一致性：即使某个视角下静态高斯被动态像素错误覆盖（因而获得正向梯度），其他视角中同一高斯覆盖静态像素时会产生更强的负向梯度，最终纠正分类——这本质上是利用了多视角之间的互补遮挡信息。优化完成后（约 1 分钟），$d > \tau_d = 7.0$ 的高斯被分类为动态，其余为静态。4D 哈希编码与变形解码。动静分离后，Swift4D 仅对动态高斯施加 4D 多分辨率哈希编码 $h_{4d}$。
\textbf{方法信息：} 代码：\url{https://github.com/wujh2001/swift4d}；GitHub 关注度较高，约 153 stars。



% METHOD: 2506.05480
\paragraph{ODE-GS（2025.06）\cite{2506_05480}}
ODE-GS 解决的是动态 3DGS 的“未来外推”而不是插值。4D-GS、Deformable 3DGS 和很多 HexPlane 方法都以时间戳 $t$ 为条件查询变形场；在观测窗口内，这种做法能补帧，但把 $t$ 推到训练范围之外时，网络没有学过“系统如何继续演化”。ODE-GS 因此先把重建和预测拆开：第一阶段仍训练 canonical Gaussians 加 time-conditioned deformation MLP，只负责在已观测时间内生成可靠轨迹；第二阶段冻结这套插值模型，把它当作轨迹生成器，密集采样每个高斯的位置、旋转和尺度序列。

外推模型把每个高斯过去一段轨迹送入 Transformer encoder，得到表示过去运动的 latent state；随后用 neural ODE 的速度场 $dz/dt=f(z)$ 在连续潜在空间中积分未来状态，再由 decoder 输出未来的高斯位置、旋转和尺度。颜色和不透明度保持不变，减少预测自由度，训练时则用 prefix-to-future sampling 让模型看不同长度的过去并预测不同 horizon，同时加入 latent smoothness、Gaussian trajectory smoothness 和 adaptive regularization。代码仓库也把流程拆成 \texttt{train\_interpolation.py} 和 \texttt{train\_extrapolation.py} 两步，对应论文中的“先重建、再学习动力系统”。相对 4D-GS，它的核心变化是把时间模块从“给定时间查询变形”改成“从过去轨迹演化未来高斯状态”。
\textbf{实现细节：} ODE-GS（2025.06） （Latent ODEs for Dynamic Scene Extrapolation with 3D Gaussian Splatting，ICLR 2026）由 Yale、TU Delft、LSU 和 RTX 联合提出，引入潜在神经常微分方程实现动态 3D 场景的未来外推。已有动态场景重建方法（如 Deformable 3D Gaussians、4DGS）依赖时间条件的变形网络，本质上是在观测时间窗口内做插值，当时间戳超出训练分布时会产生 OOD 失效。ODE-GS 的核心思路是将动态场景外推重新表述为序列到序列预测问题，从而消除对显式时间戳的依赖。解耦的两阶段架构。方法分两阶段：首先训练一个标准的规范帧加变形网络插值模型，在观测窗口内生成高斯参数的密集时序轨迹，训练完成后冻结该模型作为数据生成器；然后在生成的轨迹上训练 Transformer-Latent ODE 外推模块。形式化地，连续时间渲染算子分解为 $\mathcal{F}(t,V) = \mathcal{R}(\mathcal{G}(t), V)$，其中高斯状态在观测窗口内由规范高斯加变形网络生成（$\mathcal{G}(t) = \overline{\mathcal{G}} + \mathcal{D}_\omega(t, \overline{\mathcal{G}})$），在未来时刻由外推模块生成（$\mathcal{G}(t) = \mathcal{E}_\phi(\gamma, t)$）。每个规范高斯 $G_k = (\mu_k, q_k, s_k, c_k, \alpha_k)$ 中仅位置 $\mu_k \in \mathbb{R}^{}3$、旋转 $q_k \in \mathbb{R}^{}4$（四元数）和缩放 $s_k \in \mathbb{R}^{}3$ 是时间相关的，不透明度和球谐系数跨时间保持一致。Transformer-Latent ODE 外推模块。外推模块 $\mathcal{E}_\phi$ 将时序预测建模为序列到序列映射。
\textbf{方法信息：} 代码：\url{https://github.com/preacherwhite/ode-gs}。
\textbf{设计动机与边界：} 相对原始 3DGS 以独立高斯和图像重建误差为主的设定，ODE-GS 把问题落在以下可实现的改动上：ODE-GS三阶段：(1)插值模型：使用canonical-plus-deformation策略，学习规范高斯集和时间条件变形MLP，在观测窗口内生成密集高斯参数轨迹；(2)外推模型：Transformer编码器聚合过去轨迹为潜在状态，通过神经ODE在连续潜在空间中演化，数值积分产生平滑的未来轨迹；(3)动态采样策略：训练时随机采样不同预测时域。仅优化位置、旋转和缩放的时间依赖参数，保持不透明度和颜色不变。其设计目的不是单独增加一个模块，而是用上述表示、监督或渲染约束消除原流程在目标任务中的具体失配。提出ODE-GS，将3DGS与潜在神经ODE结合，首次实现动态3D场景的未来外推。消除时间戳依赖，将高斯参数轨迹建模为连续时间潜在动态。在D-NeRF、NVFi和HyperNeRF基准上达到SOTA外推性能，指标提升19.8%。适用时还应保留论文给出的边界：外推范围受限于运动可预测性；对高度非线性的突然运动效果有限；需要观测窗口内的充分数据。
% METHOD: 2507.19141
\paragraph{DASH（2025.07）\cite{2507_19141}}
DASH 继承 4D-GS/Grid4D 对时空编码的讨论。平面分解方法把 4D 特征拆成若干 2D plane，存储和查询都轻，但低秩假设会让复杂运动的特征互相重叠；直接把完整 4D hash 用到全场景又会把大量静态背景也塞进哈希表，带来碰撞和冗余。DASH 的改变是先自监督地区分静态与动态，再只给动态高斯使用多分辨率 4D hash deformation，静态部分仍按普通 3DGS 优化。

自监督分解阶段先用 3D hash 特征和小 MLP 为高斯预测线性运动描述量 $\Delta p$。静态高斯在优化中应收敛到 $\Delta p\approx 0$，动态高斯则会表现出更大的运动幅度；方法据此按运动幅度阈值划分两组，并对静态组加入 static constraint，避免它们被误拉进动态哈希。动态组随后查询 multiresolution 4D hash encoder，输出位置等变形残差；静态组不查询时间哈希，直接参与渲染。为抑制显式 hash deformation 的抖动，训练中加入时空平滑正则，并且静态/动态两组使用独立的 density control。相对 Swift4D，它同样避免“所有高斯都动态化”，但动态子集的底层表示从 4D Hash/HashGrid 而非 HexPlane 取得特征；相对 Grid4D，它减少了全局 hash 的碰撞来源。
\textbf{实现细节：} 从平面分解到哈希编码的范式转移。前述 DaRePlane 在平面分解（HexPlane）框架内改进平面参数化——用方向感知小波系数替代无方向性的空间域系数——但整体架构仍基于低秩假设：将 4D 时空体分解为三对 2D 平面的乘积，认为变形特征具有大量共性、可被因子化为紧凑表示。DASH（2025.07） 指出，这一低秩假设在面对空间坐标大量重叠的高斯时会产生特征重叠（feature overlap）：多个高斯共享同一平面网格单元，各自的变形模式被混杂编码，直接限制变形预测的区分度和渲染质量。Grid4D 虽将 4D 编码拆分为四个独立的 3D 哈希编码以回避低秩假设，但各 3D 哈希表之间仍然存在特征重叠，且时间信息被分散编码而非统一建模。DASH 的根本思路是：放弃平面分解，改用多分辨率 4D 哈希编码直接在统一的时空框架中表征动态高斯的变形，同时通过自监督动静分解将 4D 哈希编码的范围限制在动态区域，从而兼顾表达力和效率。问题：为什么不能直接用全场景 4D 哈希编码？多分辨率哈希编码（Instant-NGP）在静态 3D 场景中已证明高效，但直接扩展到 4D 会带来 $O(n^{}4)$ 的空间复杂度——时间维度的加入使哈希表的碰撞率急剧上升，全场景编码需要远更大的哈希表才能维持足够的特征分辨率。DASH 的核心洞察是：大多数场景包含大面积静态区域，这些区域完全不需要存储在 4D 哈希表中——静态高斯只需 3D 空间编码即可，仅真正运动的部分才需要 4D 时空编码。自监督动静分解。DASH 提出一种无需额外标注（如手动掩码或预计算分割）的动静分解策略，利用静态和动态高斯在优化过程中的固有行为差异进行区分。具体分三步：。
\textbf{方法信息：} 官方代码：\url{https://github.com/chenj02/DASH}；仓库约 27 stars，属于完整公开但关注度仍在积累的新方法。



% METHOD: 2601.05584
\paragraph{GS-DMSR（2026.01）\cite{2601_05584}}
GS-DMSR 从优化效率角度改造 Deformable 3DGS。普通动态 3DGS 会让每个高斯在每个阶段都参与变形场更新，即使其中相当一部分只是静态背景或极弱运动区域；结果是梯度和计算被低价值基元消耗。GS-DMSR 的第一步是估计每个高斯属性随时间变化的 motion saliency coefficient，把高显著性和低显著性基元分开。高显著性高斯继续通过 deformation field 迭代精修，低显著性高斯则冻结或跳过后续变形更新，使训练把梯度集中到真正运动的区域。

第二个改动是多尺度解耦流形变形。方法不使用一个 MLP 同时输出所有变形，而是设置多头解码器 $D=(D_x,D_r,D_s)$，分别预测位置、旋转和尺度残差；低频几何结构通过 coarse manifold feature 保持整体运动趋势，高频残差则补充局部细节。这样做的目标不是提出新的 rasterizer，而是改变动态 3DGS 的训练分配：先判断哪些高斯值得继续变形，再用多尺度流形特征把不同属性的运动分开优化。相对 Deformable 3DGS，它的增量在于“运动显著性驱动的梯度聚焦 + 解耦多尺度变形头”。
\textbf{方法信息：} 代码：\url{https://github.com/DustSettled/GS-DMSR}。
\textbf{设计动机与边界：} 相对原始 3DGS 以独立高斯和图像重建误差为主的设定，GS-DMSR 把问题落在以下可实现的改动上：核心方法：(1)运动显著性驱动的动态Gaussian优化——定量分析Gaussian属性的动态演化过程，通过参数化运动显著性系数量化Gaussian运动状态变化，使用预设显著性阈值将Gaussian分为高显著性和低显著性两类。高显著性Gaussian通过变形场迭代精化，低显著性Gaussian免于后续变形场更新，实现自适应梯度聚焦，显著提升训练效率；(2)多尺度解耦流形变形——采用多头Gaussian变形解码器D=(D\_x, D\_r, D\_s)，三个独立MLP分别预测位置、旋转和缩放变形。引入多尺度流形增强模块，通过动态特征维度适配、跨尺度特征金字塔构建（集成低频几何结构L\_c和高频细节H）和非线性流形映射三层操作，增强复杂变形建模能力。最终变形Gaussian G'=(x', s', r', alpha, C)用于渲染。其设计目的不是单独增加一个模块，而是用上述表示、监督或渲染约束消除原流程在目标任务中的具体失配。提出GS-DMSR，通过运动显著性驱动的动态梯度更新机制实现自适应梯度聚焦，结合多尺度流形增强模块协同优化隐式解码器和显式变形场，在合成数据集上达96FPS，同时减少存储开销和训练时间。适用时还应保留论文给出的边界：运动显著性阈值需要手动设定；多尺度流形增强模块增加了模型复杂度。
% METHOD: 2601.05138
\paragraph{VerseCrafter（4D 几何控制的视频世界模型）（2026.01）}
VerseCrafter~\cite{2601_05138} 把动态视频世界模型的控制变量从 2D 轨迹提升到 4D geometric world state。普通 camera-controlled video model 可以控制相机，但多物体运动、遮挡和三维占据关系仍然被压在投影后的 2D 图像里；VerseCrafter 的表示由静态背景点云和逐物体 3D Gaussian trajectories 组成，每个物体的轨迹同时表达运动路径和随时间变化的概率式 3D occupancy。这样，相机运动和物体运动可以在同一三维坐标系里描述。
实现流程先把输入场景整理成 4D Geometric Control：背景用 point cloud 保持静态结构，动态物体用一串带时间索引的 3D Gaussians 表达位置、尺度、占据和运动。随后把这些 4D 控制量渲染或投影成视频扩散模型可读取的 control maps，再由预训练视频生成模型合成动态真实视频。训练目标不是优化一个可直接渲染的 4DGS 场景，而是让视频输出服从给定的几何控制：相机轨迹、物体运动和遮挡关系都要与 4D control 一致。相对本章传统 4DGS 重建方法，VerseCrafter 更偏“用 3D Gaussian 轨迹控制视频世界模型”，但它直接使用高斯轨迹作为类别无关的动态原语，因此纳入动态 3DGS 谱系。
\textbf{方法信息：} 作为 2026 年扩展召回条目纳入；当前本地证据未确认完整公开代码。
\textbf{设计动机与边界：} 相对原始 3DGS 以独立高斯和图像重建误差为主的设定，VerseCrafter 把问题落在以下可实现的改动上：方法把世界状态表示为静态背景点云加逐物体 3D Gaussian trajectories，轨迹携带运动路径和概率式 3D 占据；再将该 4D 几何控制渲染为控制图，输入预训练视频扩散模型生成动态、视角一致的视频。它不是标准 3DGS 重建管线，但直接使用 3D Gaussian 轨迹作为几何控制原语。其设计目的不是单独增加一个模块，而是用上述表示、监督或渲染约束消除原流程在目标任务中的具体失配。把 3D Gaussian 轨迹作为类别无关的 4D 控制表示，用于统一相机运动和多物体运动控制。适用时还应保留论文给出的边界：输出目标主要是视频生成控制而非可渲染 3DGS 场景；全文指标尚未逐表核验。
% METHOD: 2603.11543
\paragraph{Mango-GS（2026.03）\cite{2603_11543}}
Mango-GS 针对的是 Deformable 3DGS 和 4D-GS 中常见的逐时间查询问题：模型每次只看一个时间戳，容易把每一帧当成独立状态记住，而不是学习稳定的运动趋势。它继承“规范高斯 + 运动传播”的思路，但不直接让变形 MLP 对所有高斯逐帧预测位移，而是在规范空间放置少量控制节点，用这些节点作为运动骨架来驱动密集高斯。

每个控制节点由两部分组成：规范位置 $p_i$ 提供稳定空间锚点，可学习特征 $f_i$ 提供更持久的语义/结构身份。高斯和节点的对应关系不只看欧氏距离，而是在位置和特征组成的联合空间里做 $k$-NN，并在规范空间预计算权重。这样，当人体、手臂或物体发生大幅运动时，原本空间上临近但语义上不相关的区域不容易被同一个节点错误牵引。随后，时间 Transformer 在一个短窗口内同时预测所有节点的平移、旋转和缩放轨迹，再通过预计算权重把节点运动混合传播到每个高斯。

训练上，Mango-GS 使用端到端渲染损失，同时加入输入遮蔽策略和多帧损失，让模型在遮挡帧、快速运动帧和相邻帧之间都保持一致。与 4D-GS 的差别在于，时间模块不再只是“给定 $t$ 查询一次变形”，而是一次观察多帧并输出一段节点轨迹；与 SC-GS 式节点传播相比，它用解耦节点和 Transformer 解决空间邻域漂移。论文在 HyperNeRF 和 Neural 3D Video 上报告了较好的时序一致性和实时渲染速度。
\textbf{实现细节：} Mango-GS（2026.03） （Multi-Frame Node-Guided 4D Gaussian Splatting）指出已有 4D-GS 方法（如 Deformable 3DGS、SC-GS）采用逐帧变形策略：对每一帧独立预测高斯偏移量，模型倾向于记忆各帧的瞬时状态而非学习底层运动规律，导致快速运动场景中出现模糊和闪烁。虽然时序 Transformer 天然适合捕获多帧间的运动趋势，但典型场景包含百万级高斯，直接对所有高斯施加时序注意力的计算和显存代价过高。Mango-GS 的核心思路是将百万级高斯的运动建模分解为稀疏控制节点的运动建模——用少量节点（$M \ll N$，默认 $M=2048$）捕获场景运动骨架，再通过加权和传播到全部高斯。方法引入两项关键技术。首先是解耦控制节点表示（decoupled control node）。此前 SC-GS 仅用规范空间中的欧氏距离做 k-NN 来确定节点-高斯的对应关系，但当运动幅度较大时，初始帧中邻近的点可能被带到不同运动区域，静态空间邻域不再有效（邻域漂移）。Mango-GS 将每个控制节点 $n_i$ 解耦为规范位置 $p_i \in \mathbb{R}^{}3$ 和可学习特征码 $f_i \in \mathbb{R}^D$ 两部分。对于每个高斯 $g_j$，用轻量级 MLP 将其规范参数 $\phi_j = (x_j, q_j, s_j)$（位置、旋转、缩放）映射到嵌入空间，再与节点特征码比较计算亲和度分数，取 top-$k$ 并经 softmax 归一化得到影响权重： \begin{equation*} w\_{ij} = \frac{\exp(-\mathcal{D}(g\_j, n\_i))}{\sum\_{i' \in \mathcal{K}(j)} \exp(-\mathcal{D}(g\_j, n\_{i'}))}, \quad \forall i \in \mathcal{K}(j) \end{equation*} 其中 $\mathcal{D}(\cdot)$ 衡量高斯规范参数与节点特征码在联合位置-特征空间中的距离，$\mathcal{K}(j)$ 为 $k$ 近邻节点集合（默认 $k=3$）。这种学习式亲和度产生的邻域是语义感知的——即使两点在空间上已分离，只要属于同一运动模式仍可保持关联。此外，该设计还具有参数效率优势：直接复用高斯的规范属性（位置、旋转、缩放）计算嵌入，无需为每个高斯额外引入潜在特征，避免了额外的显存开销。其次是多帧时序注意力网络。
\textbf{方法信息：} 代码：\url{https://github.com/htx0601/mango-gs}；仓库目前约 7 stars，属于有完整代码但关注度还在积累的新方法。



% METHOD: 2603.25265
\paragraph{ViewSplat（2026.03）\cite{2603_25265}}
ViewSplat 属于 pose-free 前馈 3DGS 路线，直接承接 SPFSplat/SPFSplatV2。NoPoSplat、SPFSplat 这类方法通常一次前向预测一套“对所有目标视角都固定”的 Gaussian map；问题是复杂几何、反射和局部遮挡下，一次预测出的中心、尺度、颜色很难同时满足所有视角。ViewSplat 的改动是让基础高斯在渲染某个目标视角前再做一次轻量、视角相关的属性校正，而不是把所有误差都压给一个静态 Gaussian set。

具体流程是：先由 pose-free 网络从输入图像预测基础 3D Gaussians 和相机/深度相关表示；再把目标相机相对高斯的单位方向和对数距离组织成 4D target-pose feature。视图自适应头不直接输出每个高斯的最终属性，而是预测一组动态 MLP 权重；这些小 MLP 接收 target-pose feature，为位置、尺度、旋转、不透明度和 SH/颜色产生残差。渲染时先把基础高斯加上这些残差，再用普通 3DGS rasterizer 得到目标视角。训练损失由 MSE+LPIPS 的图像重建项和重投影项组成，后者约束预测高斯中心与像素投影保持一致，避免视角残差把几何拉偏。相对 SPFSplat，它的关键不是换骨干，而是在“前馈预测”和“目标视角渲染”之间插入 view-conditioned Gaussian adaptation。
\textbf{实现细节：} 前馈与稀疏视角方面，核心挑战是把训练范式从桌面级合成场景推向真实室外数据并放宽输入假设。Flux4D 把目标场景转向真实室外驾驶并彻底消除对标注和预训练模型的依赖；Sparse4DGS 针对 5 FPS 以下稀疏帧提出纹理感知优化策略，主动引导高斯聚焦纹理丰富区域；SV-GS 引入骨架结构低维先验弥补极端稀疏时间观测；MoSca 则完全走另一条路，用 DPT/CoTracker/RAFT 等预训练视觉基础模型从随意拍摄单目视频全自动恢复 4D 场景。在位姿与视角假设方面，UFO-4D 实现从两张未定位图像的前馈 4D 重建，No Pose No Problem 进一步去除位姿先验需求，ViewSplat 提出视角自适应前馈动态高斯合成。
\textbf{方法信息：} 代码：\url{https://github.com/cvlab-uos/ViewSplat}。
\textbf{设计动机与边界：} 相对原始 3DGS 以独立高斯和图像重建误差为主的设定，ViewSplat 把问题落在以下可实现的改动上：ViewSplat的核心创新是将静态高斯回归转变为视图自适应动态泼溅：(1)基础高斯预测：基于SPFSplat架构，从无位姿图像预测标准3D高斯参数。(2)视图自适应头：预测动态MLP的权重，这些MLP接收4D目标位姿特征（3D单位视角方向+1D log距离）作为输入，为每个高斯属性（3D位置、尺度、旋转、不透明度、颜色）预测视图相关的残差更新。(3)训练目标：结合渲染损失（MSE+LPIPS）和重投影损失，后者通过惩罚2D像素坐标和投影高斯中心距离来强制几何约束。Log距离变换保持近场变化的更高敏感性。推理17FPS，渲染154FPS。其设计目的不是单独增加一个模块，而是用上述表示、监督或渲染约束消除原流程在目标任务中的具体失配。提出ViewSplat，将范式从静态原语回归转变为视图自适应动态泼溅。学习视图自适应的潜在表示，动态MLP根据目标视角坐标预测逐属性的残差更新。实现SOTA保真度，保持17FPS推理和154FPS渲染速度。适用时还应保留论文给出的边界：动态MLP增加推理开销；视图自适应机制可能在极端视角失效；依赖3D基础模型。
% METHOD: 2603.28064
\paragraph{4DSurf（2026.03）\cite{2603_28064}}
把动态3DGS用于表面重建，而不只追求新视角颜色。它仍从规范高斯出发，但为每个高斯预测速度、角速度和膨胀速度，用这些量描述从规范形状到目标时间的场景流；随后把SDF流正则接到高斯运动上：先由时间相关SDF场估计表面法线和有符号距离变化，再要求高斯运动诱导出的表面变化与SDF流一致。这样，优化目标不再只是“图像看起来对”，还要求运动后的高斯支撑一个连续、时间一致的表面。面对人体交互这类大变形序列，4DSurf把长视频切成重叠片段，每段只处理较小形变，并通过重叠时间步传递几何信息；后续增量运动微调则在控制存储的同时修正片段间误差。相对于颜色驱动的Deformable 3DGS，4DSurf的核心增量是把表面连续性和场景流几何约束直接接入动态高斯优化。
\textbf{实现细节：} 此外，4DSurf（2026.03） 实现高保真动态场景表面重建，DyGASR（2024.11） 通过广义指数泼溅和表面对齐提升动态几何精度。运动建模与扩展应用。本节介绍三条技术路线：（1）运动建模与操纵；（2）分析与概率方法；（3）其他变形与4D方法。
\textbf{实现补充：} 实现上还要核对输入、输出、参数、特征、初始化、训练损失、正则、渲染或推理接口，以及优化中的剪枝、分裂、采样、编码或解码步骤；这些抓手决定该方法到底改了 3DGS 管线的哪一环。
\textbf{技术实现补充：} 此外，4DSurf（2026.03） 实现高保真动态场景表面重建，DyGASR（2024.11） 通过广义指数泼溅和表面对齐提升动态几何精度。运动建模与扩展应用。本节介绍三条技术路线：（1）运动建模与操纵；（2）分析与概率方法；（3）其他变形与4D方法。
\textbf{方法信息：} NVIDIA、Amazon 参与了这篇论文。
\textbf{设计动机与边界：} 相对原始 3DGS 以独立高斯和图像重建误差为主的设定，\textit{4DSurf} 把问题落在以下可实现的改动上：4DSurf方法包含：(1)高斯速度场：预测每个规范高斯的速度v、角速度omega和膨胀速度delta，描述从规范形状到任意时间步的变形，显式建模场景流。(2)SDF流正则化：将点级SDF流扩展到高斯基元，推导高斯运动与诱导SDF变化的关系。通过时间依赖SDF场的表面网格估计法线，计算SDF流作为有符号距离的时间导数，约束高斯运动与表面演化一致。(3)重叠分段划分：将序列分成重叠段，每段包含小变形，通过共享重叠时间步在段间传递几何信息，解决大变形导致的表面抖动。(4)增量运动微调：最小化存储的同时保持竞争力性能。其设计目的不是单独增加一个模块，而是用上述表示、监督或渲染约束消除原流程在目标任务中的具体失配。提出无需先验的通用动态表面重建框架4DSurf，可处理多物体和大变形场景。核心创新是引入高斯变形诱导的SDF流正则化，约束高斯运动与演化表面对齐。设计重叠分段划分策略将长序列分成小变形段并通过共享重叠时间步传递几何信息，在Hi4D和CMU Panoptic数据集上Chamfer距离分别降低49%和19%。适用时还应保留论文给出的边界：依赖时间依赖的SDF场估计，重叠分段增加了优化复杂度，对初始化质量敏感。

\section{单目动态场景重建}

单目视频是动态3DGS最困难的输入形式之一：每一帧只看到一个视角，动态物体的深度、遮挡和运动同时未知。这里的主要技术路线，是在Deformable 3DGS或MoSca式规范变形之外，引入深度、光流、点跟踪、扩散模型、模糊模型或前馈先验，把欠约束的4D优化变成可收敛的问题。
% METHOD: 2309.13101
\paragraph{Deformable 3D Gaussians（规范高斯与时间条件变形）（2023.09）}
Deformable 3D Gaussians~\cite{2309_13101} 解决单目动态视频中逐帧 3DGS 互不一致的问题。它在规范空间保存一组具有中心、四元数旋转、尺度、opacity 与 SH 外观的高斯；给定该中心和时间 $t$，纯 MLP deformation field 输出位置、旋转和尺度的增量。增量后的高斯直接送入原有 tile-based Gaussian rasterizer，因此外观基元和时间运动分别由规范高斯与网络承担。

\textbf{训练路径。} 输入是带时间标签、SfM 相机位姿及其稀疏点云的单目视频，点云初始化规范高斯。每次迭代按图像时间查询 MLP，将 $G(x,r,s,\alpha)$ 变成 $G(x+\delta x,r+\delta r,s+\delta s,\alpha)$，渲染并以图像重建误差联合更新高斯与 MLP；原 3DGS 的密度控制仍用于补足表示。论文选择纯 MLP 而非低秩 grid/plane 变形编码，理由是复杂动态运动不必服从低秩时空分解。

\textbf{AST。} 为使单目监督下的时间映射先平滑、后细化，Annealing Smooth Training 在早期扰动/平滑时间编码并逐步退火。网络先被迫在邻近时间产生相近变形，随后才拟合快速细节，从而减少逐帧过拟合造成的抖动。它与 Dynamic 3D Gaussians 的核心差异是运动由共享连续函数生成而非为每粒子显式存一条时间轨迹；这带来任意时间查询与更紧凑的参数共享，也把未观测运动和规范空间的稳定性押在 MLP 及单目先验上。

% METHOD: 2404.14410
\paragraph{Guess The Unseen（扩散补全的多人单目动态 Gaussian 场景）（2024.04）}
Guess The Unseen~\cite{2404_14410} 同时重建静态世界和多个动态人物，并专门处理裁剪、遮挡或少帧可见造成的人体未观测区域。背景由世界坐标中的高斯表示；每个人在 SMPL 定义的 canonical pose 中持有一组高斯，再按各时刻人体姿态变换并与背景共同 splat。

核心不是把缺失区域当作零监督，而是在规范空间融合不同时间的稀疏可见线索，并用预训练二维扩散模型合成与观测外观一致的未见视角；这些补全视图与真实观测共同监督高斯。统一的显式表示使得移除某人的高斯集合或替换其姿态即可编辑场景。该能力依赖 SMPL 拟合和扩散补全的可信度，未见细节是先验生成而非由单目视频唯一确定。

% METHOD: 2406.00434
\paragraph{MoDGS（深度先验引导的随手单目动态重建）（2024.06）}
MoDGS~\cite{2406_00434} 处理相机静止或缓慢移动的单目视频：这类输入缺乏通过相机基线获得多视图几何的条件，普通 canonical deformation 容易学到退化几何。方法先对每帧引入单目深度先验，并使用深度和二维 motion cue 构造三维感知的 deformation 初始化，而非从随机/二维对齐直接学习。

规范空间高斯经可逆时间条件变形映射到当前帧，再以标准 splatting 重建图像。训练中的 robust depth loss 不强迫不同帧深度估计有绝对一致尺度，而是以鲁棒方式将其作为几何约束；图像误差与深度项共同更新高斯和变形场。深度先验扩展了慢相机视频可恢复的范围，但估计深度的失真、严重遮挡和大幅非刚性仍可能传递到 4D 几何。

% METHOD: 2406.18717
\paragraph{Dynamic Gaussian Marbles（单目约束下的各向同性动态粒子）（2024.06）}
Dynamic Gaussian Marbles~\cite{2406_18717} 针对随手单目视频中动态高斯高度欠约束的问题。它将各向异性高斯简化为方向固定、尺度为标量的 isotropic ``marbles''，减少单个粒子可通过形状伪造几何/运动的自由度；目标是优先恢复全局连贯轨迹，而不是立即追求最大局部表达力。

优化采用层次化 divide-and-conquer：先在粗时间/空间尺度上建立全局运动，再逐步细化；图像级先验与几何级先验共同限制表示，点跟踪结果提供跨帧轨迹损失。当前时刻的 marble 仍按普通 Gaussian 投影与合成渲染。它用较强模型约束换取单目可识别性，细小各向异性结构和未被跟踪覆盖的区域则可能受限。

% METHOD: 2411.08879
\paragraph{4D Gaussian Splatting in the Wild with Uncertainty-Aware Regularization（2024.11）\cite{2411_08879}}
处理手持单目视频中最常见的过拟合：训练视角看起来好，新视角在观测不足区域崩坏。它用每个高斯对训练图像的贡献估计不确定性，形成2D不确定性图；只在不确定区域施加扩散模型SDS先验和深度平滑，而不是对整张图统一正则。这样做的关键是保留已观测区域的重建精度，同时把生成先验集中到缺少多视角约束的位置。对于SfM难以初始化的快速运动区域，方法再利用估计深度图和场景流额外密集化高斯，使动态物体不再完全依赖静态SfM点云。相对4D-GS，它的改变是让正则化和密集化都按观测可靠性自适应，而非全局加同一项约束。
\textbf{实现细节：} 4DGS in the Wild（2024.11） 针对手持单目相机随意拍摄的真实视频，首次关注了自然条件下的单目动态重建问题。与此前方法多假设受控采集环境（如 HyperNeRF 使用不切实际的训练-测试划分，将同一视频轨迹中的帧同时用于训练和测试，使得任务更接近视频插值而非真正的新视角合成）不同，该方法面向的是相机围绕快速运动物体拍摄的真实手持视频场景。在这种场景中，4DGS 模型容易过拟合训练帧，而朴素地施加正则化虽能改善测试视角的渲染质量，却会牺牲训练图像的重建精度——二者之间存在根本性的平衡难题（trade-off）。
\textbf{实现补充：} 实现上还要核对输入、输出、参数、特征、初始化、训练损失、正则、渲染或推理接口，以及优化中的剪枝、分裂、采样、编码或解码步骤；这些抓手决定该方法到底改了 3DGS 管线的哪一环。
\textbf{方法信息：} Samsung 参与了这篇论文；当前未确认该论文完整官方代码。
\textbf{设计动机与边界：} 相对原始 3DGS 以独立高斯和图像重建误差为主的设定，4D Gaussian Splatting in the Wild with Uncertainty-Aware Regularization 把问题落在以下可实现的改动上：方法包含两个核心创新：(1)不确定性感知正则化：量化每个高斯基元对训练图像渲染的贡献作为不确定性度量，构建2D不确定性图，在不确定区域（观测不足）选择性施加扩散模型先验（SDS损失）和深度平滑约束。这解决了传统正则化在提升新视角质量时牺牲训练图像重建精度的问题；(2)动态区域密集化：针对SfM在快速运动区域无法提供可靠3D标志点的问题，利用估计深度图和场景流在动态区域初始化高斯基元。首次聚焦于自然条件下手持单目相机拍摄的日常视频。其设计目的不是单独增加一个模块，而是用上述表示、监督或渲染约束消除原流程在目标任务中的具体失配。首次针对日常单目视频的4DGS重建提出不确定性感知正则化，在不确定区域选择性施加先验约束，平衡训练重建和新视角合成质量；提出动态区域密集化方法解决SfM初始化失败问题。适用时还应保留论文给出的边界：不确定性估计和扩散模型先验增加计算开销；对极快运动仍有挑战
% METHOD: 2412.03526
\paragraph{Feed-Forward Bullet-Time Reconstruction（2024.12）\cite{2412_03526}}
把单目动态重建改造成前馈子弹时间预测。模型把上下文视频帧切成patch，加入相机Plucker嵌入和时间嵌入，经ViT处理后为每个输出token解码高斯颜色、尺度、旋转、不透明度和射线距离，3D位置由像素对齐反投影得到。Novel Time Enhancer会合成中间时间帧作为额外上下文，使模型能查询“冻结时间、移动相机”的子弹时间视图。相对于逐场景优化的Deformable 3DGS，它把动态重建能力转移到大规模训练出的前馈网络中。
\textbf{实现细节：} BTimer（2024.12） 是首个运动感知前馈动态场景重建模型，其核心创新在于"子弹时间"（bullet-time）公式：给定一组上下文帧及其相机位姿和时间戳，模型在指定的子弹时间戳 $t_b$ 重建完整场景的 3D 高斯 Splatting 表示——即将场景"冻结"在 $t_b$ 时刻的 3D 状态。迭代设置 $t_b$ 为视频中每个时间戳即可高效并行重建完整 4D 视频。关键设计是每个输入帧携带两种时间嵌入：上下文时间嵌入 $\tau_i^{\mathrm{ctx}}$（编码该帧自身的时间戳 $t_i$）和子弹时间嵌入 $\tau_i^{\mathrm{bullet}}$（编码目标重建时间 $t_b$，所有上下文帧共享）。两种时间戳先经标准正弦位置编码（sinusoidal positional encoding），再分别通过两层线性层获得对应特征。两者相加形成最终时间特征 $\tau_i^{\mathrm{time}} = \tau_i^{\mathrm{ctx}} + \tau_i^{\mathrm{bullet}}$，使模型在处理每帧时同时知道"这是什么时候的"和"要重建什么时候的"，从而学会根据目标时间戳聚合上下文信息——动态区域在 $t_b$ 附近的帧被赋予更高权重，远离 $t_b$ 的帧提供静态上下文。这一公式天然统一了静态和动态重建：当所有时间戳相等时即退化为静态场景。模型架构基于 24 层 ViT Transformer，首尾各附加 LayerNorm。
\textbf{方法信息：} NVIDIA 参与了这篇论文。
\textbf{设计动机与边界：} 相对原始 3DGS 以独立高斯和图像重建误差为主的设定，Feed-Forward Bullet-Time Reconstruction of Dynamic Scenes from Monocular Videos 把问题落在以下可实现的改动上：核心设计：(1) 子弹时间重建：基于ViT骨干网络(24个自注意力块)，将输入上下文帧分割为8x8 patches，结合相机Plucker嵌入和时间嵌入形成输入token。通过Transformer处理后，每个输出token解码为3DGS参数(颜色、尺度、旋转、不透明度、射线距离共12参数)。3D位置通过像素对齐反投影获得。(2) Novel Time Enhancer(NTE)：合成中间时间戳的插值帧作为额外上下文视图。(3) 学习课程：精心设计混合静态和动态场景数据集的训练课程，增强运动感知和时间一致性。可利用大规模静态数据预训练后迁移到动态场景。其设计目的不是单独增加一个模块，而是用上述表示、监督或渲染约束消除原流程在目标任务中的具体失配。提出首个运动感知前馈模型用于动态场景实时重建。核心创新是子弹时间(bullet-time)公式：为目标时间戳添加时间嵌入，模型聚合所有上下文帧信息重建该时间戳的完整3DGS场景。在150ms内完成重建，在静态和动态场景数据集上均达到SOTA(含与优化方法的比较)。适用时还应保留论文给出的边界：前馈模型的表达能力上限可能低于优化方法；对极端运动和遮挡的处理有待提升；需要大量混合数据训练。
% METHOD: 2412.06424
\paragraph{Deblur4DGS（2024.12）\cite{2412_06424}}
Deblur4DGS 处理模糊单目视频。普通 4DGS 默认每一帧是某个瞬时时刻的清晰观测，但真实视频中的运动模糊是曝光区间内相机位姿和动态物体状态连续变化后的时间平均。Deblur4DGS 的核心改动是把“曝光区间内连续 4D 表示估计”简化为“曝光时间参数估计”：给定一帧模糊图像，系统在曝光区间内采样多个潜在时间点，分别渲染清晰图像，再平均得到合成模糊图，与输入帧做重建损失。

具体模块分为静态相机模糊和动态物体模糊两部分。静态背景使用 static Gaussians，并通过小型 camera motion predictor 估计曝光起止位姿，中间位姿由 SE(3) 插值得到。动态区域由 SAM2 掩码定位；方法把视频切成若干片段，从每段中选择动态区域最清晰的帧作为 blur-aware variable canonical Gaussians，避免所有大运动都被迫从同一个 canonical 状态变形。曝光区间内的动态高斯由相邻整数时间状态插值得到，训练时同时优化 static Gaussians、canonical dynamic Gaussians、变形算子和每帧曝光时间参数。为了避免“随便拉长曝光也能拟合模糊”的平凡解，方法加入 exposure regularization、多帧一致性和多分辨率一致性。相对 Deformable 3DGS，它把相机抖动和物体运动模糊直接写进图像形成模型，因此同一个优化结果还能用于去模糊、插帧和稳定。
\textbf{实现细节：} 核心改动。Deblur4DGS（2024.12） 将上述"直接从模糊帧监督"的范式，改为"先建模曝光时间内的连续动态，再合成模糊帧进行监督"。直觉上，一张模糊照片是曝光时间内多张清晰照片的平均——如果能重建出这些隐含的清晰瞬态，再取平均与模糊输入比较，就能在拟合模糊外观的同时迫使模型学到清晰的底层表示。关键洞察来自 3DGS 的显式运动建模：不需要逐帧估计全部动态高斯，只需估计一组逐高斯的曝光时间参数 $\mathbf{w}_t$，再通过插值获得曝光内的连续状态。这与早期的 NeRF 方法 DyBluRF 形成对比——DyBluRF 需要为曝光内每个时刻建立独立的连续运动模型，优化复杂且难以产生高质量结果；Deblur4DGS 利用 3DGS 的显式运动表示，将连续动态估计简化为曝光时间估计。同期工作 BARD-GS 采用了类似策略，但由于优化约束不足（尤其对大运动场景），效果不及 Deblur4DGS。实现原理——静态区域。对于第 $t$ 帧模糊图像 $B_t$ 及其动态区域掩码 $M_t$（由 SAM2 提取），将曝光时间均匀划分为 $N$ 个时间戳。静态高斯 $S$ 在所有时间戳共享，只需估计连续相机位姿：优化曝光起止时刻的位姿 $P_{t,1}, P_{t,N}$，通过李代数线性插值获得中间位姿 $P_{t,i} = P_{t,1} \exp(\frac{i-1}{N-1} \log(P_{t,1}^{-1}P_{t,N}))$。为更稳定地优化位姿，Deblur4DGS 部署一个轻量 MLP 作为相机运动预测器，先用静态重建损失 $\mathcal{L}_{\mathrm{rec}}^s$ 预训练该 MLP 和静态高斯 $S$，其中 $\mathcal{L}_{\mathrm{rec}}^s = (1-\lambda)L_1(\hat{B}_t^s, B_t^s) + \lambda\mathcal{L}_{\mathrm{ssim}}(\hat{B}_t^s, B_t^s)$，$\hat{B}_t^s = (1-M_t) \odot \hat{B}_t$ 和 $B_t^s = (1-M_t) \odot B_t$ 分别为合成模糊图像和输入模糊图像的静态区域。对于每个时间戳，用对应的 $P_{t,i}$ 和 $S$ 渲染清晰图像，将 $N$ 幅图像取平均得到合成模糊图像 $\hat{I}_t$，与输入计算重建损失 $\mathcal{L}_{\mathrm{rec}} = (1-\lambda)L_1(\hat{I}_t, B_t) + \lambda_{\mathrm{ssim}}\mathcal{L}_{\mathrm{ssim}}(\hat{I}_t, B_t)$，其中 $\lambda=0.2$，与标准 3DGS 设置一致。实现原理——动态区域的模糊感知可变规范高斯。动态区域的核心挑战在于规范高斯的选择。
\textbf{设计动机与边界：} 相对原始 3DGS 以独立高斯和图像重建误差为主的设定，Deblur4DGS 把问题落在以下可实现的改动上：核心方法：1）曝光时间估计：利用3DGS的显式3D运动建模优势，将复杂的连续动态表示估计简化为曝光时间参数估计。一旦曝光时间确定，即可通过最近整数时间戳间的插值获得连续动态表示，避免DyBluRF等方法的复杂额外运动建模；2）模糊感知可变规范高斯：不同于使用单一规范高斯的方法，提出随时间变化的规范高斯。将视频均匀分为K段，根据动态区域模糊水平选择对应最清晰帧的规范高斯，降低大运动物体的建模难度；3）正则化策略：引入曝光正则化项、多帧一致性和多分辨率一致性正则化项避免平凡解；4）使用SAM2提取动态区域掩码，分别处理静态和动态区域。使用小型MLP作为相机运动预测器。为了消除原始 3DGS 在该目标任务中的具体失配，方法把上述表示、监督或渲染约束作为可学习和可执行的改动，而不是单独堆叠模块。提出Deblur4DGS框架从模糊单目视频重建4D模型。核心创新：1）将曝光时间内连续动态表示估计简化为曝光时间估计；2）模糊感知可变规范高斯适应大运动物体；3）多种正则化避免平凡解。可同时应用于去模糊、帧插值和视频稳定四项任务。适用时还应保留论文给出的边界：依赖SAM2的动态区域分割质量；对极端运动模糊可能仍有局限；曝光时间估计的精度受初始姿态质量影响。
\textbf{方法信息：} 代码：\url{https://github.com/zcsrenlongz/deblur4dgs}。



% METHOD: 2505.05356
\paragraph{Time of the Flight of the Gaussians（2025.05）\cite{2505_05356}}
把动态高斯用于连续波ToF相机。这里的输入不是普通RGB，而是间接深度测量；C-ToF传感器测到的是调制光返回后的相位/响应，优化中可能出现“深度几何错误但传感器响应仍被拟合”的坏解。方法因此把监督接口从RGB颜色改成符合C-ToF物理测量的响应，并围绕这个成像模型加入几何启发式约束，使高斯在匹配传感器读数的同时保持合理深度。它的价值不在于替换4DGS的全部动态表示，而在于说明当传感器不是相机RGB时，3DGS的渲染目标、损失和深度稳定策略都要随传感器物理重写。论文在快速运动场景中展示了相对神经体渲染方法约百倍速度优势。
\textbf{实现细节：} Time of the Flight of the Gaussians（2025.05） 致力于同时解决速度和稳定性两个问题。速度方面，该方法用 3D 高斯 Splatting 替代 NeRF 的 MLP 场景表示，将优化和渲染速度提升 100 倍。稳定性方面，论文深入分析了 C-ToF 辐射场优化中的经验偏差（empirical biases），揭示了两层核心问题：深度仅通过传感器测量值的重建被间接优化，优化对初始化和超参数高度敏感；而 NeRF 的 MLP 天然偏向低熵（简单）解——沿射线的单一高密度尖峰比多个中等密度值具有更低熵——GS 却无此内在偏差，高斯可以沿射线任意放置来拟合目标。C-ToF GS 图像形成模型。在标准 GS 渲染中，像素颜色是各高斯颜色的加权叠加。为适配 C-ToF 信号，该方法将场景辐射率替换为红外光传输模型：每个高斯的颜色 $c_k$ 被替换为光源强度 $s$ 与视角相关表面反射率 $r_k$（用 16 个球谐系数参数化）的乘积，再乘以逆平方衰减 $1/d_k^{}2$（$d_k$ 为高斯中心到相机光心距离），并加入相量调制 $W(d_k) = \exp(\mathrm{i}\psi_k)$，其中相位偏移 $\psi_k = 4\pi d_k f / \mathrm{c}$。由此得到相量形成模型：像素 $\mathbf{x}$ 处的相量为 $\mathbf{p}(\mathbf{x}) = \mathbf{p}^{\mathrm{bg}} T_N + \sum_{k=1}^{N}\frac{2sr_k}{d_k^{}2}W(d_k)\,o_k\,\mathcal{G}^{\mathrm{2D}}_k(\mathbf{x})\,T_k^{}2$，其中透射率被平方以表示光在辐射场中往返一次的衰减。进一步，为处理原始四分量帧的异步采集（四帧并非同时拍摄），用调制函数 $\phi(d_k) = \{\sin\psi_k, \cos\psi_k, -\sin\psi_k, -\cos\psi_k\}$ 替换相量，得到原始帧形成模型 $\mathbf{q}(\mathbf{x}) = \mathbf{q}^{\mathrm{bg}} T_N + \sum_{k=1}^{N}\frac{sr_k}{d_k^{}2}\phi(d_k)\,o_k\,\mathcal{G}^{\mathrm{2D}}_k(\mathbf{x})\,T_k^{}2$。关键矛盾在于：场景深度 $d(\mathbf{x}) = \sum_{k=1}^{N}d_k\,o_k\,\mathcal{G}^{\mathrm{2D}}_k(\mathbf{x})\,T_k$ 是各高斯深度的加权均值，但即使完美重建原始帧 $\mathbf{q}(\mathbf{x})$，也不保证 $d(\mathbf{x})$ 与传感器导出深度一致。两种启发式策略。与依赖额外学习型深度先验的方法不同，该方法提出的两种策略完全基于对优化过程的经验分析，不引入外部假设。
\textbf{方法信息：} Meta 参与了这篇论文。
\textbf{设计动机与边界：} 相对原始 3DGS 以独立高斯和图像重建误差为主的设定，Time of the Flight of the Gaussians 把问题落在以下可实现的改动上：该方法的核心挑战与解决方案为：(1) 问题分析：C-ToF相机间接测量深度，在使用快速高斯基元表示时优化极度不稳定。场景深度高度不准确但仍能产生高质量传感器测量重建，问题高度欠约束；(2) 两个启发式策略：改善高斯表示的场景几何精度。这些策略帮助在约束性C-ToF感知条件下稳定优化过程；(3) 应用场景：处理包括快速运动(如挥动棒球棒)的动态场景，从单个相机视角进行4D动态场景重建。其设计目的不是单独增加一个模块，而是用上述表示、监督或渲染约束消除原流程在目标任务中的具体失配。提出从单视角C-ToF传感器数据使用3DGS重建动态场景的方法，通过两个启发式策略改善深度间接优化问题，速度比神经体绘制方法快100倍。适用时还应保留论文给出的边界：深度间接测量导致优化仍有一定不稳定性；依赖特定启发式策略而非理论保证；对极端多径效应的处理有限。
% METHOD: 2506.18792
\paragraph{ViDAR（2025.06）\cite{2506_18792}}
把MoSca这类单目重建结果视为退化初稿，再用视频扩散模型补足缺失视角。它先得到静态和动态高斯，渲染出多视角但有伪影的图像；随后用DreamBooth风格个性化的SDXL作为生成增强器，把输入视频的外观和细节注入这些退化视图。为了避免扩散监督破坏几何，它使用Track Anything获得动态掩码，并设计扩散感知损失和相机姿态优化，把生成视图重新对齐到底层4D场景。相对Deformable 3DGS，ViDAR的改变是把单目不可见区域交给个性化视频生成先验，再蒸馏回高斯表示。
\textbf{实现细节：} ViDAR（2025.06） 的关键洞察是：正则化方法（如 MoSca）通过几何约束获得紧凑的几何表示，但渲染质量往往偏低；而生成式方法利用扩散模型实现高视觉质量，却难以保持时空一致性。ViDAR 将扩散模型作为外观先验注入重建流程，同时通过扩散感知策略控制其几何不一致性。据作者所知，ViDAR 是首个以几何一致方式将扩散先验融入单目视频重建的方法。具体流程分为三步。首先，使用 MoSca 作为基线获得初始 4D 表示，将场景分解为静态高斯 $\mathcal{G}_s$ 和动态高斯 $\mathcal{G}_d$（$\mathcal{G} = \mathcal{G}_d \cup \mathcal{G}_s$）。值得注意的是，MoSca 原始的对极误差分析在动/静态区域分类中会将背景误划为动态区域，导致浮点伪影（floater）。ViDAR 在此阶段引入 Track Anything 获取动态掩码 $D_t$ 替代 MoSca 的对极误差分类，更精确地约束动态高斯的位置，从而改善后续伪多视角样本的生成质量。然后用 DreamBooth 策略对预训练的 Stable Diffusion XL（SDXL）进行场景级个性化微调，将输入视频视为图像集合进行训练，使特定文本标记触发该场景的外观。
\textbf{方法信息：} Huawei 参与了这篇论文。



% METHOD: 2603.29734
\paragraph{GRVS（2026.03）\cite{2603_29734}}
是可泛化单目动态新视角合成路线。它从静态ConvGLR式框架出发，加入循环状态，让模型逐帧预测目标视频并传播隐藏表示；目标时间和输入时间不必一一对应，因此可以表达“场景冻结而相机移动”的子弹时间效果。几何侧，GRVS用动态平面扫描体把输入帧按目标相机位置做单应投影，在同一体中同时编码相机运动和场景运动；网络随后在低分辨率潜在空间用3D UNet融合这些线索。patch化和反patch化模块降低时序展开的计算量。相对优化式4DGS，GRVS不维护逐场景高斯参数，而是学习从视频到目标视角序列的前馈渲染器；它放在本章中，主要是作为4DGS路线的可泛化替代参照。
\textbf{方法信息：} Huawei 参与了这篇论文。
\textbf{设计动机与边界：} 相对原始 3DGS 以独立高斯和图像重建误差为主的设定，GRVS 把问题落在以下可实现的改动上：GRVS方法基于ConvGLR静态NVS框架适配：(1)循环回路设计：将可泛化模型包装为循环循环，逐帧预测目标视频帧同时传播循环隐藏状态。支持输入和目标视频之间的无界异步映射，不限于同步时间对应。特别支持子弹时间效果（场景冻结而视角自由变化）。引入patch化和反patch化模块提升计算效率并允许训练时的时序展开。(2)动态平面扫描：在动态输入上使用平面扫描体(PSV)解耦相机运动和场景运动。通过单应变换将输入帧投影到目标相机位置，PSV同时编码几何和运动信息，模型学习从中提取场景一致性表示。(3)架构：使用高效3D UNet在低分辨率潜在空间全局执行渲染操作，支持任意输入-目标时间映射。其设计目的不是单独增加一个模块，而是用上述表示、监督或渲染约束消除原流程在目标任务中的具体失配。提出将可泛化静态NVS框架适配到动态输入的新方法。两个关键创新：(1)循环回路实现输入和目标视频之间的无界异步映射；(2)在动态输入上高效使用平面扫描解耦相机和场景运动，实现精细的6自由度相机控制。在UCSD数据集和新构建的Kubric-4D-dyn数据集上优于四种3DGS方法和两种扩散方法。适用时还应保留论文给出的边界：循环展开的训练计算开销，对极端动态运动的处理有限，依赖平面扫描的深度分辨率。
% METHOD: 2605.04435
\paragraph{Ground4D（2026.05）\cite{2605_04435}}
Ground4D 面向非结构化越野场景，问题比普通城市动态视频更复杂：地面有岩石、沟壑和泥土等高频几何，自车在颠簸路面上产生强相机抖动，植被和尘土又带来非刚性短时变化。它不采用逐场景长时间优化，而是以 VGGT 式前馈几何骨干为基础：输入一段上下文帧后，冻结骨干先估计相机、深度和场景几何，随后若干预测头输出 pixel-aligned Gaussians、法线、动态置信度和天空/背景线索，并把像素反投影到 canonical Gaussian space。

核心模块是 voxel-grounded temporal Gaussian aggregation。系统把 canonical 空间划成稀疏体素网格，在每个非空体素内收集不同时间投来的高斯观测；面对一个查询时间，模型用查询条件的 temporal attention 给体素内候选高斯打分，并在体素内做 softmax，使一个空间位置优先选择最适合该时刻的基元，而不是把多帧观测简单平均。融合也按属性拆开：位置和颜色用加权均值，不透明度采用最大值与均值混合，尺度在 log-space 做几何平均，旋转用四元数均值。动态区域由 dynamic confidence 和 TAPIP3D 式轨迹插值补充，法线则约束高斯朝向和尺度，使其贴合崎岖地形。相对单目 4DGS，它把时间聚合的基本单位从“全局变形场”改成“局部体素中的多时刻高斯证据选择”。
\textbf{实现细节：} Ground4D（2026.05） 的关键洞察是将时间竞争限制在局部空间邻域内，消除时间选择性与空间占据之间的矛盾。该方法利用冻结的 VGGT 骨干网络处理 $T$ 帧上下文图像，产生三路互补的 token 流：聚合跨视角 token $F_t^{\mathrm{agg}}$、图像级 token $F_t^{\mathrm{img}}$ 和 DINOv2 语义 token $F_t^{\mathrm{dino}}$。骨干网络在微调时保持冻结，仅训练预测头和时序融合模块。相机头从 $F_t^{\mathrm{agg}}$ 解码相机位姿，DPT 深度头预测逐像素深度 $D_t$，通过反投影 $\mathcal{P}_t(u,v) = R_t^{}\top(K_t^{-1}D_t(u,v)[u,v,1]^{}\top - T_t)$ 获得 3D 点作为高斯基元中心。高斯头在 $F_t^{\mathrm{img}}$ 和 $F_t^{\mathrm{dino}}$ 上预测 11 通道的密集高斯属性图（颜色、不透明度、尺度和旋转），将像素反投影至共享规范高斯空间 $\mathcal{G}^s = \bigcup_t \mathcal{G}_t^s$。体素锚定的时序高斯聚合分三步。首先，将规范空间划分为体素边长 $\rho$ 的均匀网格，每个高斯基元中心 $\mu_i$ 通过 round 函数量化为整数坐标 $v_i = \mathrm{round}(\mu_i / \rho) \in \mathbb{Z}^{}3$，共面高斯归入同一体素，共产生 $M \ll N$ 个非空稀疏体素，将全部基元划分为可变大小组。体素化步骤不可微，仅用于分组，梯度通过聚合阶段的高斯属性回传。然后，对每个体素 $m$ 内的共面高斯执行查询条件的时序注意力：每个基元携带特征 $f_i^{\mathrm{gs}}$（高斯属性先编码为特征嵌入，再与高斯头最后一层的残差特征拼接），经特征编码器 $\phi_f$ 投影至共享隐空间；时间戳 $\tau_i$ 经多频率正弦编码 $\gamma(\cdot)$ 和时间 MLP $\phi_t$ 在多尺度上捕捉时间动态，二者加性融合为 $h_i = \phi_f(f_i^{\mathrm{gs}}) + \phi_t(\gamma(\tau_i))$；查询时间 $\tau^{*}$ 类似编码为 $h^{*} = \phi_t(\gamma(\tau^{*}))$，轻量注意力 MLP 输出标量对数 $a_i = \mathrm{MLP}_{\mathrm{attn}}([h_i; h^{*}])$。最后，在每个体素内独立进行温度缩放 softmax 归一化 $w_i = \exp(a_i/\beta) / \sum_{j \in \mathcal{V}_m}\exp(a_j/\beta)$，确保每个非空体素无论查询时间距上下文帧多远都能产生有效基元。
\textbf{技术实现补充：} Ground4D（2026.05） 的关键洞察是将时间竞争限制在局部空间邻域内，消除时间选择性与空间占据之间的矛盾。该方法利用冻结的 VGGT 骨干网络处理 $T$ 帧上下文图像，产生三路互补的 token 流：聚合跨视角 token $F_t^{\mathrm{agg}}$、图像级 token $F_t^{\mathrm{img}}$ 和 DINOv2 语义 token $F_t^{\mathrm{dino}}$。骨干网络在微调时保持冻结，仅训练预测头和时序融合模块。相机头从 $F_t^{\mathrm{agg}}$ 解码相机位姿，DPT 深度头预测逐像素深度 $D_t$，通过反投影 $\mathcal{P}_t(u,v) = R_t^\top(K_t^{-1}D_t(u,v)[u,v,1]^\top - T_t)$ 获得 3D 点作为高斯基元中心。高斯头在 $F_t^{\mathrm{img}}$ 和 $F_t^{\mathrm{dino}}$ 上预测 11 通道的密集高斯属性图（颜色、不透明度、尺度和旋转），将像素反投影至共享规范高斯空间 $\mathcal{G}^s = \bigcup_t \mathcal{G}_t^s$。体素锚定的时序高斯聚合分三步。首先，将规范空间划分为体素边长 $\rho$ 的均匀网格，每个高斯基元中心 $\mu_i$ 通过 round 函数量化为整数坐标 $v_i = \mathrm{round}(\mu_i / \rho) \in \mathbb{Z}^3$，共面高斯归入同一体素，共产生 $M \ll N$ 个非空稀疏体素，将全部基元划分为可变大小组。体素化步骤不可微，仅用于分组，梯度通过聚合阶段的高斯属性回传。然后，对每个体素 $m$ 内的共面高斯执行查询条件的时序注意力：每个基元携带特征 $f_i^{\mathrm{gs}}$（高斯属性先编码为特征嵌入，再与高斯头最后一层的残差特征拼接），经特征编码器 $\phi_f$ 投影至共享隐空间；时间戳 $\tau_i$ 经多频率正弦编码 $\gamma(\cdot)$ 和时间 MLP $\phi_t$ 在多尺度上捕捉时间动态，二者加性融合为 $h_i = \phi_f(f_i^{\mathrm{gs}}) + \phi_t(\gamma(\tau_i))$；查询时间 $\tau^*$ 类似编码为 $h^* = \phi_t(\gamma(\tau^*))$，轻量注意力 MLP 输出标量对数 $a_i = \mathrm{MLP}_{\mathrm{attn}}([h_i; h^*])$。最后，在每个体素内独立进行温度缩放 softmax 归一化 $w_i = \exp(a_i/\beta) / \sum_{j \in \mathcal{V}_m}\exp(a_j/\beta)$，确保每个非空体素无论查询时间距上下文帧多远都能产生有效基元。
\textbf{方法信息：} 代码：\url{https://github.com/wsnbws/ground4d}。
\textbf{设计动机与边界：} 相对原始 3DGS 以独立高斯和图像重建误差为主的设定，Ground4D 把问题落在以下可实现的改动上：Ground4D针对越野场景的前馈4D重建，解决高频几何细节（岩石、沟壑）、自运动抖动（颠簸路面）和非刚性动态（植被摆动、尘土飞扬）导致的时序高斯观测冲突问题。核心创新是体素接地时间高斯聚合（Voxel-Grounded Temporal Gaussian Aggregation）：将规范高斯空间划分为规则的三维空间体素网格，在每个体素内执行查询条件的时间注意力评分，计算每个时间戳对应的高斯激活概率。关键的体素内softmax归一化确保时间选择性（选择最佳时间戳的高斯表示）和空间占据率（覆盖完整场景的所有区域）从相互冲突的目标变为相互增强——体素内的竞争机制使高斯同时满足时间一致性和空间完整性约束。引入表面法线线索作为辅助几何指导，正则化高斯基元的几何形状使其紧密贴合越野场景复杂的地形起伏。三阶段前馈流水线：(1)冻结预训练VGGT骨干处理T帧输入构建规范高斯空间；(2)体素分区时间聚合解决时序冲突；(3)精化高斯与插值动态高斯和参数化天空背景合成最终场景表示。其设计目的不是单独增加一个模块，而是用上述表示、监督或渲染约束消除原流程在目标任务中的具体失配。提出体素接地时间高斯聚合解决越野场景前馈4D重建中的时序高斯冲突问题，空间局部化条件使时间选择性和空间占据率从冲突目标变为相互增强，实现零样本跨域泛化至未见越野场景，为非结构化环境的4D重建提供了有效解决方案。适用时还应保留论文给出的边界：体素大小超参数可能需要针对不同场景尺度调优。

\section{视频重建与时序一致性}

这里的方法不一定都重新定义4D基元，而是解决动态视频系统化使用时的两个问题：长序列如何流式更新和压缩，视频编辑或稳定任务中如何利用3DGS提供跨帧一致约束。它们通常继承4D-GS、Scaffold-GS或前景背景分解，再把压缩、增量编码、视频扩散或局部重建接入流水线。
% METHOD: 2406.02541
\paragraph{Video-3DGS（以前景--背景高斯重建约束视频编辑）（2024.06）}
Video-3DGS~\cite{2406_02541} 的目标是让零样本扩散视频编辑保持跨帧一致，而不是单独提出动态 NVS 基元。第一阶段以 Masked-and-Clipped COLMAP 处理视频片段，分别初始化动态前景与复杂背景的点云；两套 Frg-3DGS/Bkg-3DGS 各自优化，随后由可学习二维参数图合成完整画面。

第二阶段将该重建器作为编辑视频的时间约束：编辑器输出的帧必须能由一致的前景/背景 3DGS 解释，重建残差反向抑制逐帧扩散造成的漂移。前景、背景分离使动态物体不必由单一静态高斯场硬拟合。该方法依赖片段的相机与掩码估计，长时大拓扑变化或错误前景分割会削弱重建约束。

% METHOD: 2411.07541
\paragraph{HiCoM（2024.11）\cite{2411_07541}}
HiCoM 面向多视角流式动态场景，问题不是离线训练一个完整 4D 模型，而是视频帧持续到达时如何快速更新 3DGS。它继承 3DGStream 的在线思想，但指出初始 3DGS 在有限相机视角下容易过拟合并膨胀，高斯数量越多，后续每帧运动学习越慢。第一步因此是 perturbation smoothing：训练初始帧时给高斯位置加入小噪声，让模型不要把有限视角中的偶然纹理记成大量细碎基元，收敛后得到更紧凑、也更适合后续传播的初始表示。

帧间更新使用 Hierarchical Coherent Motion。方法把空间划分为多层区域，只给包含高斯的区域维护运动参数；相邻帧之间的变换由 Neural Transformation Cache 表示，但不是为每个高斯独立预测，而是在区域层级上共享和传播。由于 3DGS 分布天然不均匀，大量空区域无需任何运动参数，局部区域内的高斯又通常具有相近运动，这种层次缓存比逐高斯变形更省。对新出现或前一帧未覆盖的内容，HiCoM 会先增加额外高斯精修当前帧，再把有效新增内容合并回主集合，同时移除等量低贡献高斯以保持规模稳定。并行训练多个帧则进一步降低在线 wall time。相对 MoDec-GS 的离线紧凑建模，HiCoM 的锚点价值在“边接收边更新”的流式视频工程。

具体流程上，输入是新到达的多视角图像和上一帧高斯集合；系统先用层次缓存给当前帧初始化区域运动，再通过渲染损失优化这些区域变换，最后输出更新后的高斯集合和可继续复用的 motion cache。这个流程把最昂贵的逐点运动搜索压到局部区域和缓存更新中，因此适合连续视频流，而不是每帧从零优化。
\textbf{方法信息：} 代码：\url{https://github.com/gqk/hicom}。



% METHOD: 2412.04469
\paragraph{QUEEN（2024.12）\cite{2412_04469}}
面向自由视点视频的在线增量编码。它不强行给帧间变化套一个变形场，而是直接学习相邻帧高斯属性残差：位置、SH系数、不透明度、尺度和旋转都可以以残差形式更新。非位置属性残差通过潜在编码器/解码器量化，位置残差则由门控模块稀疏化，只把码率花在真正需要移动的基元上。静态动态分离由高斯视图空间梯度差分信号引导，这个信号会指出哪些高斯在视图中产生明显时序变化，从而决定哪些基元需要更新。相对4D-GS的离线优化，QUEEN关注每帧如何以少量增量更新高斯表示，同时满足约0.7MB/帧、数秒级训练和约350 FPS渲染这类流媒体约束。
\textbf{实现细节：} QUEEN（2024.12） 以 3DGStream（Sun et al.）为基线，解决在线流式自由视角视频的三个矛盾需求：低延迟编码、低内存带宽、高重建质量。核心改动是：移除所有结构约束，对所有高斯属性自由学习帧间残差，并设计量化-稀疏框架将残差压缩到极低存储，使模型大小随场景动态程度自适应变化。Baseline 与改进。3DGStream 使用 Instant-NGP 哈希编码结构缓存位置和旋转残差，但存在两个关键限制。第一，仅更新部分属性（位置和旋转），其余属性（不透明度、球谐系数、尺度等）保持不变，模型表达力受限。第二，施加固定的结构约束——每帧使用相同大小的 Instant-NGP 编码，无法根据场景动态程度调整存储开销。在静态场景中浪费带宽，在剧烈运动场景中又容量不足。QUEEN 的核心洞察是：不施加任何结构约束，对所有高斯属性自由学习帧间残差，同时通过量化-稀疏框架按需压缩残差，使存储开销与实际场景动态成正比。这样，变化小的帧自动获得更小的模型尺寸，变化大的帧获得更多比特。属性残差量化。直觉上，如果场景在相邻两帧间变化很小（如人物静止站立），那么绝大多数高斯的属性几乎不变，残差应集中在零附近。量化正好利用这一分布特性：将连续残差映射到少量离散整数值，大幅减少存储。具体而言，对每个时间步 $t$，QUEEN 学习全部属性残差：$\mathcal{A}_t = \mathcal{A}_{t-1} + \mathcal{R}_t$。对旋转、尺度、不透明度、球谐系数等属性残差，维护每个高斯的量化整数潜码 $\mathbf{l}_i \in \mathbb{Z}^L$，通过共享线性解码器 $D$ 恢复残差： 其中 $\hat{\mathbf{l}}_i$ 为连续潜码，STE（Straight-Through Estimator，直通估计器）在前向传递中将连续值取整为整数（$\mathrm{STE}(\cdot) = \mathrm{round}(\cdot)$），在反向传递中直接透传梯度（$\frac{\partial \mathrm{STE}}{\partial \hat{\mathbf{l}}_i} \approx 1$），使不可微的取整操作变得可微。
\textbf{锚点细化：} QUEEN：以 3DGStream（Sun et al.）为基线，解决在线流式自由视角视频的三个矛盾需求：低延迟编码、低内存带宽、高重建质量。核心改动是：移除所有结构约束，对所有高斯属性自由学习帧间残差，并设计量化-稀疏框架将残差压缩到极低存储，使模型大小随场景动态程度自适应变化。
\textbf{设计动机与边界：} 相对原始 3DGS 以独立高斯和图像重建误差为主的设定，QUEEN 把问题落在以下可实现的改动上：核心设计：(1) 直接残差学习：不施加结构约束(如变形场)，直接学习帧间高斯属性(位置、SH系数、不透明度、尺度、旋转)的残差，保持高重建质量和泛化性。(2) 量化-稀疏框架：对非位置属性残差使用学习潜在编码器-解码器进行量化；对位置残差使用学习门控模块进行稀疏化。(3) 静态/动态分离：利用高斯视图空间梯度差分向量作为信号区分场景的静态和动态内容，指导有效稀疏学习并加速训练。(4) 在线增量更新：每帧增量更新3DGS表征，满足实时约束。其设计目的不是单独增加一个模块，而是用上述表示、监督或渲染约束消除原流程在目标任务中的具体失配。提出首个高效的3DGS在线FVV流媒体框架。直接学习帧间高斯属性残差(不施加结构约束)，通过学习潜在解码器量化非位置残差、学习门控模块稀疏化位置残差。利用高斯视图空间梯度差分向量分离静态/动态内容指导稀疏学习。高度动态场景仅0.7MB/帧，训练<5秒，渲染350 FPS。适用时还应保留论文给出的边界：对场景突变(如镜头切换)处理能力有限；长时间序列可能累积误差；量化精度与压缩率的权衡需调优。
\textbf{方法信息：} NVIDIA 参与了这篇论文；官方代码：\url{https://github.com/NVlabs/queen}；GitHub 关注度较高，约 185 stars。



% METHOD: 2503.05162
\paragraph{EvolvingGS（2025.03）\cite{2503_05162}}
EvolvingGS 面向长序列体积视频，尤其是复杂人体表演、衣物和道具交互。很多动态 3DGS 系统会把长视频切成多个 GoP（Group of Pictures）分别重建，这能降低单段优化难度，却会在 GoP 边界产生闪烁，也让相邻片段重复存储相似高斯。EvolvingGS 的改变是把“每段重新建一套高斯”改成“上一帧高斯继续演化到下一帧”：绝大部分基元被继承，只在拓扑变化、遮挡解除或快速变化区域增删少量点。

方法分成变形和精修两阶段。变形阶段先把当前帧的高斯模型作为源表示，用稀疏控制节点描述从 $t$ 到 $t+1$ 的粗运动；每个控制节点带位置、四元数旋转和平移，邻近节点通过线性混合蒙皮传播到高斯。优化时同时使用图像重建损失、ARAP 局部刚性约束、四元数归一化约束和光流引导，让继承高斯先大致对齐下一时刻。精修阶段再固定继承点的外观相关属性，包括颜色、尺度和不透明度，主要优化位置等 pose 属性；对于新出现或快速变化区域，方法允许选择性新增高斯，对已经不可见或贡献很低的点则剪枝。这里剪枝不只看不透明度，而看点对实际渲染的贡献，避免透明度固定后冗余点持续膨胀。

这种表示天然适合压缩和流式传输，因为相邻时间只需要存储从上一帧到下一帧的增量，而不是完整重发一个 3DGS 场景。论文利用这种时间相干性设计压缩流程，报告超过 50 倍压缩率，并在 ActorsHQ、自采 Dancer/Lecturer 等长序列多视角数据上展示了比 GoP 独立重建更稳定的时序质量。相对 4D-GS，它的重点不在通用时间查询函数，而在体积视频工程里如何长期继承、修补、剪枝和压缩高斯。
\textbf{实现细节：} EvolvingGS（2025.03） 的核心洞察是：传统 GoP（Group of Pictures）分段处理导致段间闪烁，但连续帧之间共享大量几何和外观信息，大多数高斯只需微调属性值即可复用。基于此，EvolvingGS 提出一种无需 GoP 分段的演化式两阶段框架。变形阶段：稀疏控制点粗对齐。给定已精调好的 $t$ 帧高斯模型 $\mathcal{S}_t$，EvolvingGS 目标是将其变形为 $\hat{\mathcal{S}}_{t+1}$ 以粗略匹配 $t+1$ 帧。Baseline 逐点优化变形的方式在大规模高斯场景下计算代价高。EvolvingGS 引入 $M$ 个均匀采样的稀疏控制点，每个控制点 $C_j$ 定义一个局部仿射变换（位置 $\mathbf{c}_j$、四元数 $\mathbf{q}_j$、平移 $\mathbf{t}_j$）。高斯 $G_i$ 的变形通过最近 $K^g$ 个控制点的加权混合实现： $$\hat{\boldsymbol{\mu}}\_i=\sum\_{j\in\mathcal{N}(G\_i)}w\_j(G\_i)\bigl[\mathbf{q}\_j(\boldsymbol{\mu}\_i-\mathbf{c}\_j)+\mathbf{c}\_j+\mathbf{t}\_j\bigr],$$ 混合权重 $w_j(G_i)=(1-\|\boldsymbol{\mu}_i-\mathbf{c}_j\|/d_{\max})^{}2$ 并归一化为和为 1。变形阶段保持高斯数量不变，且所有外观特征（SH 系数、缩放、不透明度）固定不动，只优化位姿相关特征。为加速收敛，EvolvingGS 引入光流引导损失： $$E_{\text{flow}}=\frac{1}{|\mathbb{1}|}\sum_{i=1}^{N_1}T_{i,q}\|\mathbf{x}_i^{2D}-(\hat{\mathbf{x}}_i^{2D}+F_q(\mathbf{x}_i^{2D}))\|,$$ 其中 $F_q$ 是预测的光流图，$\mathbb{1}$ 是过滤静止背景的指示函数（光流值低于 $\tau_f$ 时置零），$T_{i,q}$ 是高斯透射率。总目标 $E=E_{\mathrm{data}}+\lambda_1 E_{\mathrm{flow}}+\lambda_2 E_{\mathrm{arap}}+\lambda_3 E_{\mathrm{rot}}$ 还包含 ARAP 刚性约束和四元数归一化约束，确保变形平滑且物理合理。精化阶段：双流优化与贡献裁剪。粗对齐无法处理新物体出现和拓扑变化。精化阶段维护两套高斯：继承自变形阶段的参考点 $G^{\mathrm{ref}}$（外观特征固定，仅微调位姿，保证时序稳定）和新增的扩展点 $G^{\mathrm{ext}}$（无约束优化，用于重建新物体）。由于不透明度被固定，传统基于不透明度的裁剪策略失效。EvolvingGS 改用"贡献度"作为裁剪指标： $$\text{contrib}_{i,q}=\sum_{\text{pix}}\alpha_{\text{pix},i}^{2D}\prod_{j=1}^{i-1}(1-\alpha_{\text{pix},j}^{2D}),$$ 即高斯在某个视角下对所有像素的透射贡献之和。若最大贡献低于阈值 $\tau_c$ 或所有视角贡献总和低于阈值 $\tau_s$，则裁剪该高斯。
\textbf{方法信息：} ByteDance 参与了这篇论文；目前没有可靠完整公开代码记录。
\textbf{设计动机与边界：} 相对原始 3DGS 以独立高斯和图像重建误差为主的设定，EvolvingGS 把问题落在以下可实现的改动上：两阶段方法：(1) 变形阶段（粗对齐）：引入稀疏控制节点表示变形场，每个控制节点包含位置、四元数旋转和平移向量。通过线性混合蒙皮（LBS）加权融合邻近控制节点变形。结合ARAP约束维持局部刚性，利用光流引导加速收敛。(2) 精修阶段：固定继承点的外观特征（颜色、缩放、不透明度），优化位置属性。引入基于贡献度的剪枝机制而非基于不透明度，处理拓扑变化和新出现区域。选择性在快速变化区域增减高斯点。(3) 压缩：利用时间相干性设计简单有效的压缩算法，实现50倍以上压缩。其设计目的不是单独增加一个模块，而是用上述表示、监督或渲染约束消除原流程在目标任务中的具体失配。提出EvolvingGS两阶段策略（变形+精修），避免GoP分割导致的时间不稳定。通过增量演化表示实现任意GoP长度的高质量动态序列重建。利用时间相干性提出压缩算法实现50倍以上压缩率。采用纯显式表示保持快速渲染。适用时还应保留论文给出的边界：依赖多视角同步视频输入；初始帧的3DGS质量影响后续帧；极端拓扑变化处理仍有挑战。
% METHOD: 2503.13948
\paragraph{Light4GS（2025.03）\cite{2503_13948}}
Light4GS 是 4DGS 压缩路线中的代表工作，目标不是改变动态表示的语义，而是让 4DGS 能进入存储和传输约束更强的场景。原始 4DGS 的主要体积来自两部分：可变形高斯基元本身，以及多尺度 HexPlane latent embeddings。Light4GS 先修正 HexPlane 查询：对无界场景中的高斯中心做非线性空间收缩，再用 PCA 把坐标轴对齐到场景主方向，使六个平面上的采样更集中，减少不必要的空域覆盖。

压缩分两步。第一步是 spatio-temporal significance pruning（STP）：它不只按 opacity 剪枝，而是在多个时间和视角上统计每个可变形高斯与射线/像素的相交贡献、透射率和体积影响，得到全局显著性分数；densification 结束后分两次剪掉低显著性基元，保留基元再继续微调。由于 SH 系数占高斯数据体积很大，剩余基元的 SH 还做 entropy-constrained compression。第二步是 multiscale HexPlane context model（MHCM）：低尺度到高尺度依次编码，结合 intra-scale checkerboard context、inter-scale coarse-to-fine prediction 和 hyperprior 估计概率分布，再由 CompressAI 式 entropy model 压缩 latent embeddings；代码 README 也说明 HexPlane 尺寸需要设成 4 的倍数以兼容 masked convolution。相对普通 4DGS，它的改变在于把“动态可变形表示”变成可剪枝、可概率编码、可传输的压缩资产。
\textbf{实现细节：} Light4GS（2025.03） 针对可变形 4DGS 中高维潜嵌入和大量基元导致的存储膨胀问题，提出基于显著性剪枝与深度上下文模型的轻量压缩框架。Baseline 与改进。vanilla 4DGS 使用六面体平面（hexplane）——一种将 4D 时空特征分解为六个 2D 平面（xy、xz、yz、xt、yt、zt）的紧凑编码结构——来编码时空变化，但存在两个缺陷：六面体网格通常在 SfM 点云包围盒内初始化，与无界 3D 空间不匹配；标准坐标轴不一定对齐场景主方向，非主成分投影阻碍潜特征学习。Light4GS 引入 PCA 引导的六面体查询：先用非线性空间收缩将高斯中心映射到有界空间，再对空间坐标做 PCA 旋转对齐主方向，然后用变换后坐标进行多尺度六面体双线性插值。时空显著性剪枝（STP）。vanilla 4DGS 仅通过不透明度剪枝低贡献基元，忽略了基元在不同帧和视角下的全局渲染影响。直觉上，STP 要回答的问题是：如果移除某个高斯，所有帧和所有视角的渲染质量会受到多大影响？影响越小，该高斯越冗余。
\textbf{方法信息：} 代码：\url{https://github.com/evan-sudo/light4gs}。
\textbf{设计动机与边界：} 相对原始 3DGS 以独立高斯和图像重建误差为主的设定，Light4GS 把问题落在以下可实现的改动上：Light4GS框架包含两个核心组件：(1) PCA引导的Hexplane查询：对无界3DGS场景，采用非线性空间收缩映射高斯中心到有界域，再通过PCA对齐空间坐标到主方向，改善Hexplane的特征学习能力。(2) STP时空显著性修剪：不仅考虑基元的不透明度，还评估每个基元对渲染的全局影响范围（覆盖像素数）和时变影响，精确捕捉每个基元对渲染质量的贡献，消除超过64%的冗余基元。剩余基元通过熵约束球谐压缩进一步减小存储。(3) MHCM多尺度Hexplane压缩模型：设计从粗到细的上下文结构，集成帧内预测（利用同一Hexplane内的空间相关性）和帧间预测（利用不同尺度Hexplane间的相关性），结合超先验进行概率建模，实现高效压缩。其设计目的不是单独增加一个模块，而是用上述表示、监督或渲染约束消除原流程在目标任务中的具体失配。提出Light4GS轻量级4DGS压缩框架，基于4DGS实现超过120倍压缩并提升渲染FPS达20%。两大核心组件：(1) 时空显著性修剪(STP)策略消除超过64%的可变形基元，配合熵约束球谐压缩处理剩余基元；(2) 多尺度Hexplane压缩模型(MHCM)集成帧内/帧间预测和超先验，采用从粗到细的上下文结构实现高效的多尺度潜在嵌入压缩。适用时还应保留论文给出的边界：修剪策略在极端动态场景中可能过度修剪有效基元；深度上下文模型增加编码/解码复杂度；PCA对齐对非主成分方向的信息可能损失。
% METHOD: 2505.07539
\paragraph{GIFStream（2025.05）\cite{2505_07539}}
把Scaffold-GS锚点结构扩展为沉浸式视频流。每个锚点同时存储时间无关特征$f^C$和时间依赖特征$f_t^P$，二者共同解码目标时间的若干3D高斯；可学习掩码$M_{de}$会把静态区域的时间依赖特征压到接近零，只让真正运动的锚点携带动态信息。论文报告在复杂场景中约30\%锚点需要保留时间依赖特征，简单场景仅需0.3\%，这说明压缩并不是均匀削减所有高斯，而是把码率集中给运动区域。压缩时，时间依赖特征被投影成2D视频并用视频编解码器处理，训练中加入量化感知和熵估计；时间无关特征则通过排序和PCA等步骤进一步压缩。相对MoDec-GS这类紧凑动态表示，GIFStream更明确面向网络传输，把动态特征组织成可编码的视频流。
\textbf{实现细节：} GIFStream（2025.05） 解决变形场方法在高度动态场景中表达能力不足的问题。现有基于变形的 4DGS 方法使用位置编码或低维特征平面驱动变形，难以捕获快速运动。GIFStream 基于 Scaffold-GS 的锚点结构，为每个锚点同时存储时间无关特征 $f^{*} \in \mathbb{R}^C$ 和时间相关特征流 $f_t \in \mathbb{R}^P$（$t$ 为时间戳），在每个时刻通过两个预测头分别解码高斯属性和运动：属性头 $h_{attr}: [\tilde{f}; \hat{f}_t] \to \{\alpha^i_t, s^i_t, r^i_t, c^i_t\}_{i=1}^K$ 生成 $K$ 个高斯的 opacity、缩放、旋转和颜色；运动头则先用 KNN 聚合邻近锚点特征并引入可学习混合因子 $M_{knn}$ 实现粗-细运动组合，$\tilde{f}_t = (1{-}M_{knn})\bar{f}_t + M_{knn} f_t$，最终映射为运动向量 $R_t, T_t$。压缩方面，GIFStream 将时间无关和时间相关特征分别排列为两段视频序列，支持 HEVC 等视频编码器和基于自回归熵估计的端到端压缩，配合量化感知训练（quantization-aware training）。在 RTX 4090 上以 30 Mbps 码率提供 1080p 高质量沉浸视频，实现实时渲染和快速解码。
\textbf{设计动机与边界：} 相对原始 3DGS 以独立高斯和图像重建误差为主的设定，GIFStream 把问题落在以下可实现的改动上：GIFStream的核心思路是在Scaffold-GS锚点结构上同时存储时间无关特征和时间依赖特征流。每个锚点包含时间无关特征f\^{}C和时间依赖特征f\_t\^{}P，两者结合预测不同时间戳的3D高斯原语。时间依赖特征通过可学习掩码M\_de进行运动自适应缩放，静态区域的时间依赖特征被自动压缩至零。实验表明在复杂场景中约30%锚点需要保留时间依赖特征，简单场景仅需0.3%。每个锚点解码K个高斯原语，分别通过属性预测头和运动预测头生成。压缩方面，将时间依赖特征投影为2D视频利用视频编解码器压缩，采用量化感知训练结合渲染损失和熵估计实现端到端压缩。同时采用3D到2D排序策略和PCA主成分分析优化时间无关特征压缩。其设计目的不是单独增加一个模块，而是用上述表示、监督或渲染约束消除原流程在目标任务中的具体失配。提出了一种基于规范空间和变形场的4D高斯表示，引入时间依赖特征流来增强复杂运动建模能力，同时通过运动感知剪枝和时间/空间压缩网络实现端到端压缩，在30Mbps码率下实现高质量沉浸式视频渲染。适用时还应保留论文给出的边界：需要GPU进行实时解码渲染；对极端快速运动场景的建模能力仍有提升空间；压缩率与质量之间的权衡仍需进一步优化。
\textbf{方法信息：} Intel 参与了这篇论文；官方代码：\url{https://github.com/XDimLab/GIFStream}。



% METHOD: 2506.23957
\paragraph{GaVS（2025.06）\cite{2506_23957}}
把视频稳定改写为3D局部重建和目标姿态渲染。它用前馈重建模型为每帧预测像素对齐的Gaussian Splatting基元，再根据平滑后的目标相机轨迹重渲染稳定帧。由于前馈模型在真实视频中会有跨帧漂移，GaVS在测试时再优化三类约束：相邻帧提供多视图光度监督，并显式处理动态物体；光流给出密集对应，用来匹配跨帧GS原语并在扩展时间窗口内施加一致性；视频补全负责外推源帧边界，避免稳定后出现裁剪黑边。相对于2D或2.5D稳定，GaVS的改变是先建立局部3DGS，再用新的相机路径渲染，而不是直接扭曲图像。
\textbf{实现细节：} 传统视频稳定方法在 2D 域建模相机运动（仿射变换、单应矩阵或光流），逐帧独立变换后不可避免地裁剪画面、丢失内容，且缺乏显式 3D 建模容易产生扭曲伪影。GaVS（2025.06） 的核心思路是：既然源帧（不稳定）和目标帧（稳定）的相机位姿只差一个微小变换，完全没必要全局重建整个动态场景。因此该方法将视频稳定重新表述为时间一致的"局部重建与渲染"问题——用前馈网络为每一帧预测一组逐像素 GS 基元，再在平滑后的目标位姿上渲染。具体地，给定源帧 $I_k^s$ 和对应局部重建 $\mathcal{G}_k = f(I_k^s)$（$f$ 为增强深度估计器的 GS 预测头），目标帧通过在平滑位姿 $P_k^t$ 上渲染得到：$\hat{I}_k^t = R(\mathcal{G}_k, P_k^t)$。这种逐帧前馈设计自然统一了静态和动态场景区域的处理，无需像全局 4D 优化那样显式分离动静部分。然而，逐帧预测无法保证相邻重建的时间一致性，且预训练模型在开放域视频上存在泛化问题。为此，GaVS 引入测试时优化方案，从三个维度约束每个不稳定视频的局部重建质量。
\textbf{实现补充：} 实现上还要核对输入、输出、参数、特征、初始化、训练损失、正则、渲染或推理接口，以及优化中的剪枝、分裂、采样、编码或解码步骤；这些抓手决定该方法到底改了 3DGS 管线的哪一环。
\textbf{方法信息：} Huawei 参与了这篇论文。
\textbf{设计动机与边界：} 相对原始 3DGS 以独立高斯和图像重建误差为主的设定，GaVS 把问题落在以下可实现的改动上：GaVS利用前馈重建模型为每帧预测逐像素的Gaussian Splatting原语，在目标姿态下渲染稳定帧。为解决局部重建的时序一致性和预训练模型的泛化问题，引入测试时优化：(1)多视图光度监督——从相邻帧提供监督并显式考虑动态物体；(2)跨帧正则化——通过光流的密集对应关系匹配跨帧GS原语，在膨胀时间窗口内施加一致性约束；(3)场景外推——使用视频补全方法外推源帧边界内容，确保全帧稳定。利用高斯滤波获取稳定相机轨迹。其设计目的不是单独增加一个模块，而是用上述表示、监督或渲染约束消除原流程在目标任务中的具体失配。提出GaVS，首个3D驱动的视频稳定方法，通过测试时优化的前馈重建模型预测Gaussian Splatting原语，结合多视图动态感知光度监督和跨帧正则化实现时序一致的局部重建，利用场景外推模块避免帧裁剪。适用时还应保留论文给出的边界：测试时优化增加计算开销；依赖光流估计的质量；对极端运动模糊的处理有限。
% METHOD: 2508.07557
\paragraph{Splat4D（2025.08）\cite{2508_07557}}
把4D内容创作与3DGS结合。它先用多视角扩散模型把单目视频扩展成多视角图像序列，再借助LGM、Splatter Image等预训练前馈3DGS模型获取空间和深度特征，把这些特征整理成splat图像并反投影为粗4D高斯。粗结果通常会有跨视角或跨时间不一致，因此Splat4D会再次渲染多视角结果，定位不可靠区域，用视频扩散模型和非对称U-Net对这些区域做时空精炼。最后，它微调通用高斯场预测器，使文本、图像或视频条件都能生成4D内容。相对L4GM，Splat4D更依赖扩散模型进行跨视角和跨时间修复，把生成一致性作为重建流水线的一部分，而不是只在后处理里修图。
\textbf{实现细节：} QUEEN（2024.12） 以 3DGStream（Sun et al.）为基线，解决在线流式自由视角视频的三个矛盾需求：低延迟编码、低内存带宽、高重建质量。核心改动是：移除所有结构约束，对所有高斯属性自由学习帧间残差，并设计量化-稀疏框架将残差压缩到极低存储，使模型大小随场景动态程度自适应变化。Baseline 与改进。3DGStream 使用 Instant-NGP 哈希编码结构缓存位置和旋转残差，但存在两个关键限制。第一，仅更新部分属性（位置和旋转），其余属性（不透明度、球谐系数、尺度等）保持不变，模型表达力受限。第二，施加固定的结构约束——每帧使用相同大小的 Instant-NGP 编码，无法根据场景动态程度调整存储开销。在静态场景中浪费带宽，在剧烈运动场景中又容量不足。QUEEN 的核心洞察是：不施加任何结构约束，对所有高斯属性自由学习帧间残差，同时通过量化-稀疏框架按需压缩残差，使存储开销与实际场景动态成正比。这样，变化小的帧自动获得更小的模型尺寸，变化大的帧获得更多比特。属性残差量化。直觉上，如果场景在相邻两帧间变化很小（如人物静止站立），那么绝大多数高斯的属性几乎不变，残差应集中在零附近。量化正好利用这一分布特性：将连续残差映射到少量离散整数值，大幅减少存储。具体而言，对每个时间步 $t$，QUEEN 学习全部属性残差：$\mathcal{A}_t = \mathcal{A}_{t-1} + \mathcal{R}_t$。对旋转、尺度、不透明度、球谐系数等属性残差，维护每个高斯的量化整数潜码 $\mathbf{l}_i \in \mathbb{Z}^L$，通过共享线性解码器 $D$ 恢复残差： 其中 $\hat{\mathbf{l}}_i$ 为连续潜码，STE（Straight-Through Estimator，直通估计器）在前向传递中将连续值取整为整数（$\mathrm{STE}(\cdot) = \mathrm{round}(\cdot)$），在反向传递中直接透传梯度（$\frac{\partial \mathrm{STE}}{\partial \hat{\mathbf{l}}_i} \approx 1$），使不可微的取整操作变得可微。
\textbf{方法信息：} Google / DeepMind 参与了这篇论文。
\textbf{设计动机与边界：} 相对原始 3DGS 以独立高斯和图像重建误差为主的设定，Splat4D 把问题落在以下可实现的改动上：Splat4D框架包含三个核心阶段。(1)粗4D高斯生成：将单目视频通过多视角扩散模型转换为多视角图像序列，利用预训练大规模前馈3DGS模型(LGM/Splatter Image)获取跨视角的空间和深度特征，处理为splat图像并反投影生成粗4D高斯表示。(2)时空精炼：通过多视角渲染识别时空不一致区域，使用视频扩散模型和非对称U-Net进行精炼，改善时空一致性。(3)泛化3D高斯场预测器：通过微调U-Net学习通用的3D高斯场预测能力，支持多种应用场景。框架支持从文本或图像输入生成4D内容、4D人体生成以及文本引导的内容编辑。分解场景为多个3D高斯分布同时捕获空间结构和时间演化。其设计目的不是单独增加一个模块，而是用上述表示、监督或渲染约束消除原流程在目标任务中的具体失配。提出Splat4D框架，通过多视角渲染、不一致性识别、视频扩散模型和非对称U-Net精炼实现高保真4D内容生成。利用大规模前馈3DGS模型获取空间深度特征，支持文本/图像条件4D生成、4D人体生成和文本引导编辑多种应用。适用时还应保留论文给出的边界：多阶段流水线计算成本高；对复杂宽松衣物的人体特征建模仍有挑战；视频扩散模型可能引入帧间不一致。
% METHOD: 2511.05152
\paragraph{Splatography（2025.11）\cite{2511_05152}}
Splatography 针对电影制作里的稀疏多视角动态重建。片场数据常只有少量同步相机，前景可能包含火焰、烟雾、透明物体和快速运动道具，而背景通常相对稳定但存在灯光、设备和遮挡变化。普通 4DGS 把前景和背景放进同一个变形场，容易让透明边缘、暗部背景和临时物体互相污染。Splatography 的第一步是在 $t=0$ 用稀疏掩码把初始点云拆成 foreground Gaussians 和 background Gaussians，并分别预训练 canonical foreground/background，使两者都锚定到同一个参考时间。

动态阶段使用差异化变形。前景分支使用 hex-plane deformation，预测位置、旋转和颜色变化，因为前景承担大部分复杂运动和外观变化；背景分支通常只预测位置变化，避免暗背景或片场设备在高维外观变形中被过拟合。训练时，前景 mask loss 会随机替换背景 alpha，压制前景 floaters；背景边缘则使用 blurred mask 约束，避免掩码边界过硬造成空洞。方法还加入 modified temporal opacity、正则项和 reference-free densification：当透明或动态纹理在参考帧里没有充分覆盖时，系统允许在后续时间补充基元。相对一般 4DGS，它的增量不是单个网络结构，而是把电影制作中的“前景可动、背景稳定、透明边缘难重建”写成两套 canonical 高斯和两套变形策略。
\textbf{实现细节：} Splatography（2025.11） 面向制片场景的稀疏多视角动态重建问题。在常规 4DGS 中，hex-plane 变形场 $G' = G + \Delta(G|t)$ 对所有高斯统一建模全部动态参数 $(\Delta x, \Delta r, \Delta s, \Delta c, \Delta\alpha)$，但在稀疏视角下会导致背景过度重建、前景欠重建——因为背景占据更大的图像面积，训练时梯度信号被背景主导。Splatography 的核心改动是将初始点云通过 $t{=}0$ 时刻的一组稀疏掩码 $M^{*} \in \mathbb{R}^{H \times W}$ 分割为前景 $G_f$ 和背景 $G_b$ 两个独立表示，并分别设计变形场：前景模型 $G'_f = \delta_f(G_f|t)$ 学习完整的 $(\Delta x, \Delta r, \Delta c)$，而背景模型 $G'_b = \delta_b(G_b|t)$ 仅学习 $(\Delta x)$。这一非对称设计源于制片实践的先验假设：表演区域（背景）灯光昏暗且人员运动有限，无需学习旋转和颜色变化；而前景包含丰富的纹理、透明和动态效果，需要更强的表达能力。
\textbf{方法信息：} 代码：\url{https://github.com/azzarelli/splatography-3dv-2026}。
\textbf{设计动机与边界：} 相对原始 3DGS 以独立高斯和图像重建误差为主的设定，Splatography 把问题落在以下可实现的改动上：方法包含两个核心创新：(1)规范前景-背景分离：使用t=0帧的稀疏掩码将初始点云分离为前景和背景高斯，锚定t=0进行规范预训练，开发掩码损失函数处理前景浮点和边缘伪影；(2)差异化变形场：前景变形场学习颜色、位置和旋转变化（因包含多样动态特征），背景变形场仅学习位置变化（电影制作中背景通常较暗且少动态）。引入时间不透明度函数、正则化和无参考密集化策略减少欠重建。其设计目的不是单独增加一个模块，而是用上述表示、监督或渲染约束消除原流程在目标任务中的具体失配。提出Splatography方法，将规范高斯和变形场分离为前景和背景组件，使用稀疏掩码初始化，分别训练不同损失函数和变形参数。在3D场景上达到比SOTA高3dB PSNR且模型大小减半的效果，无需密集掩码监督即可生成包含透明和动态纹理的分割动态重建。适用时还应保留论文给出的边界：依赖t=0帧的稀疏掩码质量；对前景中包含大量反射/透明物体的极端场景可能仍有困难；电影场景的特定约束限制了通用性。
% METHOD: 2512.14352
\paragraph{HGS（2025.12）\cite{2512_14352}}
把4D高斯的紧凑性建立在静态动态二值分解上。它为每个高斯引入静/动态指示器：动态基元维护随时间变化的位置、不透明度和旋转，尺度和颜色保持相对稳定并用时间函数建模；静态基元则共享一套时间不变参数，不再为每一帧存储冗余状态。初始化时，方法通过交错帧估计哪些基元长期稳定、哪些确实需要动态参数；训练上先重建整体时空场，再精修静态/动态边界，避免边界附近因为二值分解而闪烁。相对于把所有高斯都变形的4D-GS，HGS的关键改变是让静态区域完全退出时间建模，从而换来显著模型压缩，并面向4K实时和VR渲染。
\textbf{方法信息：} Microsoft Research 参与了这篇论文；当前未确认该论文完整官方代码。
\textbf{设计动机与边界：} 相对原始 3DGS 以独立高斯和图像重建误差为主的设定，HGS 把问题落在以下可实现的改动上：HGS框架：（1）混合高斯基元：将3DGS扩展到4D时空域，每个高斯参数化时间变化的位置、不透明度、旋转和不变尺度和颜色。通过二值指示器分离静态和动态基元。（2）静态-动态分解(SDD)：动态基元独立维护时间变化参数，用多项式插值参数化。静态基元共享统一的不变参数集，显著减少冗余。交错帧初始化方案分别获取两类基元。（3）两阶段训练策略：专门针对显式模型设计，增强静态-动态边界的时序连贯性。第一阶段关注整体重建，第二阶段精炼边界区域。其设计目的不是单独增加一个模块，而是用上述表示、监督或渲染约束消除原流程在目标任务中的具体失配。提出静态-动态分解(SDD)策略，动态区域用时间依赖RBF、静态区域共享不变参数。模型体积减少最多98%，RTX 3090上4K分辨率达125 FPS，RTX 3050上160 FPS，已集成到VR系统。适用时还应保留论文给出的边界：静态-动态分解依赖初始化质量；对复杂运动模式的参数化能力有限；两阶段训练增加复杂度
% METHOD: 2607.01202
\paragraph{World from Motion（2026.07）\cite{2607_01202}}
World from Motion 处理的是单目动态重建中最难补的一类缺失：相机没看见的动态区域、遮挡后方的外观，以及动态物体在新视角下应该如何运动。传统几何驱动方法如 MoSca、WorldTree 或 Shape-of-Motion 能在已观测区域保持较强几何一致性，但对未观测区域只能外推；纯视频生成模型能补图像，却不保证多视角和跨时间一致。World from Motion 的改变是把动态 3DGS 作为视频生成模型的条件，再把生成结果蒸馏回 3DGS，形成“重建给生成提供几何，生成给重建提供新观察”的闭环。

具体流程分两阶段。第一阶段先用现有单目动态重建方法得到初始动态 3DGS，虽然它能渲染输入附近视角，但会有空洞、伪影和错误 motion。系统随后沿输入相机轨迹和目标相机轨迹渲染 dense pixel-aligned conditioning，包括外观、几何、遮挡相关信息和 3D scene motion；这些条件输入视频生成模型，让模型在目标视角生成时间一致的视频帧。第二阶段不是直接把生成帧当最终 2D 输出，而是把这些虚拟观察作为监督，重新优化单一动态 3DGS，使新出现的外观、几何和运动被吸收到同一个可自由渲染的 4D 表示中。

相对 ViDAR、CAT4D 这类把生成模型用于视频或多视角补全的方法，World from Motion 的重点在于最终对象仍是显式动态 3DGS，因此可以继续做自由视角渲染和 3D motion 查询。它也不同于只对静态 3DGS 做 generative enhancement 的工作，因为条件和蒸馏都包含时间维度与动态运动。当前本地证据显示 NVIDIA 参与了这篇论文，项目页已公开，但尚未记录到完整公开代码。
\textbf{实现细节：} World from Motion。World from Motion（2026.07） 可作为"Generative completion for dynamic Gaussian reconstruction"中的新锚点候选。相较于 Dynamic 3DGS，主要改变是：视频模型条件于外观、几何和三维运动渲染，并把生成结果蒸馏回动态 3DGS。高斯基元仍承担可渲染动态场景表示。其设计动机是：构建 aligned multiview video pairs 和 dynamic 3DGS 数据集，训练模型沿输入和目标相机轨迹补全外观、几何和运动。具体而言，初始单目重建会产生伪影和缺失区域，方法用 dense pixel-aligned renderings 条件化视频模型生成修正视图，并将新观察区域和运动重新蒸馏回动态 Gaussian 表示。这一锚点判断来自当前分类中的机制关系：候选锚点：生成路线，且相关性 0.95；后续方法可围绕 将视频模型生成结果蒸馏回单一动态 3D Gaussian 表示。继续比较和扩展。适用边界是：依赖训练数据覆盖和视频模型先验；生成内容可能引入幻觉，需要蒸馏阶段约束。
\textbf{技术实现补充：} World from Motion。World from Motion（2026.07） 可作为"Generative completion for dynamic Gaussian reconstruction"中的新锚点候选。相较于 Dynamic 3DGS，主要改变是：视频模型条件于外观、几何和三维运动渲染，并把生成结果蒸馏回动态 3DGS。高斯基元仍承担可渲染动态场景表示。其设计动机是：构建 aligned multiview video pairs 和 dynamic 3DGS 数据集，训练模型沿输入和目标相机轨迹补全外观、几何和运动。具体而言，初始单目重建会产生伪影和缺失区域，方法用 dense pixel-aligned renderings 条件化视频模型生成修正视图，并将新观察区域和运动重新蒸馏回动态 Gaussian 表示。这一锚点判断来自当前分类中的机制关系：候选锚点：生成路线，且相关性 0.95；后续方法可围绕 将视频模型生成结果蒸馏回单一动态 3D Gaussian 表示。继续比较和扩展。适用边界是：依赖训练数据覆盖和视频模型先验；生成内容可能引入幻觉，需要蒸馏阶段约束。
\textbf{方法信息：} NVIDIA 参与了这篇论文；项目页：\url{https://research.nvidia.com/labs/amri/projects/world-from-motion/}。
\textbf{设计动机与边界：} 相对原始 3DGS 以独立高斯和图像重建误差为主的设定，World from Motion 把问题落在以下可实现的改动上：初始单目重建会产生伪影和缺失区域，方法用 dense pixel-aligned renderings 条件化视频模型生成修正视图，并将新观察区域和运动重新蒸馏回动态 Gaussian 表示。其设计目的不是单独增加一个模块，而是用上述表示、监督或渲染约束消除原流程在目标任务中的具体失配。构建 aligned multiview video pairs 和 dynamic 3DGS 数据集，训练模型沿输入和目标相机轨迹补全外观、几何和运动。适用时还应保留论文给出的边界：依赖训练数据覆盖和视频模型先验；生成内容可能引入幻觉，需要蒸馏阶段约束。

\section{动态物体与场景分解}

动态场景中并非所有变化都应该进入同一个4D表示。旅游照片里的行人、手持视频中的遮挡物、动态视频里的实例对象，往往需要被分离、降权或单独赋予语义身份。本节方法继承3DGS或4D-GS，但把“什么是静态场景、什么是瞬态干扰、什么是可编辑实例”作为核心变量。
% METHOD: 2406.10373
\paragraph{Wild-GS（无约束照片集合的外观与瞬态分解）（2024.06）}
Wild-GS~\cite{2406_10373} 面向旅游照片等外观、光照、相机处理和遮挡物随图变化的非受控集合。每个高斯的颜色不再只由固定 SH 给出，而分解为固有材质、每图的全局光照/相机属性和点级反射局部变化；从参考图像提取的 tri-plane 在对应局部高斯处采样，把高频参考外观显式带入三维而非停留在图像特征中。

训练同时学习这些外观因子和高斯几何，二维 visibility map 下调瞬态遮挡区域的解释权，深度正则限制几何被外观变化吸收。目标视图渲染时组合目标相机/光照条件与局部高斯特征，再走 alpha splatting。它是静态场景的鲁棒外观建模而非时间形变模型；极端外观转换、错误可见性或不可靠深度仍会造成材质与几何的分解歧义。

% METHOD: 2407.04504
\paragraph{Segment Any 4D Gaussians（时间身份特征场的 4D 分割）（2024.07）}
Segment Any 4D Gaussians~\cite{2407_04504} 在已有 4D-GS 表示上提供跨时刻对象分割与编辑。直接给 canonical Gaussian 贴固定标签会受 Gaussian drifting 影响：一个基元在不同时间的可见归属可能不稳定。SA4D 利用视频目标跟踪得到稀疏、带噪的身份 mask，并训练轻量 temporal identity feature field 预测时间相关的高斯身份特征。

将身份特征按当前时间的高斯可见性 splat 到图像后可得到任意视角/时刻的 mask；所有时间的结果再汇总成 Gaussian identity table，并经 4D segmentation refinement 去除漂浮和边界伪影。输出身份集合可用于移除、重着色或组合对象。该方法是分割接口而非重建替代，精度依赖底层 4DGS 轨迹和视频跟踪先验，在强非刚性或遮挡重现处仍可能失配。

% METHOD: 2407.08447
\paragraph{WildGaussians（DINO 不确定性与双层外观嵌入的野外 3DGS）（2024.07）}
WildGaussians~\cite{2407_08447} 解决无约束照片集合中的光照、相机处理和短暂遮挡。每张训练图像有 appearance embedding，每个高斯也有局部 appearance embedding；两者经轻量网络预测颜色的仿射调整，使同一三维材质可在不同观测外观下被解释。渲染时仍保留高斯的显式实时 splatting。

遮挡并不由大量临时高斯硬拟合：从图像提取的 DINO 特征经可学习映射给出像素不确定性，高不确定区域在 photometric loss 中被降权；该预测同时用渲染/观测的语义特征一致性训练。它将 NeRF-W 式外观变量迁到显式高斯，同时利用基础视觉特征识别 transient。工作针对静态 in-the-wild 场景，不表示时间形变；极端外观变化及 DINO 不能辨识的遮挡仍会使不确定性分解失效。

% METHOD: 2503.13176
\paragraph{DeGauss（2025.03）\cite{2503_13176}}
DeGauss 面向第一人称视频、随手拍图像集和长时间动态场景中的 distractor-free 重建。WildGaussians/SpotLessSplats 更倾向于识别不可靠像素并在损失中降权；DeGauss 则显式维护两套高斯：背景静态 3DGS 负责最终希望保留的干净场景，前景可变形 4DGS 负责人物、手部、物体交互和其它动态干扰。这样，动态内容不是简单被忽略，而是有自己的可渲染分支吸收，从而减少它们在背景中形成 floaters。

前景高斯的属性被扩展为 $G_f=\{x,s,r,\alpha,c,m_f,m_b,b\}$，其中 $m_f$、$m_b$ 是前景/背景概率属性，$b$ 是亮度控制属性。HexPlane 编码器按时间提取时空特征，多头解码器分别预测位置、旋转、尺度、opacity、颜色以及这些概率属性的变形。渲染时，前景分支不仅输出颜色，还 rasterize 出 raw foreground probability $M_f$ 和 background probability $M_b$，再归一化为 $P_f$、$P_b$；这个概率组合会抑制接近 0.5 的中间状态，推动高斯明确归入动态或静态解释。背景亮度控制分支用于吸收随手拍中的曝光、反光和亮度变化，避免动态分支因为亮度误差而过度吞掉背景。相对 HybridGS 的单图 2D 瞬态层，DeGauss 更适合长视频和第一人称交互，因为动态前景本身是时空可变形的 3D/4D 表示。
\textbf{方法信息：} 代码：\url{https://github.com/batfacewayne/degauss}。
\textbf{设计动机与边界：} 相对原始 3DGS 以独立高斯和图像重建误差为主的设定，DeGauss 把问题落在以下可实现的改动上：DeGauss采用解耦设计：(1) 前景可变形高斯：扩展高斯属性集包含前景概率m\_f、背景概率m\_b和亮度控制b。基于Hexplane的编码器提取时空特征，多头解码器预测每个属性的变形。前景高斯在时间维度上变形以捕捉动态内容。(2) 概率掩码组合：通过可微分渲染计算原始前景概率M\_f和背景概率M\_b，归一化为P\_f和P\_b。概率公式自然抑制中间值(接近0.5)，推动预测趋向0或1，产生干净的动态-静态分解。(3) 背景亮度控制：处理休闲采集中的显著光照变化，减少几何和外观的歧义。(4) 自监督：无需复杂启发式或大量监督，从随机图像集合到长时间动态第一人称视频均能鲁棒泛化。其设计目的不是单独增加一个模块，而是用上述表示、监督或渲染约束消除原流程在目标任务中的具体失配。提出DeGauss自监督框架，通过解耦的动态-静态高斯设计实现鲁棒的无干扰物3D重建。核心创新：(1) 前景可变形高斯和背景静态高斯独立优化，通过概率掩码协调组合；(2) 基于Hexplane的时空变形模块预测前景高斯的变形；(3) 亮度控制属性处理休闲采集中的光照变化。在NeRF-on-the-go、ADT、AEA、Hot3D和EPIC-Fields五个基准上均超越现有方法。适用时还应保留论文给出的边界：自监督方法的分解精度受限于损失设计；前景可变形高斯增加了训练复杂度；对极端遮挡和快速运动的处理仍有改进空间。
% METHOD: 2604.01678
\paragraph{Director（2026.04）\cite{2604_01678}}
把4DGS从渲染推进到实例级理解。每个高斯不仅有几何和外观，还带可学习语义特征；两个MLP头分别预测实例身份和语言对齐嵌入。初始化上，SAM3提供时序一致实例掩码，MLLM生成实例描述并帮助跨视角对齐身份；动静分离后，背景跨时间聚合，前景则用光流引导的统一4D表示建模位置、旋转、颜色、不透明度、尺度和语义。SDF约束用于增强表面和时序稳定性。相对SA4D这类分割后处理，Director把实例语义嵌入到高斯优化本身。
\textbf{实现细节：} Director（2026.04） 的核心主张是动态重建与场景理解不应遵循"先重建再理解"的分离流水线，而应在实例一致的高斯基元层面联合优化。现有方法要么只关注光度一致性（如 4DGS、GaussianFlow），要么仅在后期附加语义（如 LERF、4DLangSplat），导致高度动态场景中出现身份漂移、前景-背景纠缠和语义关联不稳定。Director 的解决方案是将实例级身份和语言对齐的语义嵌入直接作为每个高斯的可学习属性，与运动和外观联合优化。在初始化阶段，Director 利用 SAM3 的记忆库机制，以文本提示（如"basketball player""basketball"）引导分割，在单视频内获取时间一致的实例掩码。为建立跨视角身份一致性，将所有视角的首帧拼接为规范视频统一分割——由于相邻摄像机视点间距小、画面重叠度高，规范掩码天然对齐，再通过 IoU 匹配建立跨视角的一一对应关系，传播至全部帧，获得时空一致的实例掩码 $M^{v,t}_d$。
\textbf{技术实现补充：} Director（2026.04） 的核心主张是动态重建与场景理解不应遵循"先重建再理解"的分离流水线，而应在实例一致的高斯基元层面联合优化。现有方法要么只关注光度一致性（如 4DGS、GaussianFlow），要么仅在后期附加语义（如 LERF、4DLangSplat），导致高度动态场景中出现身份漂移、前景-背景纠缠和语义关联不稳定。Director 的解决方案是将实例级身份和语言对齐的语义嵌入直接作为每个高斯的可学习属性，与运动和外观联合优化。在初始化阶段，Director 利用 SAM3 的记忆库机制，以文本提示（如"basketball player""basketball"）引导分割，在单视频内获取时间一致的实例掩码。为建立跨视角身份一致性，将所有视角的首帧拼接为规范视频统一分割——由于相邻摄像机视点间距小、画面重叠度高，规范掩码天然对齐，再通过 IoU 匹配建立跨视角的一一对应关系，传播至全部帧，获得时空一致的实例掩码 $M^{v,t}_d$。同时用多模态大语言模型为每个实例生成双层语义描述：时间不变的全局描述 $C_i^{\mathrm{global}}$（在首帧用 SAM3 掩码渲染红色轮廓高亮实例，提示 Qwen3-VL 描述整体外观和动作特征）和逐帧动态描述 $C_i^t$（将所有视角在第 $t$ 帧的观测聚合为空间一致性视频剪辑，连同全局描述和时间提示输入 MLLM）。两种描述经 LLM 编码后，再通过轻量级自编码器压缩为紧凑的 6 维句子嵌入 $e$，作为语言对齐的语义监督信号。在场景分解上，Director 将场景显式分为静态背景 $G_{\mathrm{bg}}$ 和动态前景 $G_{\mathrm{fg}}$。背景层通过跨时间和空间维均匀采样 GT 图像全局优化，损失函数为 $E^{\mathrm{bg}}_{\mathrm{color}} = \|M_{\mathrm{bg}} \odot (\hat{C} - C)\|_1$（仅在背景掩码区域计算），配合等渗正则 $E_{\mathrm{iso}}$ 和尺寸正则 $E_{\mathrm{size}}$ 防止高斯过度拉伸（$\lambda_{\mathrm{iso}} = 0.0005$、$\lambda_{\mathrm{size}} = 0.02$），训练 20000 次迭代。背景层不仅提供全局参考，在后续每帧中还持续优化外观属性以适应光照和阴影变化。
\textbf{方法信息：} ByteDance 参与了这篇论文；当前未确认该论文完整官方代码。
\textbf{设计动机与边界：} 相对原始 3DGS 以独立高斯和图像重建误差为主的设定，Director 把问题落在以下可实现的改动上：核心设计是在高斯基元级别强制实例一致性。每个高斯配备可学习语义特征，通过两个MLP头预测实例身份和语言对齐嵌入。具体流程：(1)多模态实例初始化：SAM3提取时序一致掩码，MLLM生成实例描述，通过首帧拼接实现跨视角身份一致性；(2)动静分离：背景跨时间聚合优化，前景使用光流引导的统一4D语义表示建模位置、旋转、SH颜色、不透明度、缩放和语义特征；(3)几何约束：引入SDF约束和正则化强制表面连续性，增强前景建模的时序一致性。其设计目的不是单独增加一个模块，而是用上述表示、监督或渲染约束消除原流程在目标任务中的具体失配。首个将实例一致语义嵌入4DGS建模的统一框架，同时实现高保真渲染、实例分割和开放词汇查询。提出光流桥接2D到4D运动初始化、几何感知SDF约束和实例一致的正则化。适用时还应保留论文给出的边界：依赖SAM3和MLLM的掩码质量；对极端遮挡和快速运动仍可能失败；计算开销较大。

\section{运动建模与跟踪}

这一节的方法把动态高斯当成运动估计工具。与只要求渲染质量的方法不同，它们更关心高斯轨迹是否可解释、是否能对应到点跟踪、光流或运动迁移。技术上，它们往往继承Dynamic 3D Gaussians的持久粒子思想，或在Deformable 3DGS上加入显式运动监督。
% METHOD: 2308.09713
\paragraph{Dynamic 3D Gaussians（持久粒子的逐时刻优化）（2023.08）}
Dynamic 3D Gaussians~\cite{2308_09713} 将动态重建表述为持久的多视图合成与跟踪，而非用一个隐式变形网络查询任意时间。每个高斯在每个时刻有自己的中心和四元数旋转，但尺度、RGB、opacity logit 与背景标记跨时间固定；因此同一编号的高斯始终代表同一局部物质，时变中心和旋转就是其 6-DoF 轨迹。论文实验关闭了视角相关 SH，以避免外观随时间自由变化破坏这一对应假设。

\textbf{在线拟合与渲染。} 给定各时刻、多相机的图像和标定，第一时刻按 3DGS 优化全部属性并进行 densification。之后以上一时刻的高斯为初始化，冻结颜色、尺度、opacity 与背景，只更新当前位置和旋转；每一步仍以标准可微 Gaussian rasterizer 投影协方差、深度排序并前到后合成，与当前多相机图像比较。位置和旋转还用上一帧的速度前推初始化，优化器动量在新时刻重置；后续时刻不再增加高斯。

\textbf{从重建到跟踪的约束。} 图像误差本身不能阻止相似颜色区域的粒子互换。方法在初始邻域中取 kNN，加入短程 local-rigidity：邻居在当前帧的相对坐标应等于前帧相对坐标经中心粒子旋转后的结果；再加入邻域旋转变化相似项，以及保持初始邻居距离的长期 isometry 项。前景/背景掩码还能约束背景不运动。训练后，任意空间点由影响最大的高斯的局部坐标表示，随该高斯的时变坐标系转换即可得到跨时刻对应。代价是轨迹显式存储且依赖多视图在线优化；颜色、遮挡或拓扑突变造成的身份歧义仍需要这些物理先验缓解。

% METHOD: 2409.02104
\paragraph{DynOMo（2024.09）\cite{2409_02104}}
面向在线单目点追踪。它扩展Dynamic 3D Gaussians，让每个高斯除了静态属性外，还带有动态特征和时间偏移，并通过deformation MLP预测帧间运动；系统逐帧处理单目视频，同时估计相机位姿和优化动态高斯模型。为了让点轨迹从重建中自然涌现，DynOMo在图像重建之外加入更强的图像特征、深度监督和语义前景/背景分离，并用特征相似性调节运动正则权重：外观和语义更相近的高斯更倾向于共享运动，不相干区域则不被硬性拉在一起。相对离线动态高斯，DynOMo强调机器人等场景中的边看边建，输出不仅是新视角图像，还包括2D和3D点轨迹。
\textbf{方法信息：} NVIDIA 参与了这篇论文。
\textbf{设计动机与边界：} 相对原始 3DGS 以独立高斯和图像重建误差为主的设定，DynOMo 把问题落在以下可实现的改动上：DynOMo的核心技术：(1) 动态高斯场景表示：基于D-3DGS扩展，每个高斯除了静态属性外还有动态特征和时间偏移，通过 deformation MLP预测帧间运动。(2) 在线追踪管线：逐帧处理单目视频，同时估计相机位姿和优化动态高斯模型。(3) 重建信号增强：使用更强的图像特征和深度监督提升重建质量。(4) 基于语义的前景/背景分离：支持相机追踪。(5) 相似性增强正则化：通过特征相似性引导的运动正则化权重机制，使点轨迹自然涌现而无需对应级别监督。其设计目的不是单独增加一个模块，而是用上述表示、监督或渲染约束消除原流程在目标任务中的具体失配。提出DynOMo，首个从无位姿单目相机输入进行在线2D和3D点追踪的方法，通过3DGS动态场景在线重建实现，无需对应级别监督即可产生点追踪轨迹适用时还应保留论文给出的边界：作为首个基线方法，性能仍有提升空间；对快速运动和大变形的处理可能不够鲁棒；单目深度估计的不确定性会影响3D追踪精度
% METHOD: 2410.07707
\paragraph{MotionGS（2024.10）\cite{2410_07707}}
MotionGS 是 Deformable 3DGS 路线上的显式运动监督方法。Deformable 3DGS 主要依靠渲染误差反推变形场，在单目视频里会遇到两个歧义：物体快速运动时，颜色误差不一定能告诉高斯应该往哪里移动；相机位姿不准时，模型还可能把相机误差解释成物体变形。MotionGS 的改变是把光流变成可用的运动约束，但不直接拿原始光流监督，因为原始光流同时包含相机运动和物体运动。

具体流程分成两条数据流。2D 数据流先用现成光流网络估计相邻帧光流；同时，利用当前渲染深度和相邻相机位姿，假设场景静止，计算仅由相机运动产生的 camera flow。总光流减去 camera flow 后得到 motion flow，即在固定相机视角下物体自身运动造成的像素位移。3D 数据流则沿用规范高斯加变形网络，把高斯从时间 $t$ 变形到 $t+1$，并在 rasterizer 中投影出 Gaussian flow。训练时加入 $L_1$ 运动流损失，让 Gaussian flow 对齐 motion flow，从而把“物体应该怎样动”接到变形场梯度上。

相机位姿也不是完全固定的。MotionGS 在 COLMAP 初始位姿上增加小的旋转和平移残差，通过相机姿态精修模块交替优化高斯和相机；更新后的相机又会反过来让 camera flow 更准，减少运动流分解的误差。相对 Deformable 3DGS，它不是换掉规范空间框架，而是补上“光流解耦 + Gaussian flow 对应 + 姿态残差优化”三件事。论文在 NeRF-DS 和 HyperNeRF 等单目动态基准上展示了对快速、不规则运动的改善。
\textbf{锚点细化：} MotionGS：针对变形场仅靠外观重建损失监督在物体不规则运动（突然加速、急停）时陷入局部最优、光流同时包含相机运动和物体运动无法直接约束变形的局限，用``光流解耦模块 + 高斯流对应''替代直接光流监督，核心思想是通过深度图和位姿解析计算相机流，从原始光流减去得到纯物体运动的运动流，再在固定视角下投影高斯变形得到高斯流，与运动流 $L_1$ 对齐——避免相机运动造成的错误梯度。
	extbf{实现细节：} MotionGS 在规范空间维护高斯，并用变形网络把位置和形状映射到每个时间步。它先以当前深度和相邻相机位姿计算由相机运动产生的二维流，再从光流网络预测的总流中相减，得到只描述物体运动的残差流。渲染器从变形后的高斯产生对应的二维流，并以运动流一致性损失监督变形；同时在 COLMAP 位姿上学习小残差，交替更新相机和 3DGS 参数。这使训练能把姿态误差与真实物体运动分开，推理时则以时间编码、变形网络和标准 alpha 混合生成动态视图。
	extbf{实现细节：} MotionGS 在规范空间维护高斯，并用变形网络把位置和形状映射到每个时间步。它先以当前深度和相邻相机位姿计算由相机运动产生的二维流，再从光流网络预测的总流中相减，得到只描述物体运动的残差流。渲染器从变形后的高斯产生对应的二维流，并以运动流一致性损失监督变形；同时在 COLMAP 位姿上学习小残差，交替更新相机和 3DGS 参数。这使训练能把姿态误差与真实物体运动分开，推理时则以时间编码、变形网络和标准 alpha 混合生成动态视图。
\textbf{方法信息：} 代码：\url{https://github.com/ruijiezhu94/motiongs}；GitHub 关注度较高，约 155 stars；引用很多，Semantic Scholar 约 72 次。



% METHOD: 2511.14848
\paragraph{Gaussian See, Gaussian Do（2025.11）\cite{2511_14848}}
把源视频中的语义运动迁移到目标3D形状。它先用Stable Video Diffusion通过条件反演提取源运动嵌入，并用锚点视角感知机制保证跨视角一致；然后把该运动嵌入施加到静态目标形状渲染帧，生成目标物体的运动参考视频；最后用这些有噪监督视频训练4D高斯，控制点和MLP变形网络负责从时间和控制点编码预测3D变形，课程学习逐步扩大运动范围。相对普通动态重建，它的目标不是复原输入运动场景，而是把语义运动作为先验驱动新的高斯对象动起来。
\textbf{实现细节：} 传统运动迁移依赖骨架绑定（rigging），无法处理没有预定义关节结构的一般物体；2D 视频域的语义运动迁移（如 MotionClone）虽然能将运动与外观解耦，但缺乏 3D 一致性。Gaussian See, Gaussian Do（2025.11） 的核心改进是将 2D 条件反演提升为多视角一致的 3D 运动迁移管线——不要求源和目标共享骨架，也不需要文本描述运动。具体做法分两阶段。第一阶段（结构化多视角运动反演）：从 $N$ 视角的源视频 $V_S$ 中提取运动嵌入。为此，作者提出锚定视角感知嵌入机制：只优化 $K$ 个锚定嵌入 $\{m_k^{*}\}_{k=1}^K$，均匀分布在源视角角度范围上；对任意中间角度 $\phi \in [\phi_i, \phi_j]$，通过球面线性插值获得嵌入： 公式的含义是：角度 $\phi$ 处的运动嵌入由其两侧最近锚点的嵌入平滑混合而成，插值参数控制两个锚点的权重比例。第二阶段（视角感知语义迁移与 4D 整合）：渲染目标静态 3DGS 的随机视角图像，用对应的插值运动嵌入条件化冻结的视频扩散模型，生成目标物体在该视角下的 2D 监督视频。然后用多视角监督视频驱动 SC-GS 变形场重建 4D 场景。$w_{ik}$ 是高斯 $i$ 与控制点 $k$ 之间的蒙皮权重，由距离高斯核 $w_{ik} = \exp(-d_{ik}^{}2 / 2\sigma^{}2)$ 计算。该方法建立了首个语义 3D 运动迁移基准，在运动保真度和结构一致性上显著优于基线方法。
\textbf{方法信息：} NVIDIA 参与了这篇论文。
\textbf{设计动机与边界：} 相对原始 3DGS 以独立高斯和图像重建误差为主的设定，Gaussian See, Gaussian Do 把问题落在以下可实现的改动上：方法流程：(1)运动嵌入提取：从源视频通过条件反演(condition inversion)提取运动嵌入，使用SVD(Stable Video Diffusion)作为骨干。采用锚点视角感知运动嵌入机制，确保跨视角一致性并加速收敛。模型使用5个锚点、3000次反演迭代、15×5×1024大小的运动嵌入。(2)运动迁移：将运动嵌入应用到静态目标形状的渲染帧上，生成目标物体的运动参考视频。(3)4D高斯splatting重建：使用生成的参考视频监督动态3DGS重建。引入鲁棒的4D重建管线整合含噪监督视频。使用512个控制点（FPS初始化）和MLP变形网络编码时间和控制点到3D变形预测。可选的课程学习策略逐步扩大运动范围，解决大幅运动的学习困难。建立首个语义3D运动迁移基准。其设计目的不是单独增加一个模块，而是用上述表示、监督或渲染约束消除原流程在目标任务中的具体失配。提出首个语义3D运动迁移方法，实现无骨骼跨类别运动迁移。通过锚点视角感知运动嵌入机制确保跨视角一致性并加速收敛，建立首个语义3D运动迁移基准，运动保真度和结构一致性优于适配基线。适用时还应保留论文给出的边界：对高度关节化运动（踢腿、开合跳）效果差，逐视角运动反演会失败导致严重伪影；课程学习的必要性取决于运动幅度。
% METHOD: 2604.02586
\paragraph{TrackerSplat（2026.04）\cite{2604_02586}}
TrackerSplat 把点追踪模型作为动态高斯的运动初始化器。Deformable 3DGS 依赖渲染误差从零优化变形，大位移、遮挡和快速相机运动下很容易收敛慢或局部错误；TrackerSplat 则先让 2D tracker 给出像素轨迹，再把这些轨迹提升成高斯的 3D 运动。流程从第一帧开始：用 InstantSplat 或 COLMAP/3DGS 初始化高斯集合，再用 DOT、CoTracker 或 TAPIR 这类点跟踪器追踪初始帧像素在后续视频中的轨迹，代码实现中主要使用 DOT 管线。

对每个高斯，方法不逐像素求一个独立 3D 位移，而是把它投影 footprint 内的 2D 运动近似为仿射变换 $A,b$。Parallel weighted incremental least squares（PWI-LS）把 footprint 内像素轨迹按该高斯的 alpha contribution 加权，累积线性系统求解仿射参数；这样既避免重叠高斯互相抢像素，也不需要为每个高斯做昂贵的独立反演。随后系统在多视角中三角化更新高斯的 3D 位置、旋转和尺度；对无法稳定求解的高斯，用中值滤波和 8-nearest-neighbor 传播补运动。最后只做短程渲染精修，通常不再 densify，以免破坏追踪得到的身份。相对 Deformable 3DGS，它把“运动从哪里来”从纯梯度优化改成了“点轨迹初始化 + 加权仿射提升 + 轻量渲染细化”。
\textbf{实现细节：} 现有基于可微渲染的动态 3DGS 方法（如 D-3DGS）通过光度损失间接驱动高斯运动。TrackerSplat（2026.04） （SIGGRAPH Asia 2025）的核心诊断是：3DGS 的位置梯度计算被限制在有限局部邻域内，当物体经历大帧间运动并移出该区域时，位置梯度无法正确对齐物体运动方向，导致高斯褪色或错误变色（而非仅仅是光度损失不够强）。核心改动是：在训练之前，先用现有点跟踪模型（论文评估了 DOT、CoTracker 和 TAPIR，最终选用 DOT，因为其精度相当但速度显著更高）提取像素轨迹，通过并行加权增量最小二乘（PWI-LS）将 2D 像素轨迹直接转化为 3D 高斯的仿射运动，再经多视角三角化更新高斯的位置、旋转和缩放。设计动机是：显式的像素跟踪提供跨帧对应关系，比可微渲染的光度梯度更可靠地引导大位移运动，使得高斯在梯度优化之前已接近正确位置，训练过程更快更稳定。具体实现上，对于首帧中的每个 2D 高斯（协方差 $\Sigma_{2D}$、均值 $\mu_{2D}$），定义其运动为仿射变换 $[A|b]$，将第一帧的 2D 高斯映射到后续帧。PWI-LS 通过最小化加权最小二乘求解： 其中 $w_i = \alpha_i T_i$ 是像素 $i$ 对该高斯的不透明度-透明度权重（$\alpha_i$ 为该高斯在像素 $i$ 处的不透明度，$T_i$ 为 $\alpha$-混合 中的累积透明度；边界像素权重低，避免被无关运动误导），$P_i$ 和 $Q_i$ 是像素级别的增量统计量，可在 GPU 上并行累加。为避免数值不稳定，论文丢弃覆盖像素少于 3 个、$V_1$ 行列式低于 $10^{-12}$ 或累积 alpha 低于 $10^{-3}$ 的高斯运动。求解后得到 2D 仿射运动，再通过 SVD 三角化恢复 3D 均值 $\mu_{3D}'$（至少需要 2 个视角，为数值稳定性丢弃少于 3 个视角观察的高斯），通过多视角约束线性系统求解 3D 协方差 $\Sigma_{3D}'$（每个视角的 $\Sigma_{2D}'$ 提供 3 个约束，$\Sigma_{3D}'$ 有 6 个未知参数，至少需要 2 个不同视角），最后特征值分解提取更新后的旋转矩阵和缩放向量。
\textbf{技术实现补充：} 现有基于可微渲染的动态 3DGS 方法（如 D-3DGS）通过光度损失间接驱动高斯运动。TrackerSplat（2026.04） （SIGGRAPH Asia 2025）的核心诊断是：3DGS 的位置梯度计算被限制在有限局部邻域内，当物体经历大帧间运动并移出该区域时，位置梯度无法正确对齐物体运动方向，导致高斯褪色或错误变色（而非仅仅是光度损失不够强）。核心改动是：在训练之前，先用现有点跟踪模型（论文评估了 DOT、CoTracker 和 TAPIR，最终选用 DOT，因为其精度相当但速度显著更高）提取像素轨迹，通过并行加权增量最小二乘（PWI-LS）将 2D 像素轨迹直接转化为 3D 高斯的仿射运动，再经多视角三角化更新高斯的位置、旋转和缩放。设计动机是：显式的像素跟踪提供跨帧对应关系，比可微渲染的光度梯度更可靠地引导大位移运动，使得高斯在梯度优化之前已接近正确位置，训练过程更快更稳定。具体实现上，对于首帧中的每个 2D 高斯（协方差 $\Sigma_{2D}$、均值 $\mu_{2D}$），定义其运动为仿射变换 $[A|b]$，将第一帧的 2D 高斯映射到后续帧。PWI-LS 通过最小化加权最小二乘求解： 其中 $w_i = \alpha_i T_i$ 是像素 $i$ 对该高斯的不透明度-透明度权重（$\alpha_i$ 为该高斯在像素 $i$ 处的不透明度，$T_i$ 为 $\alpha$-混合 中的累积透明度；边界像素权重低，避免被无关运动误导），$P_i$ 和 $Q_i$ 是像素级别的增量统计量，可在 GPU 上并行累加。为避免数值不稳定，论文丢弃覆盖像素少于 3 个、$V_1$ 行列式低于 $10^{-12}$ 或累积 alpha 低于 $10^{-3}$ 的高斯运动。求解后得到 2D 仿射运动，再通过 SVD 三角化恢复 3D 均值 $\mu_{3D}'$（至少需要 2 个视角，为数值稳定性丢弃少于 3 个视角观察的高斯），通过多视角约束线性系统求解 3D 协方差 $\Sigma_{3D}'$（每个视角的 $\Sigma_{2D}'$ 提供 3 个约束，$\Sigma_{3D}'$ 有 6 个未知参数，至少需要 2 个不同视角），最后特征值分解提取更新后的旋转矩阵和缩放向量。
\textbf{方法信息：} 代码：\url{https://github.com/yindaheng98/trackersplat}。
\textbf{设计动机与边界：} 相对原始 3DGS 以独立高斯和图像重建误差为主的设定，TrackerSplat 把问题落在以下可实现的改动上：四阶段流程：(1)初始化：用InstantSplat为第一帧初始化3DGS；(2)点追踪：用现有点追踪模型提取初始帧每个像素在整个视频片段中的2D轨迹；(3)运动补偿：定义2D高斯运动为仿射变换[A|b]，通过并行加权增量最小二乘(PWI-LS)从像素轨迹求解每个高斯的仿射变换参数，利用多视角三角化更新3D高斯参数，引入运动正则化（中值滤波+邻域传播）缓解追踪误差；(4)细化：在输入视频片段上细化每帧参数。支持多GPU设备的并行管线处理。其设计目的不是单独增加一个模块，而是用上述表示、监督或渲染约束消除原流程在目标任务中的具体失配。首次将现有点追踪模型集成到3DGS动态场景重建中，通过预训练前的运动补偿有效处理大帧间位移，支持多设备并行处理而保持质量。引入PWI-LS高效求解高斯运动。适用时还应保留论文给出的边界：依赖点追踪模型的精度；对极端遮挡区域的追踪可能失败；多设备并行的一致性仍有优化空间。

\section{其他动态场景方法}

最后一组方法把动态3DGS接入特殊成像、概率建模、模糊去除、刚体分解和体积捕捉系统。它们的共同点是任务约束已经超出普通RGB新视角合成，因此需要把传感器物理、贝叶斯不确定性、曝光模型或工程流水线写入高斯表示。
% METHOD: 2311.18561
\paragraph{Periodic Vibration Gaussian（周期振动先验的城市动态重建）（2023.11）}
Periodic Vibration Gaussian~\cite{2311_18561} 面向城市视频中的树叶、旗帜等小幅周期性运动。普通动态高斯若为每个时刻自由拟合，容易把稀疏观测下的重复振动当成互不相关的位移；该方法把静态 3DGS 与周期位移模型结合，使高斯中心在时间上的变化由可学习振幅、频率和相位控制，而非逐帧独立存储。

\textbf{表示与训练。} 高斯的静态几何和外观仍按 3DGS 参数化；对被解释为动态的基元，时间查询生成周期位移，再以变换后的中心进入普通投影与 alpha 合成。训练从图像重建监督共同优化高斯和振动参数，并利用时间一致性/周期约束区分稳定背景与振动部分。因而渲染任意时间只需计算该时间的解析运动并执行 splatting，无需为每帧存新场景。该假设对反复、小位移的城市动态有效，却不适合非周期大形变、突然进入物体或长时间不规则运动。

% METHOD: 2312.00112
\paragraph{DynMF（共享基轨迹的神经运动分解）（2023.12）}
DynMF~\cite{2312_00112} 认为动态场景的逐点运动可由少量共享轨迹线性组合，而非每个高斯各自学习完整 deformation field。它学习只以时间为输入的小型网络来产生若干 basis trajectories，并为点/高斯绑定运动系数；某点在时刻 $t$ 的位移由这些轨迹按其系数相加得到。这样局部或重叠的运动被迫共享少数运动成分，从结构上约束单目动态重建的歧义。

\textbf{稀疏分解和渲染。} 系数的 sparsity loss 促使一个点只使用少量基轨迹，使不同运动可被分离和独立开关。位移更新后的 3D Gaussians 沿用标准 splatting 进行新视角合成；重建误差、运动系数稀疏性及相应优化项共同更新 basis、系数与高斯属性。与 canonical MLP 的逐坐标变形相比，网络只需按时间生成很少的全局运动函数，推理时先计算轨迹组合再光栅化。其表达能力由基数限制：基过少会把不同非刚性动作合并，基过多则失去分解带来的正则与可控性。

% METHOD: 2312.00583
\paragraph{DeformGS（高形变物体的 Gaussian scene flow）（2023.12）}
DeformGS~\cite{2312_00583} 的目标不是只渲染动态视频，而是为可变形物体操作恢复可用的三维 scene flow。它使用 3D Gaussian 表示物体，同时以跨时刻的高斯对应和运动估计把连续变形显式化；渲染误差提供外观监督，流的一致性与局部几何约束避免仅靠图像匹配产生不合理的粒子运动。

\textbf{从高斯到操作流。} 系统在观测时刻优化高斯的位置、形状和外观，并估计/传播每个高斯的运动，使它们可由当前状态变换到后续状态；将高斯场渲染为图像、深度或对应量以和输入比对。对于操作任务，scene flow 进一步提供把接触或抓取点随形变迁移的映射，而不是仅输出新视图。它与单目 canonical deformation 方法的区别是把运动场作为服务于交互的显式输出；这一优势也意味着它对标定、遮挡和高形变区域的对应质量敏感，错误 flow 会直接影响下游操作。

% METHOD: 2312.01196
\paragraph{Neural Parametric Gaussians（低秩点模板驱动的局部高斯）（2023.12）}
Neural Parametric Gaussians~\cite{2312_01196} 先用低秩参数化点模型解决单目非刚体重建的时序欠约束，再在其局部坐标系内放置精细高斯。阶段一将每一帧点集表示为少数形状基的线性组合；时间网络预测组合系数，因而基数 $K$ 直接控制变形的刚性。点模型以 mask Chamfer、邻域距离保持和由光流对齐的随机点描述子渲染项优化，得到跨时刻连续的粗模板。

\textbf{局部 volume 上的高斯。} 对每个模板点及其近邻建立 oriented local volume，方向由固定邻居构成的局部三角形随形变更新。一个 volume 可容纳不同数量的高斯；高斯位置不是世界坐标，而是对该 volume 邻居的 softmax barycentric 权重，旋转再乘局部 frame。于是点模板移动时，所有附着高斯自然随局部形变运动，而尺度、局部旋转和 SH 属性在时间上共享。

\textbf{细化与渲染。} 各 volume 的高斯用原 3DGS rasterizer 合并渲染；densification/pruning 将拆分的高斯留在原 volume，并微扰其权重。第二阶段以 $\ell_1$ 加 D-SSIM 优化权重、旋转、尺度和 SH，必要时弱化地微调第一阶段。它避免对大量高斯逐个运行时变网络，但低秩模板和固定邻域也是边界：基数小会压制真实非刚性，基数太大又会削弱正则与新视角一致性。

% METHOD: 2312.14937
\paragraph{SC-GS（稀疏控制点驱动的可编辑动态高斯）（2023.12）}
SC-GS~\cite{2312_14937} 不为每个高斯独立预测流，而以远少于高斯数的可学习 control points 表示运动。控制点位于紧凑的运动空间，周围高斯根据距离/邻域权重从控制点的位移与旋转获得变形；这把真实场景“运动稀疏、空间连续、局部近刚性”的假设直接写入表示。

训练从单目视频的图像误差优化高斯、控制点和控制权重，并以局部一致性约束防止控制点产生噪声轨迹。给定时间，先插值控制点变换、再更新高斯并照常 splat；选择或移动控制点也给编辑提供接口。与逐高斯 flow 相比，它减少动态自由度并改善可控性，但控制点数量、覆盖与邻域关联决定能否表达细小独立运动，跨越控制域的强形变仍可能被过度平滑。

% METHOD: 2312.07504
\paragraph{COLMAP-Free 3D Gaussian Splatting（联合相机配准与高斯优化）（2023.12）}
COLMAP-Free 3D Gaussian Splatting~\cite{2312_07504} 处理 3DGS 对 SfM/COLMAP 初始相机位姿和稀疏点云的依赖。它把相机外参作为优化变量，与高斯场一起从图像监督中恢复；训练不把预先成功的 COLMAP 当作输入前提。

方法以可学习相机位姿和初始高斯表示交替/联合优化：当前位姿决定高斯的投影与可微 splatting，渲染 RGB 与观测图像的误差反传到场景和相机。为缓解初始几何与初始 pose 彼此依赖的鸡生蛋问题，论文采用由粗到细的优化与对应的正则/初始化策略，逐步提高相机和场景自由度。输出仍是普通可实时光栅化的 3D Gaussian 场，而非动态变形模型；它降低预处理门槛，但没有消除 pose--geometry 的非凸性，纹理重复、弱纹理或严重错误初始化仍可能导致联合优化陷入局部解。

% METHOD: 2312.13308
\paragraph{SWinGS（滑动窗口动态 Gaussian 优化）（2023.12）}
SWinGS~\cite{2312_13308} 面向长视频的动态 Gaussian 重建。全序列同时优化会使时空参数、显存和密度控制持续增长；该方法以时间滑动窗口只维护局部连续帧的动态表示，并把窗口间共享/传播的高斯作为下一窗口初始化，从而把长序列分解为可管理的局部优化问题。

窗口内仍由图像重建驱动可微 Gaussian splatting，并通过局部时间一致性处理运动；窗口向前推进时，已稳定部分被冻结或压缩，重叠帧把新旧表示对齐，避免每段单独重建产生接缝。该思路的贡献是训练和内存调度，不是新的 Gaussian 核或全局 canonical deformation；其连续性取决于窗口长度、重叠范围和传播质量，跨越很长时间的身份或突然大位移仍可能累积误差。

% METHOD: 2312.13763
\paragraph{Align Your Gaussians（组合扩散蒸馏的 Text-to-4D）（2023.12）}
Align Your Gaussians~\cite{2312_13763} 从文本生成动态 4D 内容，而非从已拍摄视频重建。它以各向同性 3D Gaussians 和坐标--时间 MLP 位移场表示场景；位移乘以随时间从零增长的因子，保证首帧保持静态基形。局部刚性正则使邻近高斯的变形相近，显式表示也允许多个动态对象组合。

训练分两阶段。第一阶段只优化静态高斯，用多视角文本扩散先验和文本到图像先验的 score distillation 获得一致 3D 基形，并做 densification；第二阶段冻结高斯、优化 deformation MLP，以文本到视频扩散提供时序梯度，同时保留图像扩散梯度维持单帧质量。渲染视频的相机和时间被采样，潜空间的扩散分数再经可微 splatting 回传。论文另以高斯集合分布正则抑制退化的整体平移/缩放，并可用重叠区插值连续衔接多个短序列。

它与动态 3DGS 重建方法共享“高斯加变形”的渲染形式，但监督来源是生成模型，目标是符合文本先验而不是复现观测。由此可生成未观测动作，也同时受限于扩散先验、SDS/CSD 优化稳定性和短视频先验的时长。

% METHOD: 2312.15258
\paragraph{Human101（单目人体的快速 Gaussian avatar 优化）（2023.12）}
Human101~\cite{2312_15258} 针对单视角人体视频中隐式 avatar 训练和渲染缓慢的问题，采用显式 Gaussian rasterization 与人体先验结合的表示。输入视频先提供人体姿态/形状相关的参数化参考，高斯附着或在该参考的局部空间中优化，使每帧人体姿态改变时能把高斯一致地带到观察空间。

系统以图像重建监督优化高斯外观和局部几何，渲染时先根据当前姿态做人体驱动的变换，再执行深度排序的 alpha splatting。它的重点是缩短单人单目重建的优化路径，并让得到的显式高斯可直接实时渲染；单目遮挡、宽松衣物和未观测背面仍依赖人体先验及视频中的姿态覆盖，不能从单视角中被唯一恢复。

% METHOD: 2401.11535
\paragraph{EndoGS（表面对齐的可变形内窥组织 Gaussian）（2024.01）}
EndoGS~\cite{2401_11535} 针对单目内窥视频中组织非刚性形变、器械遮挡和普通体积 Gaussian 几何漂移。它以估计深度初始化表面对齐高斯，并结合已标注的器械 mask，把被器械遮住的像素从组织重建监督中剔除；组织随时间的形变由动态模型驱动，当前状态仍通过 Gaussian splatting 渲染。

表面对齐约束使高斯更贴近可见组织而非在深度方向形成体积云，图像与几何/形变一致性项共同优化高斯和动态参数。渲染时按当前时间变形高斯、投影并 alpha 合成，可输出实时新视图。它专为手术视频中的可见组织建模，工具遮挡后的区域在单目条件下没有充分三维证据；论文也明确指出单视角重建的病态性限制了直接用于需要完整解剖几何的下游任务。

% METHOD: 2401.15318
\paragraph{Gaussian Splashing（统一粒子上的物理运动与渲染）（2024.01）}
Gaussian Splashing~\cite{2401_15318} 将 3DGS 的椭球粒子同时作为可渲染基元和物理模拟离散元。它在粒子集合上注入流体/固体相互作用与约束，使同一组位置、形状和属性既可随物理状态更新，也可由 splatting 直接生成图像；目标是统一运动合成与视觉渲染，而非从视频反演一个 deformation field。

模拟步骤根据粒子材料和相互作用更新位置、速度及相关状态，渲染步骤将当前粒子状态投影并 alpha 合成。高斯的各向异性和显式邻域既提供局部几何，也让物理约束能够接在粒子层。该路线可把外力或交互写入场景演化，但物理真实性取决于材料模型、接触处理与初始属性，视觉拟合本身不能保证模拟参数正确。

% METHOD: 2402.03307
\paragraph{4D-Rotor Gaussian Splatting（以 4D rotor 参数化时空协方差）（2024.02）}
4D-Rotor Gaussian Splatting~\cite{2402_03307} 采用原生时空 Gaussian 表示动态场景，并重点改进四维旋转的参数化。4D 协方差需要同时描述空间轴、时间轴及其耦合；直接使用一般矩阵既冗余又难保持有效旋转。论文以 4D rotor 表示旋转，将时空均值、尺度与 rotor 组成可优化的 4D Gaussian。

给定时间，4D 分布被条件化/切片为对应的 3D Gaussian，再走标准投影和 alpha blending。训练从渲染误差端到端更新时空属性和密度控制，使运动由协方差的空间--时间相关性产生，而非 canonical MLP 输出。其优势是紧凑而连续的时空方向建模；但四维基元的参数与优化仍更复杂，极长序列或拓扑变化需要更多基元和稳定的时空采样。

% METHOD: 2403.02751
\paragraph{Splat-Nav（高斯地图上的安全规划与定位）（2024.03）}
Splat-Nav~\cite{2403_02751} 将已有 Gaussian splatting 地图转为机器人可用的几何接口，而不改变场景的渲染表示。Splat-Plan 利用每个高斯的椭球支撑进行碰撞检测，构造满足安全裕量的 corridor，并在目标变化或观测更新时在线重规划；这避免先把高斯场粗化为低分辨率 occupancy grid。

Splat-Loc 则以当前相机图像和由高斯地图渲染的预测视图对齐，优化/跟踪相机位姿，形成定位--规划闭环。因而该工作的方法单元是“椭球几何查询、走廊规划、渲染辅助位姿估计”，不是动态 4D 重建。安全性依赖地图已足够覆盖障碍物以及高斯椭球边界与保守裕量的设定；未建模的动态障碍或错误地图仍需要外部感知处理。

% METHOD: 2403.08321
\paragraph{ManiGaussian（面向多任务操作的时空 Gaussian 表示）（2024.03）}
ManiGaussian~\cite{2403_08321} 将 Gaussian splatting 用作语言条件机器人操作的时空感知接口。单帧或静态三维表征难描述抓取、推动等交互中物体状态的变化；该方法从观测序列构造动态 Gaussian 表示，提取跨视角、跨时间的三维特征，并将其与语言任务条件送入操作策略以预测动作。

高斯场负责以可微/显式方式保留工作台几何和外观，动态特征分支编码物体随操作的状态；策略训练将这些时空特征与示范动作的监督关联。它的技术贡献在于把动态渲染表示接到多任务 policy，而非仅提高 NVS 指标。能力受表示重建质量、训练任务的交互覆盖及视觉遮挡影响；若动态状态未被观测或高斯地图失配，动作条件也会随之退化。

% METHOD: 2403.11447
\paragraph{Motion-aware 3D Gaussian Splatting（运动感知约束的动态重建）（2024.03）}
Motion-aware 3D Gaussian Splatting~\cite{2403_11447} 指出仅逐帧最小化光度误差时，显式 Gaussian 的运动仍是欠约束的：相同图像可由错误的轨迹或几何解释。方法把视频中的运动线索接入动态高斯优化，使时间相邻帧的位移、几何与图像观测不只通过颜色间接耦合。

系统在规范高斯和时变运动的基础上，利用连续帧的 motion cue 约束高斯变换，并与渲染重建共同优化。当前时间的高斯仍走普通投影和 alpha 合成，因此运动约束作用于表示层而不是替换 rasterizer。该类监督可减少快速运动或稀疏观测中的轨迹歧义，但其可靠性受二维运动估计、遮挡处理和相机运动分离影响；错误 motion cue 也会向显式场传递偏差。

% METHOD: 2403.12365
\paragraph{GaussianFlow（由高斯运动溅射得到的光流监督）（2024.03）}
GaussianFlow~\cite{2403_12365} 解决的是“像素 optical flow 如何监督多高斯 alpha 合成后的 3D 动态”这一对应问题。单个像素通常同时由多个高斯贡献，像素流不能直接赋给其中一个粒子。论文将每个高斯的平移、旋转和尺度变化投影到图像，并以与颜色渲染相同的可见性/合成权重将这些二维运动聚合，定义为 Gaussian flow。

Gaussian flow 通过 tile-based splatting 计算，因而可微且开销接近常规渲染；把它与输入视频估计的 optical flow 匹配，梯度便可直接约束高斯运动。该损失可附加在动态新视角重建或扩散引导的 4D 内容生成上，和 RGB/SDS 监督互补。它不是新的 4D 表示，而是将现有时变高斯的动力学映射到可监督的像素量；光流在遮挡、无纹理或生成视频中不可靠时，监督边界也随之存在。

% METHOD: 2403.12722
\paragraph{HUGS（城市场景的背景--对象分解与语义高斯）（2024.03）}
HUGS~\cite{2403_12722} 面向驾驶城市场景的整体理解，而非只渲染 RGB。它将场景拆为静态背景和多个动态对象的 Gaussian 表示：背景在全局坐标中优化，对象在各自局部坐标中建模并由每帧对象位姿带到世界坐标。预测的 2D 语义、光流与 3D tracks 提供弱监督，连同图像重建共同恢复外观、几何、语义和运动。

渲染时将背景与经当前位姿变换的对象高斯一起投影、排序、合成；语义/深度等可由同一 splatting 管线输出。对象姿态和表示在训练中联合优化，避免要求人工精确 3D box。该分解使动态物体可跟踪并支持场景级语义查询，但依赖单目检测、跟踪和语义预测的质量，漏检或错误关联会污染对象--背景分配。

% METHOD: 2403.14939
\paragraph{STAG4D（时空锚定的生成式 4D Gaussian）（2024.03）}
STAG4D~\cite{2403_14939} 处理单图/单目视频生成 4D 内容时，未观测视角与相邻时间帧缺少一致监督的问题。其动态 Gaussian 表示由 4D 高斯和变形模块承载；生成训练不固定使用一个全局 densification 阈值，而在每次操作中选择累计梯度排名靠前的固定比例高斯，以适应不同生成对象的不确定性。

空间一致性来自多视角扩散先验：多视角模型的输出被用作图像到图像扩散 SDS 的条件。为让后续时间帧保留首帧的结构锚点，方法在无训练 temporal attention 中把后续帧 self-attention 的 key/value 与首帧 key/value 按权重混合，再以时序一致的多视角视频监督动态高斯。渲染仍是当前时间变形后高斯的标准 splatting。它属于生成式先验驱动，不等价于从真实视频直接确定几何；attention 混合与扩散先验可提升一致性，也会限制大幅姿态变化或引入先验偏差。

% METHOD: 2411.18613
\paragraph{CAT4D（动态场景与时间建模）（2024.11）}
CAT4D~\cite{2411_18613} 关注的问题是：建模动态场景、形变或 4D Gaussian 表示 CAT4D 是宽召回补全中发现的论文，主要根据标题、摘要和 arXiv 类别判定其与 3DGS 的关系。
相对原始 3DGS，CAT4D 的变化不是只替换损失函数，而是把表示、初始化、约束或网络预测路径重新组织到任务流程中：该工作被召回是因为论文文本包含 Gaussian/3D Gaussian 相关表述；当前分类先确定其任务域和与 3DGS 的相关性，后续证据流水线再补代码、企业和引用信息。
\textbf{实现细节：} 具体实现可以按输入、核心模块和渲染/优化路径来读。输入端首先取得任务所需的图像、相机、文本、时间序列、语义或几何线索，并把这些线索转换成可以约束高斯中心、尺度、旋转、不透明度、颜色或附加特征的条件。核心结构中，方法使用的主要机制是：从摘要层面看，CAT4D 围绕“建模动态场景、形变或 4D Gaussian 表示”展开。若其正文或代码后续可读，应进一步核对它具体改动的是表示、初始化、渲染、优化、损失、动态建模、压缩或应用接口中的哪一环。 训练或推理时，这些模块通常先给出初始化或中间特征，再通过可微 Gaussian rasterization、重建损失、感知/语义约束、几何正则或任务损失回传到高斯参数；如果方法包含前馈网络，则网络输出的是高斯属性、变形场、压缩码、语义特征或控制变量，而不是只输出最终 2D 图像。渲染路径仍以显式高斯投影、深度排序和 alpha 合成为基础，只是在采样、属性预测、局部编辑、动态变形、压缩编码或下游任务接口处加入额外约束。这样写入本章时，它与前后方法的关系是：相对原始 3DGS 和本章已有锚点方法，它可以理解为 3DGS extends\_or\_uses 宽召回分类阶段的关系判断，后续需要按论文方法部分复核。
\textbf{证据与边界：} 本轮召回分类尚未逐项核对实验表格；后续证据补全阶段需要读取论文源码/PDF或项目页确认数据集、指标和对比方法。 该条来自宽召回批次；源码下载受 arXiv e-print 网络限制，部分论文暂时只能依据标题、摘要和类别分类，需要在入可信书前做人工/LLM精读复核。 本书只采用本地 paper pack、数据库分类、源码/引用登记和论文摘要中可核验的信息；没有本地证据支撑的速度、排行榜或开源完整性不在正文中强化。\textbf{方法信息：} CAT4D project page lists Google DeepMind affiliation for authors 参与；引用信号：167 citations (8.6/month, domain-relative cutoff 5.3/month, semantic\_scholar)。
\textbf{设计动机与边界：} 相对原始 3DGS 以独立高斯和图像重建误差为主的设定，CAT4D 把问题落在以下可实现的改动上：从摘要层面看，CAT4D 围绕“建模动态场景、形变或 4D Gaussian 表示”展开。若其正文或代码后续可读，应进一步核对它具体改动的是表示、初始化、渲染、优化、损失、动态建模、压缩或应用接口中的哪一环。其设计目的不是单独增加一个模块，而是用上述表示、监督或渲染约束消除原流程在目标任务中的具体失配。该工作被召回是因为论文文本包含 Gaussian/3D Gaussian 相关表述；当前分类先确定其任务域和与 3DGS 的相关性，后续证据流水线再补代码、企业和引用信息。适用时还应保留论文给出的边界：该条来自宽召回批次；源码下载受 arXiv e-print 网络限制，部分论文暂时只能依据标题、摘要和类别分类，需要在入可信书前做人工/LLM精读复核。
% METHOD: 2412.09043
\paragraph{DrivingRecon（动态场景与时间建模）（2024.12）}
DrivingRecon~\cite{2412_09043} 关注的问题是：建模动态场景、形变或 4D Gaussian 表示 DrivingRecon 是宽召回补全中发现的论文，主要根据标题、摘要和 arXiv 类别判定其与 3DGS 的关系。
作为 supplementary 类条目，DrivingRecon 不应被写成单纯的新渲染算子。它更像是把 3DGS 放进 动态场景与4D 的协议、系统或应用边界里：该工作被召回是因为论文文本包含 Gaussian/3D Gaussian 相关表述；当前分类先确定其任务域和与 3DGS 的相关性，后续证据流水线再补代码、企业和引用信息。
\textbf{实现细节：} 具体实现可以按输入、核心模块和渲染/优化路径来读。输入端首先取得任务所需的图像、相机、文本、时间序列、语义或几何线索，并把这些线索转换成可以约束高斯中心、尺度、旋转、不透明度、颜色或附加特征的条件。核心结构中，方法使用的主要机制是：从摘要层面看，DrivingRecon 围绕“建模动态场景、形变或 4D Gaussian 表示”展开。若其正文或代码后续可读，应进一步核对它具体改动的是表示、初始化、渲染、优化、损失、动态建模、压缩或应用接口中的哪一环。 训练或推理时，这些模块通常先给出初始化或中间特征，再通过可微 Gaussian rasterization、重建损失、感知/语义约束、几何正则或任务损失回传到高斯参数；如果方法包含前馈网络，则网络输出的是高斯属性、变形场、压缩码、语义特征或控制变量，而不是只输出最终 2D 图像。渲染路径仍以显式高斯投影、深度排序和 alpha 合成为基础，只是在采样、属性预测、局部编辑、动态变形、压缩编码或下游任务接口处加入额外约束。这样写入本章时，它与前后方法的关系是：相对原始 3DGS 和本章已有锚点方法，它可以理解为 3DGS extends\_or\_uses 宽召回分类阶段的关系判断，后续需要按论文方法部分复核。
\textbf{证据与边界：} 本轮召回分类尚未逐项核对实验表格；后续证据补全阶段需要读取论文源码/PDF或项目页确认数据集、指标和对比方法。 该条来自宽召回批次；源码下载受 arXiv e-print 网络限制，部分论文暂时只能依据标题、摘要和类别分类，需要在入可信书前做人工/LLM精读复核。 本书只采用本地 paper pack、数据库分类、源码/引用登记和论文摘要中可核验的信息；没有本地证据支撑的速度、排行榜或开源完整性不在正文中强化。\textbf{方法信息：} 代码：\url{https://github.com/envision-research/driverecon}；代码关注：envision-research/driverecon has 191 stars (10.1/month, domain-relative cutoff 19.5/month)。
\textbf{设计动机与边界：} 相对原始 3DGS 以独立高斯和图像重建误差为主的设定，DrivingRecon 把问题落在以下可实现的改动上：从摘要层面看，DrivingRecon 围绕“建模动态场景、形变或 4D Gaussian 表示”展开。若其正文或代码后续可读，应进一步核对它具体改动的是表示、初始化、渲染、优化、损失、动态建模、压缩或应用接口中的哪一环。其设计目的不是单独增加一个模块，而是用上述表示、监督或渲染约束消除原流程在目标任务中的具体失配。该工作被召回是因为论文文本包含 Gaussian/3D Gaussian 相关表述；当前分类先确定其任务域和与 3DGS 的相关性，后续证据流水线再补代码、企业和引用信息。适用时还应保留论文给出的边界：该条来自宽召回批次；源码下载受 arXiv e-print 网络限制，部分论文暂时只能依据标题、摘要和类别分类，需要在入可信书前做人工/LLM精读复核。
% METHOD: 2412.09608
\paragraph{Representing Long Volumetric Video with Temporal Gaussian Hierarchy（动态场景与时间建模）（2024.12）}
Representing Long Volumetric Video with Temporal Gaussian Hierarchy~\cite{2412_09608} 关注的问题是：建模动态场景、形变或 4D Gaussian 表示 Representing Long Volumetric Video with Temporal Gaussian Hierarchy 是宽召回补全中发现的论文，主要根据标题、摘要和 arXiv 类别判定其与 3DGS 的关系。
作为 supplementary 类条目，Representing Long Volumetric Video with Temporal Gaussian Hierarchy 不应被写成单纯的新渲染算子。它更像是把 3DGS 放进 动态场景与4D 的协议、系统或应用边界里：该工作被召回是因为论文文本包含 Gaussian/3D Gaussian 相关表述；当前分类先确定其任务域和与 3DGS 的相关性，后续证据流水线再补代码、企业和引用信息。
\textbf{实现细节：} 具体实现可以按输入、核心模块和渲染/优化路径来读。输入端首先取得任务所需的图像、相机、文本、时间序列、语义或几何线索，并把这些线索转换成可以约束高斯中心、尺度、旋转、不透明度、颜色或附加特征的条件。核心结构中，方法使用的主要机制是：从摘要层面看，Representing Long Volumetric Video with Temporal Gaussian Hierarchy 围绕“建模动态场景、形变或 4D Gaussian 表示”展开。若其正文或代码后续可读，应进一步核对它具体改动的是表示、初始化、渲染、优化、损失、动态建模、压缩或应用接口中的哪一环。 训练或推理时，这些模块通常先给出初始化或中间特征，再通过可微 Gaussian rasterization、重建损失、感知/语义约束、几何正则或任务损失回传到高斯参数；如果方法包含前馈网络，则网络输出的是高斯属性、变形场、压缩码、语义特征或控制变量，而不是只输出最终 2D 图像。渲染路径仍以显式高斯投影、深度排序和 alpha 合成为基础，只是在采样、属性预测、局部编辑、动态变形、压缩编码或下游任务接口处加入额外约束。这样写入本章时，它与前后方法的关系是：相对原始 3DGS 和本章已有锚点方法，它可以理解为 3DGS extends\_or\_uses 宽召回分类阶段的关系判断，后续需要按论文方法部分复核。
\textbf{证据与边界：} 本轮召回分类尚未逐项核对实验表格；后续证据补全阶段需要读取论文源码/PDF或项目页确认数据集、指标和对比方法。 该条来自宽召回批次；源码下载受 arXiv e-print 网络限制，部分论文暂时只能依据标题、摘要和类别分类，需要在入可信书前做人工/LLM精读复核。 本书只采用本地 paper pack、数据库分类、源码/引用登记和论文摘要中可核验的信息；没有本地证据支撑的速度、排行榜或开源完整性不在正文中强化。\textbf{方法信息：} 引用信号：71 citations (3.8/month, domain-relative cutoff 5.3/month, semantic\_scholar)。
\textbf{设计动机与边界：} 相对原始 3DGS 以独立高斯和图像重建误差为主的设定，Representing Long Volumetric Video with Temporal Gaussian Hierarchy 把问题落在以下可实现的改动上：从摘要层面看，Representing Long Volumetric Video with Temporal Gaussian Hierarchy 围绕“建模动态场景、形变或 4D Gaussian 表示”展开。若其正文或代码后续可读，应进一步核对它具体改动的是表示、初始化、渲染、优化、损失、动态建模、压缩或应用接口中的哪一环。其设计目的不是单独增加一个模块，而是用上述表示、监督或渲染约束消除原流程在目标任务中的具体失配。该工作被召回是因为论文文本包含 Gaussian/3D Gaussian 相关表述；当前分类先确定其任务域和与 3DGS 的相关性，后续证据流水线再补代码、企业和引用信息。适用时还应保留论文给出的边界：该条来自宽召回批次；源码下载受 arXiv e-print 网络限制，部分论文暂时只能依据标题、摘要和类别分类，需要在入可信书前做人工/LLM精读复核。
% METHOD: 2501.04140
\paragraph{Spatiotemporal Gaussian Optimization for 4D Cone Beam CT Reconstruction from Sparse Projec（动态场景与时间建模）（2025.01）}
Spatiotemporal Gaussian Optimization for 4D Cone Beam CT Reconstruction from Sparse Projec~\cite{2501_04140} 关注的问题是：建模动态场景、形变或 4D Gaussian 表示 Spatiotemporal Gaussian Optimization for 4D Cone Beam CT Reconstruction from Spa 是宽召回补全中发现的论文，主要根据标题、摘要和 arXiv 类别判定其与 3DGS 的关系。
作为 supplementary 类条目，Spatiotemporal Gaussian Optimization for 4D Cone Beam CT Reconstruction from Sparse Projec 不应被写成单纯的新渲染算子。它更像是把 3DGS 放进 动态场景与4D 的协议、系统或应用边界里：该工作被召回是因为论文文本包含 Gaussian/3D Gaussian 相关表述；当前分类先确定其任务域和与 3DGS 的相关性，后续证据流水线再补代码、企业和引用信息。
\textbf{实现细节：} 具体实现可以按输入、核心模块和渲染/优化路径来读。输入端首先取得任务所需的图像、相机、文本、时间序列、语义或几何线索，并把这些线索转换成可以约束高斯中心、尺度、旋转、不透明度、颜色或附加特征的条件。核心结构中，方法使用的主要机制是：从摘要层面看，Spatiotemporal Gaussian Optimization for 4D Cone Beam CT Reconstruction from Spa 围绕“建模动态场景、形变或 4D Gaussian 表示”展开。若其正文或代码后续可读，应进一步核对它具体改动的是表示、初始化、渲染、优化、损失、动态建模、压缩或应用接口中的哪一环。 训练或推理时，这些模块通常先给出初始化或中间特征，再通过可微 Gaussian rasterization、重建损失、感知/语义约束、几何正则或任务损失回传到高斯参数；如果方法包含前馈网络，则网络输出的是高斯属性、变形场、压缩码、语义特征或控制变量，而不是只输出最终 2D 图像。渲染路径仍以显式高斯投影、深度排序和 alpha 合成为基础，只是在采样、属性预测、局部编辑、动态变形、压缩编码或下游任务接口处加入额外约束。这样写入本章时，它与前后方法的关系是：相对原始 3DGS 和本章已有锚点方法，它可以理解为 3DGS extends\_or\_uses 宽召回分类阶段的关系判断，后续需要按论文方法部分复核。
\textbf{证据与边界：} 本轮召回分类尚未逐项核对实验表格；后续证据补全阶段需要读取论文源码/PDF或项目页确认数据集、指标和对比方法。 该条来自宽召回批次；源码下载受 arXiv e-print 网络限制，部分论文暂时只能依据标题、摘要和类别分类，需要在入可信书前做人工/LLM精读复核。 本书只采用本地 paper pack、数据库分类、源码/引用登记和论文摘要中可核验的信息；没有本地证据支撑的速度、排行榜或开源完整性不在正文中强化。\textbf{方法信息：} 代码：\url{https://github.com/fuyabo/4dgs\_for\_4dcbct}；本条的源码/章节抽取不完整，正文只使用摘要、分类和已登记证据。
\textbf{设计动机与边界：} 相对原始 3DGS 以独立高斯和图像重建误差为主的设定，Spatiotemporal Gaussian Optimization for 4D Cone Beam CT Reconstruction from Sparse Projec 把问题落在以下可实现的改动上：从摘要层面看，Spatiotemporal Gaussian Optimization for 4D Cone Beam CT Reconstruction from Spa 围绕“建模动态场景、形变或 4D Gaussian 表示”展开。若其正文或代码后续可读，应进一步核对它具体改动的是表示、初始化、渲染、优化、损失、动态建模、压缩或应用接口中的哪一环。其设计目的不是单独增加一个模块，而是用上述表示、监督或渲染约束消除原流程在目标任务中的具体失配。该工作被召回是因为论文文本包含 Gaussian/3D Gaussian 相关表述；当前分类先确定其任务域和与 3DGS 的相关性，后续证据流水线再补代码、企业和引用信息。适用时还应保留论文给出的边界：该条来自宽召回批次；源码下载受 arXiv e-print 网络限制，部分论文暂时只能依据标题、摘要和类别分类，需要在入可信书前做人工/LLM精读复核。
% METHOD: 2502.03465
\paragraph{Seeing World Dynamics in a Nutshell（动态场景与时间建模）（2025.02）}
Seeing World Dynamics in a Nutshell~\cite{2502_03465} 关注的问题是：建模动态场景、形变或 4D Gaussian 表示 Seeing World Dynamics in a Nutshell 是宽召回补全中发现的论文，主要根据标题、摘要和 arXiv 类别判定其与 3DGS 的关系。
相对原始 3DGS，Seeing World Dynamics in a Nutshell 的变化不是只替换损失函数，而是把表示、初始化、约束或网络预测路径重新组织到任务流程中：该工作被召回是因为论文文本包含 Gaussian/3D Gaussian 相关表述；当前分类先确定其任务域和与 3DGS 的相关性，后续证据流水线再补代码、企业和引用信息。
\textbf{实现细节：} 具体实现可以按输入、核心模块和渲染/优化路径来读。输入端首先取得任务所需的图像、相机、文本、时间序列、语义或几何线索，并把这些线索转换成可以约束高斯中心、尺度、旋转、不透明度、颜色或附加特征的条件。核心结构中，方法使用的主要机制是：从摘要层面看，Seeing World Dynamics in a Nutshell 围绕“建模动态场景、形变或 4D Gaussian 表示”展开。若其正文或代码后续可读，应进一步核对它具体改动的是表示、初始化、渲染、优化、损失、动态建模、压缩或应用接口中的哪一环。 训练或推理时，这些模块通常先给出初始化或中间特征，再通过可微 Gaussian rasterization、重建损失、感知/语义约束、几何正则或任务损失回传到高斯参数；如果方法包含前馈网络，则网络输出的是高斯属性、变形场、压缩码、语义特征或控制变量，而不是只输出最终 2D 图像。渲染路径仍以显式高斯投影、深度排序和 alpha 合成为基础，只是在采样、属性预测、局部编辑、动态变形、压缩编码或下游任务接口处加入额外约束。这样写入本章时，它与前后方法的关系是：相对原始 3DGS 和本章已有锚点方法，它可以理解为 3DGS extends\_or\_uses 宽召回分类阶段的关系判断，后续需要按论文方法部分复核。
\textbf{证据与边界：} 本轮召回分类尚未逐项核对实验表格；后续证据补全阶段需要读取论文源码/PDF或项目页确认数据集、指标和对比方法。 该条来自宽召回批次；源码下载受 arXiv e-print 网络限制，部分论文暂时只能依据标题、摘要和类别分类，需要在入可信书前做人工/LLM精读复核。 本书只采用本地 paper pack、数据库分类、源码/引用登记和论文摘要中可核验的信息；没有本地证据支撑的速度、排行榜或开源完整性不在正文中强化。\textbf{方法信息：} 公开链接：\url{https://github.com/Nut-World/NutWorld}，完整性按本地证据保守处理；代码关注：nut-world/nutworld has 114 stars (6.6/month, domain-relative cutoff 19.5/month)。
\textbf{设计动机与边界：} 相对原始 3DGS 以独立高斯和图像重建误差为主的设定，Seeing World Dynamics in a Nutshell 把问题落在以下可实现的改动上：从摘要层面看，Seeing World Dynamics in a Nutshell 围绕“建模动态场景、形变或 4D Gaussian 表示”展开。若其正文或代码后续可读，应进一步核对它具体改动的是表示、初始化、渲染、优化、损失、动态建模、压缩或应用接口中的哪一环。其设计目的不是单独增加一个模块，而是用上述表示、监督或渲染约束消除原流程在目标任务中的具体失配。该工作被召回是因为论文文本包含 Gaussian/3D Gaussian 相关表述；当前分类先确定其任务域和与 3DGS 的相关性，后续证据流水线再补代码、企业和引用信息。适用时还应保留论文给出的边界：该条来自宽召回批次；源码下载受 arXiv e-print 网络限制，部分论文暂时只能依据标题、摘要和类别分类，需要在入可信书前做人工/LLM精读复核。
% METHOD: 2503.03115
\paragraph{NTR-Gaussian（2025.03）\cite{2503_03115}}
NTR-Gaussian 把 4DGS 从 RGB 动态外观扩展到夜间热红外温度场。普通 4DGS 学的是颜色随时间变化，而热红外图像测到的是表面辐射温度；如果只把温度当作颜色通道拟合，模型只能记住离散时间点，不能解释或预测温度随环境散热的连续变化。NTR-Gaussian 因此把高斯颜色替换为 infrared radiation temperature，并把对流换热、发射率和热容量等热力学参数接入时间预测。

方法分两阶段训练。第一阶段用 Temp-Net 学习离散时间和空间位置到高斯温度的映射，同时优化高斯几何和属性，使模型能在已有时间戳渲染热红外温度图。第二阶段把从初始时刻到目标时刻的时间区间划成若干段，Thermal-Net 接收当前温度、时间编码、位置编码和 Feature-GS 式语义特征，预测每段温度增量；这些增量通过热力学方程的积分策略累积到目标时刻。损失不再只是 RGB 重建，而是温度图的 $L_1$/SSIM 监督，并辅以平滑项，使同一位置的温度变化符合夜间散热的连续性。相对 ThermalGaussian 这类静态热重建，它的核心是从“某一时刻的温度贴图”变成“可在任意夜间时刻查询的 4D 温度场”。
\textbf{实现细节：} 热红外成像具有全天候工作能力，可在无光、雾霾等恶劣条件下非接触测量物体表面温度，广泛用于安防巡检、建筑监测和能源管理。然而，现有的 3D 热场景重建方法——无论是基于 NeRF 的 ThermoNeRF，还是基于 3DGS 的 Thermal Gaussian Splatting 和 Thermal3DGS——都只针对单一时刻做静态重建，无法建模或预测场景温度随时间的变化。NTR-Gaussian（2025.03） 首次将动态 4D 高斯表示引入热红外领域，目标是实现夜间户外场景的多视角、多时刻温度重建与新视角温度图合成。该方法的核心挑战在于：夜间没有太阳辐射输入，场景温度完全由散热过程决定，且温度随时间持续变化。在标准 4DGS 中，位置和协方差随时间变形，而颜色 $c$ 和不透明度 $\alpha$ 保持常数。NTR-Gaussian 将颜色属性替换为红外辐射温度 $T$ 后，温度也必须成为时变量——否则无法表达夜间持续降温的过程——这正是 NTR-Gaussian 需要建模的关键量。方法的输入为带时间戳的多视角热红外（TIR）图像及对应的合成可见光图像。方法分两阶段训练。第一阶段为粗粒度温度学习：Temp\_Net（MLP）建立离散时间到高斯温度的映射，输入为时间编码（对应热监督图像的采集时刻）和高斯位置编码（来自倾斜摄影模型的点云），输出每个高斯的辐射温度 $\mathrm{Temp}(t)$。第二阶段引入热力学先验，是整个方法的核心创新。夜间物体的温度变化遵循物理散热规律，可以用热力学微分方程约束高斯的温度预测。夜间场景的热量释放由两部分组成：对流换热 $Q_{\mathrm{con}} = C \cdot (T(t) - T_m)$ 和辐射散热 $Q_{\mathrm{rad}} = E \cdot \sigma \cdot (T(t))^{}4$，其中 $C$ 为对流传热系数（取决于物体材质和辐射温度）、$E$ 为发射率（取决于物体材质）、$\sigma$ 为 Stefan-Boltzmann 常数、$T_m$ 为环境温度。温度变化率满足微分方程 $T'(t) = (Q_{\mathrm{rad}} + Q_{\mathrm{con}}) / H$，其中 $H$ 为热容。
\textbf{方法信息：} 代码：\url{https://github.com/ykykykykyk/ntr-gaussian}。
\textbf{设计动机与边界：} 相对原始 3DGS 以独立高斯和图像重建误差为主的设定，NTR-Gaussian 把问题落在以下可实现的改动上：方法分为两阶段：(1) 粗优化阶段：优化Temp\_Net，基于离散时间编码和倾斜摄影模型的位置编码，获取高斯球的热温度和属性值。将3DGS中的颜色替换为红外辐射温度。(2) 精优化阶段：将时间区间划分为N段，利用热力学方程，通过Thermal\_Net预测每段温度变化。结合Feature\_GS的语义特征，联合优化网络以预测任意时刻的红外温度值。引入对流换热系数、发射率和热容量等热力学参数的神经网络预测。其设计目的不是单独增加一个模块，而是用上述表示、监督或渲染约束消除原流程在目标任务中的具体失配。首次将热力学原理（对流换热、辐射散热）融入4D Gaussian Splatting，提出NTR-Gaussian方法，能够从多视角热红外图像学习场景温度分布的4D时空表示，实现任意时刻、任意视角的温度图渲染，预测误差在1摄氏度以内。适用时还应保留论文给出的边界：需要时间戳标注的热红外图像；当前仅在夜间场景验证；热力学模型假设可能不适用于所有环境条件。
% METHOD: 2504.15122
\paragraph{MoBGS（2025.04）\cite{2504_15122}}
MoBGS 面向模糊单目动态视频，基线是 SplineGS 式静态/动态高斯表示。普通动态 NVS 默认输入帧是瞬时清晰观测；而运动模糊来自曝光区间内相机位姿和物体状态的连续积分。Deblur4DGS 已经把曝光模型接入 4DGS，MoBGS 的进一步改变是把“相机模糊程度”显式作为潜在相机轨迹估计的先验，并让同一个曝光区间同时约束全局相机运动和局部物体运动。

具体流程从一帧模糊图像开始。系统先用 DFT 中低频能量占比估计 blur score，再由 blur feature 和当前相机位姿进入 Blur-adaptive Latent Camera Estimation（BLCE）：一个 Neural ODE solver 在潜在空间生成 $N_l$ 个曝光期间的 latent camera poses，模糊越严重，轨迹可以覆盖更长、更复杂的路径。然后 Latent Camera-induced Exposure Estimation（LCEE）根据这些潜在相机姿态估计帧级曝光时间区间；在该区间内采样潜在时间戳，使用 SplineGS 的动态高斯预测物体状态，并分别渲染多个 latent sharp frames，最后平均成合成模糊帧与输入监督。训练使用 RGB 重建和深度项，LCEE 在前期几千次迭代后启用，代码默认每帧采样 $N_l=9$ 个潜在图像。相对把 2D 去模糊作为预处理的方法，MoBGS 把模糊形成过程直接写进 3DGS 渲染循环。
\textbf{实现细节：} MoBGS（2025.04） 从另一个角度切入运动模糊问题：不直接估计虚拟相机位姿序列，而是通过神经 ODE 建模连续的潜在相机轨迹。其核心模块是模糊自适应潜在相机估计（Blur-adaptive Latent Camera Estimation, BLCE），利用一个模糊自适应神经 ODE 求解器从模糊图像中推断曝光时间内的连续相机运动轨迹。此外，MoBGS 提出潜在相机引导曝光估计（Latent Camera-induced Exposure Estimation, LCEE）方法，确保全局相机运动和局部物体运动的去模糊一致性：利用估计的潜在相机轨迹推导曝光参数，将相机运动建模的精度传递到动态物体的去模糊流程中。DBMovi-GS：稀疏控制点的模糊感知动态重建。
\textbf{方法信息：} 代码：\url{https://github.com/kaist-viclab/mobgs}。



% METHOD: 2506.05348
\paragraph{FreeTimeGS（动态场景与时间建模）（2025.06）}
FreeTimeGS~\cite{2506_05348} 关注的问题是：建模动态场景、形变或 4D Gaussian 表示 FreeTimeGS 是宽召回补全中发现的论文，主要根据标题、摘要和 arXiv 类别判定其与 3DGS 的关系。
相对原始 3DGS，FreeTimeGS 的变化不是只替换损失函数，而是把表示、初始化、约束或网络预测路径重新组织到任务流程中：该工作被召回是因为论文文本包含 Gaussian/3D Gaussian 相关表述；当前分类先确定其任务域和与 3DGS 的相关性，后续证据流水线再补代码、企业和引用信息。
\textbf{实现细节：} 具体实现可以按输入、核心模块和渲染/优化路径来读。输入端首先取得任务所需的图像、相机、文本、时间序列、语义或几何线索，并把这些线索转换成可以约束高斯中心、尺度、旋转、不透明度、颜色或附加特征的条件。核心结构中，方法使用的主要机制是：从摘要层面看，FreeTimeGS 围绕“建模动态场景、形变或 4D Gaussian 表示”展开。若其正文或代码后续可读，应进一步核对它具体改动的是表示、初始化、渲染、优化、损失、动态建模、压缩或应用接口中的哪一环。 训练或推理时，这些模块通常先给出初始化或中间特征，再通过可微 Gaussian rasterization、重建损失、感知/语义约束、几何正则或任务损失回传到高斯参数；如果方法包含前馈网络，则网络输出的是高斯属性、变形场、压缩码、语义特征或控制变量，而不是只输出最终 2D 图像。渲染路径仍以显式高斯投影、深度排序和 alpha 合成为基础，只是在采样、属性预测、局部编辑、动态变形、压缩编码或下游任务接口处加入额外约束。这样写入本章时，它与前后方法的关系是：相对原始 3DGS 和本章已有锚点方法，它可以理解为 3DGS extends\_or\_uses 宽召回分类阶段的关系判断，后续需要按论文方法部分复核。
\textbf{证据与边界：} 本轮召回分类尚未逐项核对实验表格；后续证据补全阶段需要读取论文源码/PDF或项目页确认数据集、指标和对比方法。 该条来自宽召回批次；源码下载受 arXiv e-print 网络限制，部分论文暂时只能依据标题、摘要和类别分类，需要在入可信书前做人工/LLM精读复核。 本书只采用本地 paper pack、数据库分类、源码/引用登记和论文摘要中可核验的信息；没有本地证据支撑的速度、排行榜或开源完整性不在正文中强化。\textbf{方法信息：} FreeTimeGS project page lists Geely Automobile Research Institute affiliation for two authors 参与；引用信号：50 citations (3.8/month, domain-relative cutoff 5.3/month, semantic\_scholar)。
\textbf{设计动机与边界：} 相对原始 3DGS 以独立高斯和图像重建误差为主的设定，FreeTimeGS 把问题落在以下可实现的改动上：从摘要层面看，FreeTimeGS 围绕“建模动态场景、形变或 4D Gaussian 表示”展开。若其正文或代码后续可读，应进一步核对它具体改动的是表示、初始化、渲染、优化、损失、动态建模、压缩或应用接口中的哪一环。其设计目的不是单独增加一个模块，而是用上述表示、监督或渲染约束消除原流程在目标任务中的具体失配。该工作被召回是因为论文文本包含 Gaussian/3D Gaussian 相关表述；当前分类先确定其任务域和与 3DGS 的相关性，后续证据流水线再补代码、企业和引用信息。适用时还应保留论文给出的边界：该条来自宽召回批次；源码下载受 arXiv e-print 网络限制，部分论文暂时只能依据标题、摘要和类别分类，需要在入可信书前做人工/LLM精读复核。
% METHOD: 2507.16608
\paragraph{Dyna3DGR（动态场景与时间建模）（2025.07）}
Dyna3DGR~\cite{2507_16608} 关注的问题是：建模动态场景、形变或 4D Gaussian 表示 Dyna3DGR 是宽召回补全中发现的论文，主要根据标题、摘要和 arXiv 类别判定其与 3DGS 的关系。
相对原始 3DGS，Dyna3DGR 的变化不是只替换损失函数，而是把表示、初始化、约束或网络预测路径重新组织到任务流程中：该工作被召回是因为论文文本包含 Gaussian/3D Gaussian 相关表述；当前分类先确定其任务域和与 3DGS 的相关性，后续证据流水线再补代码、企业和引用信息。
\textbf{实现细节：} 具体实现可以按输入、核心模块和渲染/优化路径来读。输入端首先取得任务所需的图像、相机、文本、时间序列、语义或几何线索，并把这些线索转换成可以约束高斯中心、尺度、旋转、不透明度、颜色或附加特征的条件。核心结构中，方法使用的主要机制是：从摘要层面看，Dyna3DGR 围绕“建模动态场景、形变或 4D Gaussian 表示”展开。若其正文或代码后续可读，应进一步核对它具体改动的是表示、初始化、渲染、优化、损失、动态建模、压缩或应用接口中的哪一环。 训练或推理时，这些模块通常先给出初始化或中间特征，再通过可微 Gaussian rasterization、重建损失、感知/语义约束、几何正则或任务损失回传到高斯参数；如果方法包含前馈网络，则网络输出的是高斯属性、变形场、压缩码、语义特征或控制变量，而不是只输出最终 2D 图像。渲染路径仍以显式高斯投影、深度排序和 alpha 合成为基础，只是在采样、属性预测、局部编辑、动态变形、压缩编码或下游任务接口处加入额外约束。这样写入本章时，它与前后方法的关系是：相对原始 3DGS 和本章已有锚点方法，它可以理解为 3DGS extends\_or\_uses 宽召回分类阶段的关系判断，后续需要按论文方法部分复核。
\textbf{证据与边界：} 本轮召回分类尚未逐项核对实验表格；后续证据补全阶段需要读取论文源码/PDF或项目页确认数据集、指标和对比方法。 该条来自宽召回批次；源码下载受 arXiv e-print 网络限制，部分论文暂时只能依据标题、摘要和类别分类，需要在入可信书前做人工/LLM精读复核。 本书只采用本地 paper pack、数据库分类、源码/引用登记和论文摘要中可核验的信息；没有本地证据支撑的速度、排行榜或开源完整性不在正文中强化。\textbf{方法信息：} 代码：\url{https://github.com/windrise/dyna3dgr}。
\textbf{设计动机与边界：} 相对原始 3DGS 以独立高斯和图像重建误差为主的设定，Dyna3DGR 把问题落在以下可实现的改动上：从摘要层面看，Dyna3DGR 围绕“建模动态场景、形变或 4D Gaussian 表示”展开。若其正文或代码后续可读，应进一步核对它具体改动的是表示、初始化、渲染、优化、损失、动态建模、压缩或应用接口中的哪一环。其设计目的不是单独增加一个模块，而是用上述表示、监督或渲染约束消除原流程在目标任务中的具体失配。该工作被召回是因为论文文本包含 Gaussian/3D Gaussian 相关表述；当前分类先确定其任务域和与 3DGS 的相关性，后续证据流水线再补代码、企业和引用信息。适用时还应保留论文给出的边界：该条来自宽召回批次；源码下载受 arXiv e-print 网络限制，部分论文暂时只能依据标题、摘要和类别分类，需要在入可信书前做人工/LLM精读复核。
% METHOD: 2511.18794
\paragraph{ChronoGS（动态场景与时间建模）（2025.11）}
ChronoGS~\cite{2511_18794} 关注的问题是：建模动态场景、形变或 4D Gaussian 表示 ChronoGS 是宽召回补全中发现的论文，主要根据标题、摘要和 arXiv 类别判定其与 3DGS 的关系。
作为 supplementary 类条目，ChronoGS 不应被写成单纯的新渲染算子。它更像是把 3DGS 放进 动态场景与4D 的协议、系统或应用边界里：该工作被召回是因为论文文本包含 Gaussian/3D Gaussian 相关表述；当前分类先确定其任务域和与 3DGS 的相关性，后续证据流水线再补代码、企业和引用信息。
\textbf{实现细节：} 具体实现可以按输入、核心模块和渲染/优化路径来读。输入端首先取得任务所需的图像、相机、文本、时间序列、语义或几何线索，并把这些线索转换成可以约束高斯中心、尺度、旋转、不透明度、颜色或附加特征的条件。核心结构中，方法使用的主要机制是：从摘要层面看，ChronoGS 围绕“建模动态场景、形变或 4D Gaussian 表示”展开。若其正文或代码后续可读，应进一步核对它具体改动的是表示、初始化、渲染、优化、损失、动态建模、压缩或应用接口中的哪一环。 训练或推理时，这些模块通常先给出初始化或中间特征，再通过可微 Gaussian rasterization、重建损失、感知/语义约束、几何正则或任务损失回传到高斯参数；如果方法包含前馈网络，则网络输出的是高斯属性、变形场、压缩码、语义特征或控制变量，而不是只输出最终 2D 图像。渲染路径仍以显式高斯投影、深度排序和 alpha 合成为基础，只是在采样、属性预测、局部编辑、动态变形、压缩编码或下游任务接口处加入额外约束。这样写入本章时，它与前后方法的关系是：相对原始 3DGS 和本章已有锚点方法，它可以理解为 3DGS extends\_or\_uses 宽召回分类阶段的关系判断，后续需要按论文方法部分复核。
\textbf{证据与边界：} 本轮召回分类尚未逐项核对实验表格；后续证据补全阶段需要读取论文源码/PDF或项目页确认数据集、指标和对比方法。 该条来自宽召回批次；源码下载受 arXiv e-print 网络限制，部分论文暂时只能依据标题、摘要和类别分类，需要在入可信书前做人工/LLM精读复核。 本书只采用本地 paper pack、数据库分类、源码/引用登记和论文摘要中可核验的信息；没有本地证据支撑的速度、排行榜或开源完整性不在正文中强化。\textbf{方法信息：} 代码：\url{https://github.com/zhongtaowang/chronogs}。
\textbf{设计动机与边界：} 相对原始 3DGS 以独立高斯和图像重建误差为主的设定，ChronoGS 把问题落在以下可实现的改动上：从摘要层面看，ChronoGS 围绕“建模动态场景、形变或 4D Gaussian 表示”展开。若其正文或代码后续可读，应进一步核对它具体改动的是表示、初始化、渲染、优化、损失、动态建模、压缩或应用接口中的哪一环。其设计目的不是单独增加一个模块，而是用上述表示、监督或渲染约束消除原流程在目标任务中的具体失配。该工作被召回是因为论文文本包含 Gaussian/3D Gaussian 相关表述；当前分类先确定其任务域和与 3DGS 的相关性，后续证据流水线再补代码、企业和引用信息。适用时还应保留论文给出的边界：该条来自宽召回批次；源码下载受 arXiv e-print 网络限制，部分论文暂时只能依据标题、摘要和类别分类，需要在入可信书前做人工/LLM精读复核。
% METHOD: 2512.15719
\paragraph{Fast Volumetric Capture and Reconstruction Pipeline（2025.12）\cite{2512_15719}}
Fast Volumetric Capture and Reconstruction Pipeline 更像完整工程系统而非单一新算法。相对研究型 4DGS 或逐场景 Gaussian Splatting 重建，它的基线问题不是某个变形场是否更准，而是多 RGB-D/RGB 相机数据如何稳定进入体积捕捉流水线。它把动态高斯放到“采集、预处理、重建、导出、查看”的制作链路中：系统先统一执行人体/前景分割、深度到彩色图对齐、光流估计和联合双边滤波；这些中间结果不只服务一种输出，而是同时供点云重建和 Gaussian splat reconstruction 使用。

重建端有两条导出路径。点云路径从多相机深度和彩色融合出逐帧点云，可导出 PLY 或 MPEG V-PCC 相关格式；高斯路径则改造 GPS-Gaussian 式回归器，对旋转参数化、尺度预测和域内微调/checkpoint 做增强，输出可用于 Web、Unity、Unreal 或通用 viewer 的 PLY/SPLAT 资产。硬件上，论文讨论 6--12 台 Orbbec Femto Bolt 沿半圆弧布置的捕捉系统，因此它关心的不只是 PSNR，而是帧数据如何进入 GPU 预处理、如何快速形成可预览资产、如何进入现有引擎。当前官方仓库 \url{https://github.com/irc-hslu/capturestudio} 已公开 CaptureStudio 和 viewer 主流程，但 README 仍显示部分 reconstruction/export scripts 处于整理状态；因此这里把它视为公开工程流水线，而非已完全成熟的端到端研究代码。
\textbf{实现细节：} Fast Volumetric Capture（2025.12） 面向的核心问题是：如何用少量廉价相机在不受控环境中快速完成人体体积捕获。传统多视角系统依赖密集相机阵列和受控照明，部署成本高且灵活性差。该工作的出发点是将点云和高斯 splat 两条重建路径统一在同一预处理流程之下，从而在 6--12 台 Orbbec Femto Bolt RGB-D 相机组成的半圆弧形阵列下实现实时预览。在点云分支中，baseline 做法是将深度图直接反投影到世界坐标系。该系统沿用这一思路，但针对野外噪声增加了半径邻域有效性检验：对于每个世界空间点 $p$，在半径 $r = 20\,\mathrm{cm}$ 内搜索邻居，仅当邻居数 $|N(p)| \geq N_\theta = 30$ 时保留该点，孤立离群点被有效滤除。高斯 splat 分支是该系统的核心创新。传统 GPS-Gaussian 等前馈回归器将高斯旋转预测在相机坐标系下，却将其当作世界坐标系旋转直接使用。这一错误在圆形、等高布局的相机阵列中不易暴露（因为高度各向异性的"针状"splat 恰好与训练几何对齐），但在任意稀疏布局下会产生严重的视角依赖拉伸伪影。该工作对此进行了关键的坐标系修正：设网络预测的相机坐标系下高斯旋转为 ${_g^c R}$，相机到世界的标定旋转为 ${_c^w R}$，则正确的世界坐标系旋转应为 $${}_g^w R = {}_c^w R \cdot {}_g^c R,$$ 即相机到世界旋转必须左乘预测的相机坐标系旋转。相应地，世界坐标系下的协方差矩阵由 ${^w \Sigma} = {}_c^w R \cdot {}^c \Sigma \cdot ({}_c^w R)^{}\top$ 给出，其中 ${^c \Sigma} = {}_g^c R \cdot \mathrm{diag}(s^{}2) \cdot ({}_g^c R)^{}\top$ 是相机坐标系下的协方差。这一修正消除了视角依赖的拉伸伪影。在尺度参数化方面，该系统也做了改进。GPS-Gaussian 原本使用 softplus 激活，该系统改为先对预激活值施加双曲正切 $\tanh$ 有界化，再进行逐轴实例归一化并施加可学习的仿射参数： $$s\_{a,i} = s\_\text{max} \cdot \text{softplus}(\gamma\_a \hat{z}\_{a,i} + \beta\_a),$$ 其中 $\hat{z}_{a,i} = (z_{a,i} - \mu_a)/\sigma_a$ 是归一化后的预激活值，$\gamma_a, \beta_a$ 为可学习参数。这一改动使尺度参数能够跨不同相机距离和内参迁移。总损失为 $L = L_\mathrm{rec} + L_\mathrm{ent}$，其中熵正则化权重 $\lambda_\mathrm{ent} = 10^{-2}$。
\textbf{方法信息：} 代码：\url{https://github.com/irc-hslu/capturestudio}。
\textbf{设计动机与边界：} 相对原始 3DGS 以独立高斯和图像重建误差为主的设定，A Fast Volumetric Capture and Reconstruction Pipeline for Dynamic Point Clouds and Gaussia 把问题落在以下可实现的改动上：系统包含统一的预处理管线：执行分割、深度到色彩对齐、光流估计和联合双边深度滤波。从共享中间结果同时生成：(1)逐相机点云重建导出PLY和MPEG V-PCC压缩格式；(2)逐相机高斯溅射导出PLY和SPLAT格式。对GPS-Gaussian回归器进行多项改进，包括引入高斯旋转和缩放的增强处理。系统支持6-12台Orbbec Femto Bolt相机沿半圆弧排列，输出格式兼容现有播放器和游戏引擎。提供Web端查看器和Unity/Unreal插件。其设计目的不是单独增加一个模块，而是用上述表示、监督或渲染约束消除原流程在目标任务中的具体失配。提出快速高效的体积捕捉和重建系统，支持RGB-D和RGB输入生成点云和高斯溅射，改进GPS-Gaussian回归器实现高质量重建，支持野外操作、灵活相机配置，提供5-10FPS实时预览并开源完整框架。适用时还应保留论文给出的边界：实时预览帧率较低(5-10FPS)；高斯溅射质量受限于前馈回归器；复杂运动场景下可能出现时序不一致。
% METHOD: 2601.19843
\paragraph{Graphical X Splatting（2026.01）\cite{2601_19843}}
GraphiXS 把 4DGS 的确定性优化改写成概率图模型。标准 4DGS 往往默认相机同步、视角足够、帧没有缺失，优化目标就是找到一组能解释所有图像的高斯参数；但真实多相机系统会有稀疏视角、缺帧、相机异步、故障相机等问题。GraphiXS 的核心差异是把这些输入问题显式放进生成过程和潜在变量里，而不是把它们都当作普通重建误差。

它从混合模型视角理解 GS：一个场景是许多 component 的混合，渲染时给定相机和时间，先决定哪些 component 会被投到图像，再决定哪些 component 影响某个像素，最后由这些 component 合成颜色。GraphiXS 把 component 参数 $\theta$ 作为待推断潜变量，位置、速度、加速度等动态属性都可以放进其中；相机位姿和时间也可以视为随机变量，从而表达相机同步和采样不确定性。具体实例包括 GraphiGS 和 GraphiTS：前者使用高斯 component，后者使用 Student-t component，并通过自由度参数处理重尾噪声，使异常观测对模型的破坏更小。

优化上，GraphiXS 不只是做一次 MAP 点估计，而是引入 SGHMC 这类随机梯度哈密顿蒙特卡洛采样来近似后验推理；视角相关颜色仍可通过时间相关位置上的球谐函数表达。相对 4D-GS，它不是一个新的渲染器，而是一个能“升级”现有 4DGS 的不确定性框架：在标准、稀疏视角、稀疏帧、不同步相机和故障相机等设置下，模型能用先验和后验推理抵消数据缺口。
\textbf{实现细节：} 现有 4DGS 方法普遍假设数据质量充足——足够多的相机覆盖所有视角、良好的标定和同步。但在实际场景（如交通路口监控）中，这些假设常常不成立。GraphiXS（2026.01） 的核心贡献是将 4DGS 重新表述为一个概率图模型，把相机位姿 $\mathcal{C}$、帧时间 $\mathcal{T}$、视频帧 $\mathcal{I}$ 都视为随机变量，而非常量——直觉上，这意味着承认"我不确定相机在哪里、帧是何时拍的"，并将这种不确定性显式纳入优化过程。这使得框架能显式建模视角稀疏、帧缺失、相机异步等多种数据不确定性。GraphiXS 将 4DGS 的渲染流程重新表述为生成过程：从混合模型 $F = \sum_{i=1}^{M}\psi_i$（$M \to \infty$）出发，给定相机 $\mathcal{C}$ 和时间 $\mathcal{T}$，先采样投影到图像平面的组件子集 $\alpha_i$，再为每个像素的光线采样命中子集 $\beta_j$，最终生成颜色。原始 3DGS 可视为该框架的特例（无 $\mathcal{C}, \mathcal{T}$ 随机性，组件为高斯分布）。关键创新是所有可学习参数 $\Theta$ 被视为潜在变量，通过最大后验（MAP）推断： $$p(\Theta | \mathcal{C}, \mathcal{T}, \mathcal{I}) \propto p(\mathcal{I}|\Theta, \mathcal{C}, \mathcal{T}) \cdot p(\Theta) \cdot p(\mathcal{C}) \cdot p(\mathcal{T})$$ 其中 $p(\mathcal{C})$ 和 $p(\mathcal{T})$ 的引入使得相机位姿和时间戳的不确定性被纳入优化。该框架可实例化为不同变体：GraphiGS 使用高斯基元 $\phi_{\mathrm{GS}}(\mathbf{x}; \boldsymbol{\mu}, \boldsymbol{\Sigma}) = \exp(-\frac{1}{2}(\mathbf{x}-\boldsymbol{\mu})^T\boldsymbol{\Sigma}^{-1}(\mathbf{x}-\boldsymbol{\mu}))$；GraphiTS 使用 Student-t 分布，额外引入可学习的自由度参数 $\nu \in [1, +\infty)$ 控制尾部厚度 $\phi_{\mathrm{TS}}(\mathbf{x}; \boldsymbol{\mu}, \boldsymbol{\Sigma}, \nu) \propto (1+\frac{1}{\nu})^{-\frac{\nu+3}{2}}$，在 Cauchy 分布（重尾、高不确定性容忍）和高斯分布之间平滑插值。优化采用随机梯度哈密顿蒙特卡洛（SGHMC，一种结合随机梯度下降和哈密顿动力学的采样优化方法），4D 组件的添加/回收通过重要性得分 $s_i^t = \lambda_{sg} g_{sg,i}^t + \lambda_o o_i^t$ 调控。
\textbf{实现补充：} 实现上还要核对输入、输出、参数、特征、初始化、训练损失、正则、渲染或推理接口，以及优化中的剪枝、分裂、采样、编码或解码步骤；这些抓手决定该方法到底改了 3DGS 管线的哪一环。
\textbf{方法信息：} Baidu 参与了这篇论文；代码：\url{https://github.com/realcrane/graphical-x-splatting}；仓库目前约 3 stars，属于较新的公开实现。



% METHOD: 2602.22096
\paragraph{WeatherCity（动态场景与时间建模）（2026.02）}
WeatherCity~\cite{2602_22096} 关注的问题是：建模动态场景、形变或 4D Gaussian 表示 WeatherCity 是宽召回补全中发现的论文，主要根据标题、摘要和 arXiv 类别判定其与 3DGS 的关系。
作为 supplementary 类条目，WeatherCity 不应被写成单纯的新渲染算子。它更像是把 3DGS 放进 动态场景与4D 的协议、系统或应用边界里：该工作被召回是因为论文文本包含 Gaussian/3D Gaussian 相关表述；当前分类先确定其任务域和与 3DGS 的相关性，后续证据流水线再补代码、企业和引用信息。
\textbf{实现细节：} 具体实现可以按输入、核心模块和渲染/优化路径来读。输入端首先取得任务所需的图像、相机、文本、时间序列、语义或几何线索，并把这些线索转换成可以约束高斯中心、尺度、旋转、不透明度、颜色或附加特征的条件。核心结构中，方法使用的主要机制是：从摘要层面看，WeatherCity 围绕“建模动态场景、形变或 4D Gaussian 表示”展开。若其正文或代码后续可读，应进一步核对它具体改动的是表示、初始化、渲染、优化、损失、动态建模、压缩或应用接口中的哪一环。 训练或推理时，这些模块通常先给出初始化或中间特征，再通过可微 Gaussian rasterization、重建损失、感知/语义约束、几何正则或任务损失回传到高斯参数；如果方法包含前馈网络，则网络输出的是高斯属性、变形场、压缩码、语义特征或控制变量，而不是只输出最终 2D 图像。渲染路径仍以显式高斯投影、深度排序和 alpha 合成为基础，只是在采样、属性预测、局部编辑、动态变形、压缩编码或下游任务接口处加入额外约束。这样写入本章时，它与前后方法的关系是：相对原始 3DGS 和本章已有锚点方法，它可以理解为 3DGS extends\_or\_uses 宽召回分类阶段的关系判断，后续需要按论文方法部分复核。
\textbf{证据与边界：} 本轮召回分类尚未逐项核对实验表格；后续证据补全阶段需要读取论文源码/PDF或项目页确认数据集、指标和对比方法。 该条来自宽召回批次；源码下载受 arXiv e-print 网络限制，部分论文暂时只能依据标题、摘要和类别分类，需要在入可信书前做人工/LLM精读复核。 本书只采用本地 paper pack、数据库分类、源码/引用登记和论文摘要中可核验的信息；没有本地证据支撑的速度、排行榜或开源完整性不在正文中强化。\textbf{方法信息：} 代码：\url{https://github.com/irmvlab/weathercity}。
\textbf{设计动机与边界：} 相对原始 3DGS 以独立高斯和图像重建误差为主的设定，WeatherCity 把问题落在以下可实现的改动上：从摘要层面看，WeatherCity 围绕“建模动态场景、形变或 4D Gaussian 表示”展开。若其正文或代码后续可读，应进一步核对它具体改动的是表示、初始化、渲染、优化、损失、动态建模、压缩或应用接口中的哪一环。其设计目的不是单独增加一个模块，而是用上述表示、监督或渲染约束消除原流程在目标任务中的具体失配。该工作被召回是因为论文文本包含 Gaussian/3D Gaussian 相关表述；当前分类先确定其任务域和与 3DGS 的相关性，后续证据流水线再补代码、企业和引用信息。适用时还应保留论文给出的边界：该条来自宽召回批次；源码下载受 arXiv e-print 网络限制，部分论文暂时只能依据标题、摘要和类别分类，需要在入可信书前做人工/LLM精读复核。
% METHOD: 2605.23672
\paragraph{RiGS（2026.05）\cite{2605_23672}}
RiGS 把单目动态场景分成静态、刚体和瞬态三类高斯，解决单一变形场难以同时表达长期平滑运动和短期复杂扰动的问题。静态高斯表示固定背景；rigid Gaussians 表示长期、低频、近似刚性的运动，带有时间中心、生命周期、线速度和角速度等参数；transient Gaussians 表示短期、高频、出现又消失的复杂变化。这个拆分让方法不必用同一个 deformation MLP 同时解释背景、物体整体运动、遮挡和局部非刚体细节。

关键在于如何决定高斯属于哪一类。RiGS 聚合 SAM2、ViPE、MegaSAM 等 2D/3D 基础模型提供的远程光流、深度和跟踪点，形成对象级 dynamic mask 与 scene-flow supervision。静态区域保持固定；刚体高斯用 SE(3) 式低频运动解释长时间一致的物体运动；如果某个刚体基元生命周期短、运动不稳定或不能被低频模型解释，它可以过渡为 transient Gaussian。训练中先建立静态/动态分解，再用场景流和渲染损失约束刚体与瞬态分支，底层实现采用 2DGS/gsplat 相关组件以提高训练和渲染效率。相对 Deformable 3DGS，RiGS 的增量是把运动时间尺度显式写入表示：长期由刚体基元承担，短期由瞬态基元吸收。
\textbf{实现细节：} 现有单目 4DGS 方法要么逐帧优化（时间一致性差），要么采用刚体假设（无法建模快速非刚体运动）。RiGS（2026.05） 分析了基于刚体假设的高斯表示中动态高斯的时间持续时长分布，发现其呈现显著双峰模式：长时峰对应稳定低频运动，短时峰对应快速高频运动。这启发了将动态内容进一步分解为刚体和瞬态两种高斯基元。RiGS 引入三种高斯基元。静态高斯 $g^s = (\boldsymbol{\mu}^s, s^s, q^s, o^s, c^s)$ 表示不变背景。刚体高斯 $g^r = (\boldsymbol{\mu}^r, s^r, q^r, o^r, c^r, w, \delta^r, \tau^r)$ 通过 $K$ 个运动基 $\{T_{j,t} \in SE(3)\}$ 建模长期低频运动，查询时间 $t$ 的变换为 $\boldsymbol{\mu}^r(t) = (\sum_{j=1}^{K}w_j T_{j,t})\boldsymbol{\mu}^r$，其中 $\|w\|_2^2=1$ 确保归一化权重。瞬态高斯 $g^t = (\boldsymbol{\mu}^t, s^t, q^t, o^t, c^t, v, \delta^t, \tau^t)$ 通过可学习速度 $v$ 捕获短期高频运动：$\boldsymbol{\mu}^t(t) = \boldsymbol{\mu}^t + v(t - \tau^t)$。两者共享软门控时间衰减函数 $o(t) = o \cdot \sigma(\delta - |t - \tau^r|)$（$\delta$ 控制边界锐度，设为 3.0），使高斯仅在有限时间窗口内活跃。在动态掩码方面，RiGS 提出物体级掩码而非逐像素掩码：对 SAM2 分割的每个物体，计算其帧级运动分数 $s_{i,t}$（结合光流不确定性和 Sampson 误差），经时间滤波 $\tau^{\mathrm{temp}} = 10^{-4}$ 后聚合为整体运动分数 $s_i$，超过动态阈值 $\tau^{\mathrm{dyn}}$ 的整个物体被标记为动态。这种设计避免了同一物体被静态和动态高斯不一致地建模。训练中，刚体高斯可根据 $\delta^r$ 缩短自然退化为瞬态高斯，运动基数量 $K=10$。
\textbf{技术实现补充：} 现有单目 4DGS 方法要么逐帧优化（时间一致性差），要么采用刚体假设（无法建模快速非刚体运动）。RiGS（2026.05） 分析了基于刚体假设的高斯表示中动态高斯的时间持续时长分布，发现其呈现显著双峰模式：长时峰对应稳定低频运动，短时峰对应快速高频运动。这启发了将动态内容进一步分解为刚体和瞬态两种高斯基元。RiGS 引入三种高斯基元。静态高斯 $g^s = (\boldsymbol{\mu}^s, s^s, q^s, o^s, c^s)$ 表示不变背景。刚体高斯 $g^r = (\boldsymbol{\mu}^r, s^r, q^r, o^r, c^r, w, \delta^r, \tau^r)$ 通过 $K$ 个运动基 $\{T_{j,t} \in SE(3)\}$ 建模长期低频运动，查询时间 $t$ 的变换为 $\boldsymbol{\mu}^r(t) = (\sum_{j=1}^{K}w_j T_{j,t})\boldsymbol{\mu}^r$，其中 $\|w\|_2^2=1$ 确保归一化权重。瞬态高斯 $g^t = (\boldsymbol{\mu}^t, s^t, q^t, o^t, c^t, v, \delta^t, \tau^t)$ 通过可学习速度 $v$ 捕获短期高频运动：$\boldsymbol{\mu}^t(t) = \boldsymbol{\mu}^t + v(t - \tau^t)$。两者共享软门控时间衰减函数 $o(t) = o \cdot \sigma(\delta - |t - \tau^r|)$（$\delta$ 控制边界锐度，设为 3.0），使高斯仅在有限时间窗口内活跃。在动态掩码方面，RiGS 提出物体级掩码而非逐像素掩码：对 SAM2 分割的每个物体，计算其帧级运动分数 $s_{i,t}$（结合光流不确定性和 Sampson 误差），经时间滤波 $\tau^{\mathrm{temp}} = 10^{-4}$ 后聚合为整体运动分数 $s_i$，超过动态阈值 $\tau^{\mathrm{dyn}}$ 的整个物体被标记为动态。这种设计避免了同一物体被静态和动态高斯不一致地建模。训练中，刚体高斯可根据 $\delta^r$ 缩短自然退化为瞬态高斯，运动基数量 $K=10$。
\textbf{方法信息：} 代码：\url{https://github.com/ladvu/rigs}。

\begin{takeaway}
动态3DGS的发展可以理解为三次问题重排。第一，Dynamic 3D Gaussians把高斯变成可跟踪粒子，Deformable 3DGS和4D-GS把动态场景转成规范空间加时间变形，SpacetimeGS则把高斯本身提升为时空基元。第二，单目和野外场景暴露了欠约束问题，于是深度、光流、点跟踪、扩散模型、不确定性和静动态分解成为必要补充。第三，当动态高斯进入视频流、电影制作、自动驾驶、热红外或体积捕捉系统时，重点从单帧PSNR转向长序列一致性、可压缩、可编辑、可追踪和与传感器物理一致。
\end{takeaway}
\textbf{设计动机与边界：} 相对原始 3DGS 以独立高斯和图像重建误差为主的设定，RiGS 把问题落在以下可实现的改动上：RiGS引入三种高斯基元：(1)静态高斯：表示时间不变的背景，具有固定属性（位置、尺度、旋转、不透明度、颜色）；(2)刚体高斯：建模长期低频平滑运动，具有时序中心、前向/后向线性和角速度；(3)瞬态高斯：表示短期高频复杂运动。关键创新：(1)对象级动态掩码：聚合远程时空运动信息引导静态/动态区域分解；(2)刚体-瞬态过渡：刚体高斯根据时序寿命可过渡为瞬态高斯；(3)场景流引导：利用2D基础模型（VIPE, MegaSAM）获取光流、深度和2D跟踪点作为密集监督信号。训练使用多阶段策略。其设计目的不是单独增加一个模块，而是用上述表示、监督或渲染约束消除原流程在目标任务中的具体失配。首次提出多时间尺度的高斯表示（静态+刚体+瞬态），刚体到瞬态的过渡机制实现对长期低频和短期高频运动的统一建模，场景流引导提供密集3D运动监督。适用时还应保留论文给出的边界：依赖2D基础模型提供的监督质量；刚体-瞬态过渡阈值需要调参；长视频的计算开销较大
